\documentclass[11pt]{amsart}
\usepackage[margin=2.5cm]{geometry}
\usepackage[T1]{fontenc}
\usepackage{lmodern}
\usepackage{amsmath,amssymb,amsthm,mathtools}
\usepackage[shortlabels]{enumitem}
\usepackage{booktabs,array}
\usepackage[hidelinks]{hyperref}
\newcommand{\R}{\mathbb R}
\newcommand{\Z}{\mathbb Z}
\newcommand{\cG}{\mathcal G}
\newcommand{\cA}{\mathcal A}
\newcommand{\cT}{\mathcal T}
\newcommand{\cC}{\mathcal C}
\newcommand{\Reg}{\operatorname{Reg}}
\newcommand{\cs}{c^{*}}
\newcommand{\Wmin}{W_{\min}}
\newcommand{\ca}{\cos\alpha}
\newcommand{\sa}{\sin\alpha}
\newcommand{\ta}{\tan\alpha}
\newcommand{\se}{\sec\alpha}
\newcommand{\caz}{\cos\alpha_0}
\newcommand{\saz}{\sin\alpha_0}
\newcommand{\taz}{\tan\alpha_0}
\newcommand{\sez}{\sec\alpha_0}
\newcommand{\Kb}{\bar K}
\newtheorem{theorem}{Theorem}[section]
\newtheorem{proposition}[theorem]{Proposition}
\newtheorem{lemma}[theorem]{Lemma}
\newtheorem{corollary}[theorem]{Corollary}
\theoremstyle{definition}
\newtheorem{definition}[theorem]{Definition}
\theoremstyle{remark}
\newtheorem{remark}[theorem]{Remark}
\numberwithin{equation}{section}

\title[An upper bound of order $k^{3/8}$ for the deficiency]{Packing $k^2-c$ unit squares:\\ an upper bound of order $k^{3/8}$ for the deficiency}
\author{Sungjoon Ryu}
\date{Preprint, version 1.1, 10 October 2026}

\raggedbottom
\begin{document}
\begin{abstract}
For $n\ge1$ let $s(n)$ be the side of the smallest square that contains $n$ unit squares with pairwise disjoint
interiors, and for an integer $k\ge2$ let $\cs(k)$ be the largest integer $c$ with $s(k^2-c)=k$. We prove a tiered
bound $\cs(k)<C_ik^{3/8}+a_i$ for every integer $k\ge k_i$, for example
\[
\cs(k)<48.38\,k^{3/8}+0.017\quad(k\ge6.1\cdot10^{14})\qquad\text{and}\qquad
\cs(k)<40.62\,k^{3/8}+1.4\cdot10^{-5}\quad(k\ge3.4\cdot10^{26}),
\]
and, with two further lemmas on the wall filler, $\cs(k)<45.08\,k^{3/8}+0.026$ for every integer
$k\ge6.4\cdot10^{13}$;
with much better constants for moderate $k$ we prove $\cs(k)<20.668\,k^{2/5}+0.623$ for every integer
$k\ge1.6\cdot10^7$.
The tiered bound is smaller than $20.668\,k^{2/5}+0.623$ for every $k\ge6.1\cdot10^{14}$, with two further lemmas for every
$k\ge5.07\cdot10^{13}$ (tier B1s), and with a sawtooth accounting of the wall rows and a choice of the lifts (tiers
$\ddagger$, for example $\cs(k)<41.19\,k^{3/8}+0.04$ for $k\ge5.22\cdot10^{12}$) for every $k\ge2.48\cdot10^{12}$
(tier C3).
The frame is an L-shaped arrangement in the spirit of Erd\H{o}s and Graham and of
Chung and Graham, in which two bands of width slightly less than an integer $b$ are filled with rows of $b$ unit squares
tilted by slightly more than $2\arctan(1/b)$; this frame alone gives $\cs(k)\le8\lceil\sqrt{k-4}\,\rceil-1$ for $k\ge6$,
which we also prove. The rows are lifted, and the four ends of the bands are filled by a ``staircase'': columns of
squares with one common small tilt, whose lowest squares are replaced by horizontal layers so that the columns follow
the rising end of the band. This gives an uncovered area $O(b^{2/3})$ per end and the exponent $2/5$. Filling the two
walls of each end by chains of staircases whose tilt is lowered from block to block gives $O(b^{3/5})$ per end and the
exponent $3/8$. These ideas are due to Bui; our proofs are self-contained. The constants of the $3/8$ bounds are
certified by interval arithmetic; their range of validity cannot be tested by explicit computation, and the explicit
checks test the construction at smaller sizes. Version~1.1 replaces the single bound $43.06\,k^{3/8}+2\cdot10^{-5}$
($k\ge4\cdot10^{26}$) of version~1.0 by the tiers. Explicit packings, checked in exact rational arithmetic, also by
checkers written independently of the generating program, give for example $\cs(10^5)\le1583$ and
$\cs(10^8)\le24790$. Together with the lower bound of order $k^{1/3}$ of [T] (a preprint, not refereed),
$k^{1/3}\ll\cs(k)\ll k^{3/8}$; the true exponent is open. The paper has not been refereed.
\end{abstract}
\maketitle

\section{Introduction and results}\label{sec:intro}

A \emph{packing} of $n$ unit squares in a closed square $K$ is a family of $n$ closed unit squares, in any orientation,
contained in $K$ and with pairwise disjoint interiors. For $n\ge1$ let $s(n)$ be the infimum of the sides of the squares
in which $n$ unit squares can be packed. Removing squares from a packing leaves a packing, so $s$ is non-decreasing.
By area $s(k^2)\ge k$, and the grid gives $s(k^2)\le k$; hence $s(k^2-c)\le k$ for $0\le c<k^2$. For an integer $k\ge2$
let
\[
\cs(k):=\max\{c\in\Z:\ 0\le c\le k^2-1,\ s(k^2-c)=k\},
\]
the \emph{deficiency}: removing up to $\cs(k)$ squares from the $k\times k$ grid does not allow a smaller container,
while removing $\cs(k)+1$ squares does. The wasted area of Erd\H{o}s and Graham [EG] is $W(x):=x^2-M(x)$, where $M(x)$ is
the largest number of unit squares that can be packed in a square of side $x$. In Section~\ref{sec:prelim} we record
the elementary relation
\begin{equation}\label{eq:translation}
\cs(k)=\lim_{x\to k^-}W(x)-1 ,
\end{equation}
so $\cs(k)$ measures the wasted area just below an integer.

\emph{Lower bounds.} The author showed $\cs(k)\to\infty$ in [R], with a bound of order $\log k$; lower bounds of
order $k^{1/4}$ and $k^{1/3}$ followed in [Q] and [T]. These are preprints that have not been refereed; see
Section~\ref{sec:lower}.

\emph{Upper bounds.} From the literature on $W(x)$ one gets $\cs(k)=O(k^{3/5})$ (preprints [Bui], [McC]); the refereed
literature gives $O(k^{7/11})$ [EG] and $O(k^{(3+\sqrt2)/7}\log k)$ [CG], and Arslanov, Mustafin and Shangitbayev
[AMS] proved $\cs(k)\le k-1$ for $k>12$. See Section~\ref{sec:lit}.

\subsection{Results}
The main result is the following.

\begin{theorem}[computer-assisted constants]\label{thm:main38}
For each tier $i$ of Table~\ref{tab:tiers} and every integer $k\ge k_i$,
\[
\cs(k)\ <\ C_i\,k^{3/8}+a_i .
\]
Here $C_i\ge\frac{32}{5}\bigl(\frac53\bigr)^{3/8}C_{E,i}^{5/8}$, $k_i\ge\frac35C_{E,i}b_{0,i}^{8/5}$ and
$a_i\ge\frac{12}5C_{E,i}b_{0,i}^{-2/5}$, where $b_{0,i}$ and $C_{E,i}$ are the threshold and the end constant of
Corollary~\ref{cor:b35} for the parameters of tier $i$ (Table~\ref{tab:certtiers}).
The tiers marked $\dagger$ use, in addition to the lemmas used by the other tiers, the construction ZC$'$ and
Lemmas~\ref{lem:P3a}--\ref{lem:Wpp} of Section~\ref{ssec:refine}.
The tiers marked $\ddagger$ use, in addition to these, Lemmas~\ref{lem:SL}, \ref{lem:T1}, \ref{lem:W4} and~\ref{lem:T2}
and Corollary~\ref{cor:b35c} (Sections~\ref{ssec:sawtooth} and~\ref{ssec:liftT2}); for them \(b_{0,i}\) and \(C_{E,i}\) are those of Corollary~\ref{cor:b35c} (Table~\ref{tab:certC}).
\end{theorem}

\begin{table}[ht]
\centering
\small
\begin{tabular}{lrrrrrr}
\toprule
tier $i$ & $b_{0,i}$ & $C_{E,i}$ & $C_i$ & $k_i$ & $a_i$ & better than Thm.~\ref{thm:main} for $k\ge$\\
\midrule
B1$^\dagger$ & $10^{7.7}$ & $17.461$ & $46.31$ & $2.2\cdot10^{13}$ & $0.035$ & $1.04\cdot10^{14}$\\
B1s$^\dagger$ & $10^{7.9}$ & $16.968$ & $45.49$ & $4.5\cdot10^{13}$ & $0.029$ & $5.07\cdot10^{13}$\\
B2$^\dagger$ & $10^{8}$ & $16.721$ & $45.08$ & $6.4\cdot10^{13}$ & $0.026$ & $6.4\cdot10^{13}$\\
A1 & $10^{8.5}$ & $18.961$ & $48.76$ & $4.6\cdot10^{14}$ & $0.019$ & $8.14\cdot10^{14}$\\
A1s & $10^{8.58}$ & $18.722$ & $48.38$ & $6.1\cdot10^{14}$ & $0.017$ & $6.1\cdot10^{14}$\\
A2 & $10^{9}$ & $17.523$ & $46.42$ & $2.7\cdot10^{15}$ & $0.011$ & $2.7\cdot10^{15}$\\
A3 & $10^{10}$ & $15.859$ & $43.61$ & $9.6\cdot10^{16}$ & $0.0039$ & $9.6\cdot10^{16}$\\
A4 & $10^{12}$ & $14.778$ & $41.73$ & $1.5\cdot10^{20}$ & $5.7\cdot10^{-4}$ & $1.5\cdot10^{20}$\\
A5 & $10^{16}$ & $14.153$ & $40.62$ & $3.4\cdot10^{26}$ & $1.4\cdot10^{-5}$ & $3.4\cdot10^{26}$\\
C1$^\ddagger$ & $10^{7.3617}$ & $14.473$ & $41.19$ & $5.22\cdot10^{12}$ & $0.04$ & $5.22\cdot10^{12}$\\
C2$^\ddagger$ & $10^{7.2}$ & $14.805$ & $41.78$ & $2.95\cdot10^{12}$ & $0.047$ & $2.95\cdot10^{12}$\\
C3$^\ddagger$ & $10^{7.15}$ & $14.983$ & $42.09$ & $2.48\cdot10^{12}$ & $0.05$ & $2.48\cdot10^{12}$\\
\bottomrule
\end{tabular}
\medskip
\caption{The tiers of Theorem~\ref{thm:main38}. The constants are rounded up. The last column is
$\max(k_i,k_{x,i})$, where $k_{x,i}:=(C_i/20.668)^{40}$ (rounded up to three digits) is the point from which
$C_ik^{3/8}<20.668\,k^{2/5}$; as $a_i<0.623$, from there on the tier gives a smaller bound than
Theorem~\ref{thm:main}.
Tiers marked $\dagger$ use ZC$'$ and Lemmas~\ref{lem:P3a}--\ref{lem:Wpp} (Section~\ref{ssec:refine}).
Tiers marked $\ddagger$ use, in addition, Lemmas~\ref{lem:SL}--\ref{lem:W4} and~\ref{lem:T2} (Sections~\ref{ssec:sawtooth}
and~\ref{ssec:liftT2}); their parameters are in Table~\ref{tab:certC}.
}\label{tab:tiers}
\end{table}

The proof is in Section~\ref{sec:38}. Its constants are certified by interval arithmetic (Corollaries~\ref{cor:b35} and~\ref{cor:b35c}).
The range $k\ge k_i$, where one end of a band has more than $10^{7}$ pieces, has not been tested by explicit
computation, except for $15$ end packings of the tiers $\ddagger$ near $b_0$ (``Verification of version~1.1''); the
explicit checks of Remark~\ref{rem:38exact} test the validity of the construction and of its lemmas at smaller sizes.

\emph{Version~1.0.} Version~1.0 of this paper [U] stated $\cs(k)<43.06\,k^{3/8}+2\cdot10^{-5}$ for every integer
$k\ge4\cdot10^{26}$ (with $C_E=15.54$ and $b_0=10^{16}$). That statement remains proved (it is the case
$(b_0,c_y,c_D,c_R,a_w,\varpi)=(10^{16},1/12,4,4,3/2,0.65)$ of Section~\ref{sec:38}, see
Section~\ref{ssec:proofthm38}), and it is superseded by the tier A5: $\cs(k)<40.62\,k^{3/8}+1.4\cdot10^{-5}$ for
$k\ge3.4\cdot10^{26}$.

For moderate $k$ the staircase alone gives a smaller bound with much better constants:

\begin{theorem}\label{thm:main}
For every integer $k\ge k_0:=1.6\cdot10^7$,
\[
\cs(k)\ <\ \frac{20}{3}\Bigl(\frac32\Bigr)^{2/5}5.03^{3/5}\,k^{2/5}+\frac83\cdot5.03\cdot10^{-4/3}
\ <\ 20.668\,k^{2/5}+0.623 .
\]
\end{theorem}

The proof (Section~\ref{sec:stair}) gives $\cs(k)\le(18.89+o(1))\,k^{2/5}$ as $k\to\infty$. At $k=k_0$ the bound is
$15738.5$; the bound of Theorem~\ref{thm:sqrt} below is $31999$ there.

\emph{Which bound to use.} Of the two bounds, only Theorem~\ref{thm:main} is proved for $1.6\cdot10^7\le k<2.48\cdot10^{12}$
($k<2.2\cdot10^{13}$ without Sections~\ref{ssec:sawtooth} and~\ref{ssec:liftT2}, and $k<4.6\cdot10^{14}$ without
Section~\ref{ssec:refine} either).
From there on both theorems apply, and for each $k$ one may take the smallest of the bounds that apply.
With the tiers that use only Lemmas~\ref{lem:W} and~\ref{lem:E38} (no $\dagger$), Theorem~\ref{thm:main38} gives
a smaller bound than $20.668\,k^{2/5}+0.623$ (Theorem~\ref{thm:main}) for every $k\ge6.1\cdot10^{14}$ (tier A1s). With
the tiers marked $\dagger$
as well, it does so for every $k\ge5.07\cdot10^{13}$ (tier B1s).
With the tiers marked $\ddagger$ as well, it does so for every $k\ge2.48\cdot10^{12}$ (tier C3).
Below that point Theorem~\ref{thm:main} is the better bound (or the only one), and it also comes with explicit packings
for moderate $k$ (Theorem~\ref{thm:cert} below); Theorem~\ref{thm:main38} gives the exponent $3/8$. Both bounds are
far above the values of explicit packings, and neither is claimed to be close to $\cs(k)$.

The frame of the proof is an L-shaped packing. For an integer $b\ge2$ put $D:=b^2+1$,
\[
\gamma_0:=\frac{b^2-1}{D},\qquad h_0:=\frac{3b^2-1}{D},\qquad
R(L):=\max\bigl(0,\ \lceil (L-h_0)\gamma_0\rceil\bigr)\quad(L\in\Z),
\]
and for integers $2\le b\le k$
\[
N(k,b):=(k-b)^2+b\,\bigl(R(k)+R(k-b)\bigr).
\]

\begin{theorem}\label{thm:A}
Let $k\ge2$ and $2\le b\le k$ be integers with $(k,b)\ne(2,2)$. Then $N(k,b)$ unit squares can be packed in a square of
side less than $k$. Consequently, for all $2\le b\le k$,
\[
\cs(k)\ \le\ k^2-N(k,b)-1\ <\ 6b+\frac{4k}{b}-3 .
\]
\end{theorem}

\begin{theorem}\label{thm:sqrt}
For every integer $k\ge6$,
\[
\cs(k)\ \le\ 8\bigl\lceil\sqrt{k-4}\,\bigr\rceil-1\ <\ 8\sqrt{k-4}+7 .
\]
\end{theorem}

The bound of Theorem~\ref{thm:sqrt} is the best that the packings of Theorem~\ref{thm:A} give, up to an additive
constant:

\begin{proposition}\label{prop:D}
For all integers $2\le b\le k$ one has $k^2-N(k,b)>8\sqrt k-14$.
\end{proposition}

The packings of Theorem~\ref{thm:A} lose area at the four ends of the two tilted bands, about $b$ at each end
(Section~\ref{sec:count}). In Section~\ref{sec:ends} we lift the rows of each band and show that the deficiency is
exactly the body cost of the bands, about $4k/b$, plus the uncovered areas of four end regions
(Proposition~\ref{prop:waste}). Section~\ref{sec:stair} fills an end region by a staircase of columns with uncovered
area $O(b^{2/3})$; with $b\asymp k^{3/5}$ this gives Theorem~\ref{thm:main}. Section~\ref{sec:38} replaces the two
walls of the staircase by chains of staircases with slowly decreasing tilt; the uncovered area becomes $O(b^{3/5})$,
and with $b\asymp k^{5/8}$ this gives Theorem~\ref{thm:main38}. With tuned parameters the staircase construction gives
explicit packings for moderate $k$:

\begin{theorem}[computer-assisted]\label{thm:cert}
For each row of Table~\ref{tab:cert}, $k^2-c$ unit squares can be packed in a square of side less than $k$. Hence
$\cs(k)\le c-1$, as listed.
\end{theorem}

\begin{table}[ht]
\centering
\begin{tabular}{rrrrrrr}
\toprule
$k$ & $b$ & $y_0$ & $c=k^2-N$ & bound & Thm.~\ref{thm:main} & $8\lceil\sqrt{k-4}\rceil-1$\\
\midrule
$10^5$ & $623$ & $257/4$ & $1584$ & $\cs(k)\le1583$ & $2067$ & $2535$\\
$10^6$ & $2481$ & $653/4$ & $3973$ & $\cs(k)\le3972$ & $5192$ & $7999$\\
$10^7$ & $9875$ & $1663/4$ & $10040$ & $\cs(k)\le10039$ & $13041$ & $25303$\\
$3.825\cdot10^7$ & $22086$ & $710$ & $16989$ & $\cs(k)\le16988$ & $22302$ & $49479$\\
$10^8$ & $39314$ & $1041$ & $24791$ & $\cs(k)\le24790$ & $32757$ & $79999$\\
\bottomrule
\end{tabular}
\medskip
\caption{Certified packings (Theorem~\ref{thm:cert}). The band width is $b-\varepsilon$ with
$\tan(\theta/2)=b/(b^2-1)$, and $y_0$ is the lift of the first row of each band. The column ``Thm.~\ref{thm:main}''
gives the integer part of $20.668\,k^{2/5}+0.623$ (Theorem~\ref{thm:main} itself applies only for $k\ge1.6\cdot10^7$),
and the last column is the bound of Theorem~\ref{thm:sqrt}. These staircase certificates are the best values we have
at these $k$. The construction of Section~\ref{sec:38} gives larger values here (for example $\cs(10^8)\le44299$,
Remark~\ref{rem:38exact}); it wins only for very large $k$.}\label{tab:cert}
\end{table}

For comparison, the explicit bound of Wang, Dong and Li [WDL] would give about $8.1\cdot10^4$ at $k=10^5$
(Section~\ref{sec:lit}).

To the best of our knowledge, no upper bound of order $o(k^{3/5})$ for $\cs(k)$, or for $W(x)$ just below an integer,
has been stated in the sources we searched (listed in Section~\ref{sec:lit}). The ingredients of
Theorem~\ref{thm:A} are standard, and the staircase and the chained wall stairs are due to Bui
(Section~\ref{ssec:credit}); what seems not to have been written down is their use just below an integer, with bands of
width $b-\varepsilon$.

\emph{The order of $\cs(k)$.} Theorem~\ref{thm:main38} (any one tier) and the lower bound of [T] give
\[
k^{1/3}\ \ll\ \cs(k)\ \ll\ k^{3/8}.
\]
By a theorem of Roth and Vaughan, a bound $W(x)\le Cx^{\alpha}$ valid for all large real $x$ must have
$\alpha\ge1/2$, so by \eqref{eq:translation} such bounds cannot give $\cs(k)=o(\sqrt k)$
(Section~\ref{ssec:structural}); Theorems~\ref{thm:main38} and~\ref{thm:main} use that the band width is just below an
integer. These are the bounds proved here and in [T]; the true exponent of $\cs(k)$ is open: it lies between $1/3$ and
$3/8$ (Section~\ref{sec:open}).

\subsection{Credits}\label{ssec:credit}
The skeleton of the packing of Theorem~\ref{thm:A} (an axis-parallel square, and two bands along two of its sides
filled with tilted blocks that touch both sides of the band) is the construction of Erd\H{o}s and Graham [EG,
Figs.~1--2]; the same L-shaped skeleton is the Type~1 packing of Chung and Graham [CG]. The block (a $1\times b$
rectangle of $b$ unit squares tilted by about $2\arctan(1/b)$ so that it fits in a band of width slightly less than $b$)
is used by Arslanov, Mustafin and Shangitbayev [AMS] for small $b$ and in a lemma for even $b$; they do not chain the
blocks into long bands. The staircase of Section~\ref{sec:stair}, in which tilted columns of one fixed tilt keep only
their upper squares and their lower squares are replaced by horizontal axis-parallel layers, so that one tilt can
follow a rising ceiling, is the ``primitive'' construction of H.~D.~Bui [Bui, Section~3]. The wall filler of
Section~\ref{sec:38}, in which a thin trapezoid is filled by a chain of such staircases whose tilt is lowered from block
to block so that the layers close in phase, follows the ideas of [Bui, Sections~3--4] (arXiv:2508.04603v2). Our proofs
of the end bounds are self-contained and do not use any argument of [Bui]; the write-up and the constants are ours.

\subsection{Status}\label{ssec:status}
Theorems~\ref{thm:A} and~\ref{thm:sqrt} and Proposition~\ref{prop:D} are proved in full in
Sections~\ref{sec:prelim}--\ref{sec:count}, Theorem~\ref{thm:main} in Sections~\ref{sec:ends}
and~\ref{sec:stair}, and Theorem~\ref{thm:main38} in Sections~\ref{sec:ends} and~\ref{sec:38}.
Theorem~\ref{thm:main} relies on one evaluation in interval arithmetic,
$\Phi(10^{-4/3})\le5.0163$ (Corollary~\ref{cor:b23}), done by the author's program and by three independent
computations; the other numerical constants were certified in the same way. The constants of
Theorem~\ref{thm:main38} (Corollaries~\ref{cor:b35} and~\ref{cor:b35c}) are certified, tier by tier, by interval arithmetic on covers of the
whole range $b\ge b_{0,i}$, by the author's covering and by three coverings written independently from the text: two
by verification agents (\texttt{mpmath.iv} and \texttt{python-flint} arb) and one by a reviewer (\texttt{mpmath.iv}), and all three implement the same normalised inequalities (the lists of Appendices~\ref{app:enc} and~\ref{app:encC}), so they check the arithmetic and the coverage but are only a weak check of the transcription of these inequalities from the lemmas; the point checks of the un-normalised formulas (at \(90\), \(120\) and \(5850\) points) partly address that.
All programs were written by AI systems of the same family (see ``Verification of version~1.1'' at the end of the
paper). The range of validity cannot be tested by
explicit computation, and the explicit checks test the construction at smaller $b$ (Remark~\ref{rem:38exact}). The
changes of version~1.1 in Section~\ref{sec:38}
(the parameters, the tiers and Section~\ref{ssec:refine})
received two independent blank-slate AI reviews, which found them correct with minor fixes; the fixes are included.
Sections~\ref{ssec:sawtooth} and~\ref{ssec:liftT2} (the tiers marked $\ddagger$) received two further blank-slate AI
reviews: one found Lemmas~\ref{lem:SL}, \ref{lem:T1}, \ref{lem:W4} and~\ref{lem:T2} correct and
Corollary~\ref{cor:b35c} and the tiers $\ddagger$ correct after fixes (included), the other found all of it correct.
Theorem~\ref{thm:cert} consists of the structural part of Section~\ref{sec:ends} and a finite part (explicit end
packings), checked in exact rational arithmetic by the program that produced it and by checkers written independently
from a description of the construction (Section~\ref{sec:cert}). The proofs were written and checked with AI
assistance and have not been reviewed by a human expert. The paper has not been refereed.

\emph{Organization.} Section~\ref{sec:prelim} proves \eqref{eq:translation} and the reduction used throughout.
Section~\ref{sec:block} treats the tilted block, Section~\ref{sec:lshape} the L-shaped packing, and
Section~\ref{sec:count} the count and the proofs of Theorems~\ref{thm:A}, \ref{thm:sqrt} and Proposition~\ref{prop:D}.
Section~\ref{sec:ends} lifts the rows and proves the waste identity, Section~\ref{sec:stair} constructs the staircase
end and proves Theorem~\ref{thm:main}, Section~\ref{sec:38} constructs the chained wall fillers and proves
Theorem~\ref{thm:main38}, and Section~\ref{sec:cert} proves Theorem~\ref{thm:cert}.
Sections~\ref{sec:lit} and~\ref{sec:lower} compare with the literature, and Section~\ref{sec:open} collects remarks on
other end constructions and open problems.

\section{The reduction and the wasted area}\label{sec:prelim}

For $x\ge1$ let $M(x)$ be the largest number of unit squares that can be packed in a square of side $x$, and
$W(x):=x^2-M(x)$. A packing in a square of side $x$ is also a packing in a larger square containing it, so $M$ is
non-decreasing; it is integer-valued and $M(x)\le x^2$ by area.

\begin{lemma}\label{lem:translation}
Let $k\ge2$ be an integer and $m(k):=\max_{1\le x<k}M(x)$. Then $\cs(k)=k^2-m(k)-1$, and $W(x)\to k^2-m(k)$ as
$x\to k^-$. In particular \eqref{eq:translation} holds.
\end{lemma}

\begin{proof}
The maximum exists because $M$ is non-decreasing, integer-valued and bounded by $k^2$ on $[1,k)$; and $M(x)=m(k)$ on
some interval $[x_1,k)$, so $W(x)=x^2-m(k)\to k^2-m(k)$. Let $1\le n\le k^2$. If $s(n)<k$, then by the definition of the
infimum $n$ unit squares can be packed in a square of some side $x<k$, so $n\le M(x)\le m(k)$. Conversely, if $n\le
m(k)=M(x)$ with $x<k$, then $s(n)\le x<k$. Since $s(n)\le k$ for $n\le k^2$, we get, for $0\le c\le k^2-1$,
\[
s(k^2-c)=k\iff k^2-c>m(k)\iff c\le k^2-m(k)-1 .
\]
As $1\le m(k)<k^2$ (one square fits in a square of side $1<k$, and $M(x)\le x^2<k^2$), the largest such $c$ is
$k^2-m(k)-1$.
\end{proof}

\begin{lemma}[reduction]\label{lem:reduction}
If $N\ge1$ unit squares can be packed in a square of side $S<k$, then $\cs(k)\le k^2-N-1$.
\end{lemma}

\begin{proof}
By assumption $N\le M(S)\le m(k)$, and Lemma~\ref{lem:translation} gives $\cs(k)=k^2-m(k)-1\le k^2-N-1$.
\end{proof}

This is the source of the ``$-1$'' in Theorems~\ref{thm:A} and~\ref{thm:cert}: a packing of $N$ squares in a container of
side less than $k$ shows $\cs(k)\le c-1$ with $c:=k^2-N$. In the notation of [Q] and [T], $m(k)$ is also the largest
number of unit squares in $[0,k]^2$ that are pairwise disjoint as closed sets, and $\cs(k)=\Wmin(k)-1$ with
$\Wmin(k):=k^2-m(k)$ ([Q, Section~1]).

\section{The tilted block}\label{sec:block}

Fix an integer $b\ge2$ and an angle $\theta\in(0,\pi/2)$, and write $\gamma:=\cos\theta$, $\sigma:=\sin\theta$,
$u:=(\gamma,\sigma)$ and $v:=(-\sigma,\gamma)$. For $P\in\R^2$ and $0\le i\le b-1$ let
\[
Q_i(P):=\{P+\alpha u+\beta v:\ \alpha\in[i,i+1],\ \beta\in[0,1]\},\qquad B(P):=\bigcup_{i=0}^{b-1}Q_i(P).
\]
Each $Q_i(P)$ is a unit square, and $B(P)$ is a $b\times1$ rectangle tilted by $\theta$, with lowest vertex $P$. Put
\[
f(\theta):=b\gamma+\sigma\ \ (\text{width}),\qquad h(\theta):=b\sigma+\gamma\ \ (\text{height}),\qquad
d(\theta):=\sec\theta\ \ (\text{row spacing}).
\]

\begin{lemma}[projections]\label{lem:proj}
The projection of $B(P)$ to the $x$-axis is $[P_x-\sigma,\,P_x+b\gamma]$, of length $f(\theta)$, and its projection to the
$y$-axis is $[P_y,\,P_y+b\sigma+\gamma]$, of length $h(\theta)$.
\end{lemma}

\begin{proof}
The coordinates $x=P_x+\alpha\gamma-\beta\sigma$ and $y=P_y+\alpha\sigma+\beta\gamma$ are linear on
$[0,b]\times[0,1]\ni(\alpha,\beta)$, and $\gamma,\sigma>0$. So $x$ is smallest at $(0,1)$ and largest at $(b,0)$, and $y$
is smallest at $(0,0)$ and largest at $(b,1)$.
\end{proof}

\begin{lemma}[width]\label{lem:width}
Let $\theta_0:=2\arctan(1/b)$. Then $0<\theta_0<\pi/2$,
\[
\cos\theta_0=\gamma_0=\frac{b^2-1}{D},\quad \sin\theta_0=\frac{2b}{D},\quad f(\theta_0)=b,\quad
h(\theta_0)=h_0=\frac{3b^2-1}{D},\quad d(\theta_0)=d_0:=\frac{D}{b^2-1},
\]
with $D=b^2+1$, and $f(\theta)<b$ for every $\theta\in(\theta_0,\pi/2]$.
\end{lemma}

\begin{proof}
With $t=\tan(\theta_0/2)=1/b$ the half-angle formulas give $\cos\theta_0=(1-t^2)/(1+t^2)=(b^2-1)/D$ and
$\sin\theta_0=2t/(1+t^2)=2b/D$; hence $f(\theta_0)=(b(b^2-1)+2b)/D=b$ and $h(\theta_0)=(2b^2+b^2-1)/D$. As $b\ge2$,
$\arctan(1/b)<\pi/4$, so $\theta_0<\pi/2$. With $\varphi:=\arctan(1/b)$ we have $f(\theta)=\sqrt{b^2+1}\cos(\theta-\varphi)$.
For $\theta\in[\theta_0,\pi/2]$, $\theta-\varphi\in[\varphi,\pi/2-\varphi]\subset[0,\pi]$, where $\cos$ is strictly
decreasing; so $f(\theta)<f(\theta_0)=b$ for $\theta>\theta_0$.
\end{proof}

\begin{lemma}[stacking]\label{lem:stack}
Let $P_0\in\R^2$ and $P_j:=P_0+(0,\,j\sec\theta)$ for $j\in\Z$. If $(i,j)\ne(i',j')$, then $Q_i(P_j)$ and
$Q_{i'}(P_{j'})$ have disjoint interiors. More precisely, in the coordinates $\alpha'=u\cdot(z-P_0)$,
$\beta'=v\cdot(z-P_0)$,
\[
Q_i(P_j)=[\,i+j\tan\theta,\ i+1+j\tan\theta\,]\times[\,j,\ j+1\,].
\]
\end{lemma}

\begin{proof}
The map $z\mapsto(u\cdot(z-P_0),v\cdot(z-P_0))$ is a rigid motion. Since $v\cdot(0,j\sec\theta)=j$ and
$u\cdot(0,j\sec\theta)=j\tan\theta$, the square $Q_i(P_j)$ has the stated form. If $j\ne j'$ the $\beta'$-intervals have
disjoint interiors; if $j=j'$ and $i\ne i'$ the $\alpha'$-intervals do.
\end{proof}

Thus consecutive rows touch along their long sides: the $v$-component of the shift $(0,\sec\theta)$ is exactly $1$.

\begin{lemma}[one band]\label{lem:band}
Let $a\in\R$, $L\ge0$, $\theta\in(0,\pi/2)$ and an integer $R\ge0$ with $R=0$ or $(R-1)d(\theta)+h(\theta)\le L$. Then the
$bR$ squares $Q_i(P_j)$, $0\le i<b$, $0\le j<R$, with $P_j=(a+\sigma,\,j\sec\theta)$, have pairwise disjoint interiors
and lie in $[a,\,a+f(\theta)]\times[0,L]$.
\end{lemma}

\begin{proof}
By Lemma~\ref{lem:proj} the row $B(P_j)$ has $x$-projection $[a,a+b\gamma+\sigma]$ and $y$-projection
$[jd,jd+h]\subset[0,(R-1)d+h]\subset[0,L]$. Disjointness is Lemma~\ref{lem:stack}.
\end{proof}

\section{The L-shaped packing}\label{sec:lshape}

Let $k\ge2$ and $2\le b\le k$ be integers and $\theta\in(\theta_0,\pi/2)$. By Lemma~\ref{lem:width},
\[
\varepsilon:=b-f(\theta)>0,\qquad S:=k-\varepsilon=k-b+f(\theta)<k .
\]
Let $R_1,R_2\ge0$ be integers satisfying the \emph{length conditions}
\[
\text{(L1)}\ \ R_1=0\ \text{ or }\ (R_1-1)d(\theta)+h(\theta)\le S,\qquad
\text{(L2)}\ \ R_2=0\ \text{ or }\ (R_2-1)d(\theta)+h(\theta)\le k-b .
\]
The packing consists of three parts (Figure~1 of [EG] and the Type~1 packing of [CG] have the same skeleton).
\begin{enumerate}[label=(\arabic*)]
\item The \emph{grid} $\cG$: the $(k-b)^2$ squares $[i,i+1]\times[j,j+1]$, $0\le i,j\le k-b-1$, in $G:=[0,k-b]^2$.
\item The \emph{right band} $\cA$: the $bR_1$ squares of Lemma~\ref{lem:band} with $a=k-b$, $L=S$, $R=R_1$; they lie in
$A:=[k-b,S]\times[0,S]$, since $a+f(\theta)=S$.
\item The \emph{top band} $\cT$: the image under the reflection $\iota(x,y):=(y,x)$ of the $bR_2$ squares of
Lemma~\ref{lem:band} with $a=k-b$, $L=k-b$, $R=R_2$. Before reflection they lie in $[k-b,S]\times[0,k-b]$, so $\cT$ lies
in $T:=[0,k-b]\times[k-b,S]$.
\end{enumerate}

\begin{lemma}\label{lem:valid}
The $N=(k-b)^2+b(R_1+R_2)$ squares of $\cG\cup\cA\cup\cT$ form a packing in $[0,S]^2$, and $S<k$.
\end{lemma}

\begin{proof}
Since $f(\theta)>0$ we have $k-b<S$, so $G,A,T\subset[0,S]^2$; and $S<k$. The interiors $(0,k-b)^2$,
$(k-b,S)\times(0,S)$ and $(0,k-b)\times(k-b,S)$ of $G$, $A$, $T$ are pairwise disjoint ($G$ and $A$ are separated by
$x=k-b$, $G$ and $T$ by $y=k-b$, and $A$ and $T$ by $x=k-b$); the corner $[k-b,S]^2$ belongs to $A$ only. The interior of
a square contained in one of these regions lies in the interior of that region, so squares from different parts have
disjoint interiors. Inside $\cG$ this is clear, inside $\cA$ it is Lemma~\ref{lem:band}, and inside $\cT$ it is
Lemma~\ref{lem:band} together with the fact that $\iota$ is an isometry.
\end{proof}

If $k=b$, the grid is empty, $T$ is degenerate and $R_2=0$; the lemma remains true.

\subsection{An explicit angle}
For an integer $L$ let $R(L)$ be as in Section~\ref{sec:intro}. Then $R(L)$ is the number of integers $j\ge0$ with
$jd_0+h_0<L$: if $L\le h_0$ both are $0$; otherwise the condition is $j<(L-h_0)/d_0=(L-h_0)\gamma_0$, and the number of
such $j\ge0$ is $\lceil(L-h_0)\gamma_0\rceil$.

\begin{lemma}\label{lem:angle}
Let $M:=14kb^2+1$, $p:=M$, $q:=bM-1$ and $\theta:=2\arctan(p/q)$. Then $\theta\in(\theta_0,\pi/2)$, and $R_1:=R(k)$,
$R_2:=R(k-b)$ satisfy (L1) and (L2). Moreover $\cos\theta=(q^2-p^2)/(q^2+p^2)$ and $\sin\theta=2pq/(p^2+q^2)$ are
rational, so all vertices of the packing have rational coordinates.
\end{lemma}

\begin{proof}
(i) We have $p/q>1/b$ (as $bM>bM-1$) and $p/q<1$ (as $(b-1)M>1$), so $\theta_0<\theta<\pi/2$. Since the derivative of
$\arctan$ is at most $1$,
\[
\theta-\theta_0=2\Bigl(\arctan\frac pq-\arctan\frac1b\Bigr)\le2\Bigl(\frac{M}{bM-1}-\frac1b\Bigr)=\frac{2}{b(bM-1)}
\le\frac{1}{7kb^4}=:\tau ,
\]
because $bM-1=14kb^3+b-1\ge14kb^3$.

(ii) \emph{Slack.} Let $L$ be an integer with $R:=R(L)\ge1$, and $g(L):=L-h_0-(R-1)d_0$. Then $g(L)>0$ by the
characterization of $R(L)$ above, and
\[
(b^2-1)D\,g(L)=(b^4-1)L-(b^2-1)(3b^2-1)-(R-1)D^2
\]
is an integer; hence $g(L)\ge1/(b^4-1)$.

(iii) \emph{Derivatives.} For $b\ge2$, $\gamma_0\ge3/5$, and $\tau\le1/224$; so $\cos\ge3/5-\tau\ge1/2$ on
$[\theta_0,\theta_0+\tau]$. There $d'=\sin/\cos^2\le4$, $h'=b\cos-\sin\le b$ and $|f'|=|\cos-b\sin|\le\sqrt{b^2+1}\le b+1$.

(iv) If $R_1=R(k)\ge1$ (otherwise there is nothing to prove), then by the mean value theorem, using $f(\theta_0)=b$,
\begin{align*}
(R_1-1)d(\theta)+h(\theta)-S&=\bigl[(R_1-1)d_0+h_0-k\bigr]+(R_1-1)(d(\theta)-d_0)\\
&\qquad+(h(\theta)-h_0)+(f(\theta_0)-f(\theta))\\
&\le-g(k)+(\theta-\theta_0)\bigl(4(R_1-1)+2b+1\bigr).
\end{align*}
Here $R_1-1<(k-h_0)\gamma_0<k$ and $b\le k$, so the last factor is at most $7k$, and the right side is at most
$-1/(b^4-1)+7k\tau=-1/(b^4-1)+1/b^4<0$. This is (L1). Similarly, if $R_2\ge1$,
$(R_2-1)d(\theta)+h(\theta)-(k-b)\le-g(k-b)+\tau(4k+b)\le-1/(b^4-1)+5/(7b^4)<0$, which is (L2).
\end{proof}

\begin{lemma}[maximality within the family]\label{lem:max}
If $\theta\in(\theta_0,\arctan b)$, then every pair $R_1,R_2$ satisfying (L1) and (L2) has $R_1\le R(k)$ and
$R_2\le R(k-b)$. For $\theta\in(\theta_0,\theta_0+\tau]$ equality is possible.
\end{lemma}

\begin{proof}
Note $\tan\theta_0=2b/(b^2-1)<b$, so the interval is non-empty. On it $d'>0$ and $h'=b\cos\theta-\sin\theta>0$, so $d$ and
$h$ are strictly increasing. If $R_1\ge R(k)+1$, then $(R_1-1)d_0+h_0\ge k$ by the characterization of $R(k)$, hence
$(R_1-1)d(\theta)+h(\theta)>k>S$, contradicting (L1). The top band is the same with $k-b$ in place of $S$ and $k$. The
proof of Lemma~\ref{lem:angle} works for every $\theta\in(\theta_0,\theta_0+\tau]$.
\end{proof}

So $N(k,b)$ is the exact number of squares of the packings of this section with $\theta$ close to $\theta_0$ and with the
rows of each band starting at the bottom of the band. Remark~\ref{rem:extra} shows that the family can be enlarged.

\section{The count: proofs of the bounds of order \texorpdfstring{$\sqrt k$}{sqrt k}}\label{sec:count}

\begin{proof}[Proof of Theorem~\ref{thm:A}]
By Lemmas~\ref{lem:angle} and~\ref{lem:valid}, $N(k,b)$ unit squares can be packed in a square of side $S<k$. If $k>b$
then $N(k,b)\ge(k-b)^2\ge1$; if $k=b$ then $N(k,b)=bR(k)$, and $R(k)=0$ only if $k\le h_0<3$, that is $k=b=2$. So
$N(k,b)\ge1$ for $(k,b)\ne(2,2)$, and Lemma~\ref{lem:reduction} gives $\cs(k)\le k^2-N(k,b)-1$; for $(k,b)=(2,2)$ this
bound is $\cs(2)\le3$, which holds by definition.

For the second inequality we compute the count exactly. Put
\[
\eta:=\frac{2k-10}{D}+\frac{8}{D^2},\qquad \eta':=\eta-\frac{2b}{D}.
\]
We claim
\begin{equation}\label{eq:exact}
k^2-N(k,b)\ \le\ 6b+b\bigl(\lfloor\eta\rfloor+\lfloor\eta'\rfloor\bigr),\qquad\text{with equality if }k-b\ge2.
\end{equation}
For an integer $L$, using $h_0=3-4/D$ and $\gamma_0=1-2/D$,
\[
(L-h_0)\gamma_0=L-3+\frac4D-\frac{2L}{D}+\frac6D-\frac{8}{D^2}=L-3-\eta_L,\qquad \eta_L:=\frac{2L-10}{D}+\frac{8}{D^2}.
\]
Since $L$ is an integer, $\lceil L-3-\eta_L\rceil=L-3-\lfloor\eta_L\rfloor$, so
$R(L)=\max(0,L-3-\lfloor\eta_L\rfloor)\ge L-3-\lfloor\eta_L\rfloor$. Here $\eta_k=\eta$ and $\eta_{k-b}=\eta'$. Hence
\[
k^2-N(k,b)=2kb-b^2-b\bigl(R(k)+R(k-b)\bigr)\le2kb-b^2-b\bigl(2k-b-6-\lfloor\eta\rfloor-\lfloor\eta'\rfloor\bigr)
=6b+b\bigl(\lfloor\eta\rfloor+\lfloor\eta'\rfloor\bigr).
\]
If $L\ge2$, then $L-h_0>-1$ and $0<\gamma_0<1$, so $(L-h_0)\gamma_0>-1$ and the ceiling is at least $0$; the maximum is
then not active. This applies to $L=k$ always and to $L=k-b$ when $k-b\ge2$, which gives the equality case.

Now $\lfloor\eta\rfloor+\lfloor\eta'\rfloor\le2\eta-2b/D$, and therefore
\[
k^2-N(k,b)\le6b+\frac{(4k-20)b}{D}+\frac{16b}{D^2}-\frac{2b^2}{D}
=6b+\frac{4kb}{D}-2+\frac{2-20b+16b/D}{D}.
\]
Here $4kb/D<4k/b$, and for $b\ge2$ we have $16b/D\le32/5<7$, so $2-20b+16b/D<9-20b<0$. Hence
$k^2-N(k,b)<6b+4k/b-2$, and $\cs(k)\le k^2-N(k,b)-1<6b+4k/b-3$.
\end{proof}

\begin{corollary}\label{cor:B}
For every integer $k\ge6$, $\cs(k)<4\sqrt6\,\sqrt k\approx9.798\sqrt k$.
\end{corollary}

\begin{proof}
The function $g(x)=6x+4k/x$ has its minimum $g(x^*)=4\sqrt6\sqrt k$ at $x^*=\sqrt{2k/3}$. Let $b=\lceil x^*\rceil=x^*+\delta$
with $0\le\delta<1$. Using $4k/x^{*2}=6$,
\[
g(b)-g(x^*)=6\delta-\frac{4k\delta}{x^*(x^*+\delta)}=\frac{6\delta^2}{x^*+\delta}<\frac{6}{x^*}\le3,
\]
because $x^*\ge2$ for $k\ge6$. Also $2\le b\le x^*+1\le k$. Theorem~\ref{thm:A} gives $\cs(k)<g(b)-3<4\sqrt6\sqrt k$.
\end{proof}

\begin{proof}[Proof of Theorem~\ref{thm:sqrt}]
Let $b:=\lceil\sqrt{k-4}\,\rceil$. Then $b\ge\lceil\sqrt2\,\rceil=2$ and $b\le\sqrt{k-4}+1\le k$. Moreover $D=b^2+1\ge k-3>0$
and $2k-10>0$, so
\[
\eta\le\frac{2k-10}{k-3}+\frac{8}{(k-3)^2}=2-\frac{4(k-3)-8}{(k-3)^2}<2,
\]
since $k-3>2$. Hence $\lfloor\eta\rfloor\le1$ and $\lfloor\eta'\rfloor\le\lfloor\eta\rfloor\le1$, and \eqref{eq:exact} gives
$k^2-N(k,b)\le8b$. Since $(k,b)\ne(2,2)$, Theorem~\ref{thm:A} gives $\cs(k)\le8b-1<8\sqrt{k-4}+7$.
\end{proof}

For $k>15$ Theorem~\ref{thm:sqrt} is stronger than Corollary~\ref{cor:B}, since $8\sqrt k+7<4\sqrt6\sqrt k$ for
$\sqrt k>7/(4\sqrt6-8)\approx3.89$.

\begin{proof}[Proof of Proposition~\ref{prop:D}]
If $k-b\ge2$, equality holds in \eqref{eq:exact}, and $\lfloor z\rfloor>z-1$ gives
\[
k^2-N(k,b)>6b+b\Bigl(2\eta-\frac{2b}{D}-2\Bigr)=4b+\frac{(4k-20)b}{D}+\frac{16b}{D^2}-\frac{2b^2}{D}
\ge4b+\frac{4kb}{D}-12,
\]
using $16b/D^2\ge0$, $2b^2/D\le2$ and $20b/D\le10$. With $x:=b+1/b$ we have $kb/D=k/x$ and $4b=4x-4/b\ge4x-2$, so
$k^2-N(k,b)>4x+4k/x-14\ge8\sqrt k-14$ by the inequality of arithmetic and geometric means. If $b\in\{k-1,k\}$, then
$R(k-b)\in\{R(0),R(1)\}=\{0\}$ (as $h_0\ge11/5$) and $R(k)\le\lceil k-h_0\rceil\le k-2$, so $N(k,b)\le1+k(k-2)$ and
$k^2-N(k,b)\ge2k-1$; and $2k-1-(8\sqrt k-14)=2(\sqrt k-2)^2+5>0$.
\end{proof}

Hence the best bound obtainable from Theorem~\ref{thm:A} satisfies
$8\sqrt k-15<\min_b\bigl(k^2-N(k,b)-1\bigr)\le8\lceil\sqrt{k-4}\,\rceil-1$.

\begin{remark}[where $8\sqrt k$ comes from]\label{rem:eight}
The exact count \eqref{eq:exact} is $6b+b(\lfloor\eta\rfloor+\lfloor\eta'\rfloor)$ with $\eta\approx2k/b^2$. The
estimate $\lfloor z\rfloor\le z$ gives $6b+4k/b$, minimized at about $9.80\sqrt k$; but for $b=\lceil\sqrt{k-4}\rceil$ and
$k\ge23$ both floors equal $1$ while $\eta$ is just below $2$, and the count is $8b$. (For $k\ge31$: $\eta<2$ by the
proof of Theorem~\ref{thm:sqrt}, and with $s:=\sqrt{k-4}$, $b<s+1$ gives $b^2+2b<k-1+4s$, so
$D\eta'>2k-10-2b>b^2+1=D$ as soon as $k-10\ge4s$, which holds for $k\ge31$; hence $1<\eta'\le\eta<2$. For
$23\le k\le30$ this was checked exactly, as was the whole range $23\le k\le2\cdot10^5$; $k=22$ fails. Nothing
else in the paper uses this.) Geometrically (Section~\ref{sec:ends}),
each band loses about $2L/b$ along its length $L$, because the row spacing $\sec\theta$ exceeds $1$ by about $2/b^2$,
and about $b$ at each of its two ends, where the first and the last row cross the band with a rise of about $2$. The
total $4k/b+4b$ is at least $8\sqrt k$.
\end{remark}

\begin{remark}[a larger family]\label{rem:extra}
Lemma~\ref{lem:max} concerns the family in which the top band stays in $x\le k-b$. One more row in the top band
protrudes into the corner $[k-b,S]^2$, but often fits into the empty wedge above the last row of the right band. An
exact computer check (all $6\le k\le30$ with $k-b\ge2$) found $285$ of $375$ such enlarged packings valid; for example
$k=12$, $b=3$ gives $126$ squares instead of $N(12,3)=123$. We have no general proof, and this remark is not used.
\end{remark}

\section{Lifted rows and the waste identity}\label{sec:ends}

\subsection{Lifted rows and end regions}\label{ssec:frame}
Throughout this section $b\ge2$, $\theta\in(\theta_0,\pi/2)$, $\gamma=\cos\theta$, $\sigma=\sin\theta$, and
$w:=f(\theta)=b\gamma+\sigma<b$, $h:=h(\theta)$, $d:=\sec\theta$. A band is described in its own \emph{strip frame}: the
strip is $\Sigma:=[0,w]\times[0,L_b]$, where $L_b>0$ is the length of the band. Given a \emph{lift} $y_0\ge0$ and an
integer $R\ge1$, the rows are $B(P_j)$, $0\le j<R$, with
\[
P_j:=(\sigma,\ y_0+jd),\qquad\text{and we put}\qquad y_1:=L_b-h-(R-1)d-y_0 .
\]
We assume $y_1\ge0$. For $y\ge0$ the \emph{end region} is
\[
\Reg(y):=\{(x',y'):\ 0\le x'\le w,\ \ 0\le y'\le y+(x'-\sigma)\tan\theta\},
\]
and $\rho(x',y'):=(w-x',\,L_b-y')$ is the rotation by $\pi$ about the centre of $\Sigma$. The upper side of $\Reg(y_0)$
is the line through the lower side of the first row $B(P_0)$; on $[0,\sigma]$ the first row is higher than this line,
so $\Reg(y_0)$ is slightly smaller than the part of $\Sigma$ below the first row.

\begin{lemma}\label{lem:ends}
Let $\cC_0$ and $\cC_1$ be finite families of unit squares with pairwise disjoint interiors, contained in $\Reg(y_0)$ and
$\Reg(y_1)$ respectively. Then the $bR$ squares $Q_i(P_j)$ ($0\le i<b$, $0\le j<R$), the squares of $\cC_0$ and the
squares of $\rho(\cC_1)$ form a packing in $\Sigma$.
\end{lemma}

\begin{proof}
\emph{Containment.} By Lemma~\ref{lem:proj} the row $B(P_j)$ has $x$-projection $[0,w]$ and $y$-projection
$[y_0+jd,\,y_0+jd+h]\subset[y_0,\,L_b-y_1]\subset[0,L_b]$. On $\Reg(y)$ we have
$y'\le y+(w-\sigma)\tan\theta=y+b\sigma\le y+h$; since $L_b=y_0+y_1+(R-1)d+h$ is at least $y_0+h$ and at least
$y_1+h$, both end regions lie in $\Sigma$, and $\rho(\Sigma)=\Sigma$.

\emph{Separation.} Let $\beta'(z):=v\cdot(z-P_0)$. By Lemma~\ref{lem:stack} (with this $P_0$), every row square lies in
$\{0\le\beta'\le R\}$. The upper side of $\Reg(y_0)$ is the line through $P_0$ with direction $u$, and
$\beta'(x',y')=-\sigma(x'-\sigma)+\gamma(y'-y_0)\le0$ is equivalent to $y'\le y_0+(x'-\sigma)\tan\theta$; so
$\Reg(y_0)\subset\{\beta'\le0\}$. Next, the highest vertex of the last row is
$P_{R-1}+bu+v=(b\gamma,\ y_0+(R-1)d+h)$, and
\begin{equation}\label{eq:rot}
\rho\bigl(P_{R-1}+bu+v\bigr)=(w-b\gamma,\ L_b-y_0-(R-1)d-h)=(\sigma,\ y_1)=:P_0'.
\end{equation}
As above, $\Reg(y_1)\subset\{z:\ v\cdot(z-P_0')\le0\}$. The linear part of $\rho$ is $-\mathrm{id}$, so for
$z\in\Reg(y_1)$
\[
\beta'(\rho(z))=v\cdot\bigl(\rho(P_0')-P_0\bigr)-v\cdot(z-P_0')=v\cdot(P_{R-1}+bu+v-P_0)-v\cdot(z-P_0')\ge(R-1)+1=R,
\]
using $v\cdot(P_{R-1}-P_0)=R-1$, $v\cdot u=0$ and $v\cdot v=1$. Hence $\rho(\Reg(y_1))\subset\{\beta'\ge R\}$. The three
sets $\{\beta'\le0\}$, $\{0\le\beta'\le R\}$ and $\{\beta'\ge R\}$ have pairwise disjoint interiors; inside each group the
interiors are disjoint by Lemma~\ref{lem:stack} and by assumption ($\rho$ is an isometry).
\end{proof}

\begin{corollary}\label{cor:assembly}
Let $k>b$ be an integer, $S:=k-b+w$, and $y_0\ge0$. For the right band put $L_b:=S$ and for the top band $L_b:=k-b$, and
suppose that for each band an integer $R\ge1$ is given with $y_1\ge0$. Let $\cC_0,\cC_1$ be as in Lemma~\ref{lem:ends}
for each band. Map the right band into the plane by $(x',y')\mapsto(k-b+x',y')$, and the top band by
$(x',y')\mapsto\iota(k-b+x',y')=(y',k-b+x')$. Together with the grid $\cG$ these squares form a packing in $[0,S]^2$, with
$S<k$ and
\[
N=(k-b)^2+b\,(R_{\rm right}+R_{\rm top})+\sum_{\text{4 ends}}|\cC| .
\]
\end{corollary}

\begin{proof}
The strips are mapped onto $A$ and $T$ of Section~\ref{sec:lshape}. Lemma~\ref{lem:ends} gives a packing inside each of
$A$ and $T$, and the separation of $G$, $A$, $T$ is as in the proof of Lemma~\ref{lem:valid}.
\end{proof}

\subsection{The waste identity}
For $y\ge0$ let
\[
A_0(y):=yw+\tfrac12\gamma\sigma(1+b^2).
\]
If $y=y_0$, this is the area of the part of $\Sigma$ below the lower boundary of the first row: that boundary is the
segment from $(0,y_0+\gamma)$ to $(\sigma,y_0)$ followed by the segment from $(\sigma,y_0)$ to $(w,y_0+b\sigma)$, and
\[
\int_0^\sigma\Bigl(y_0+\gamma-\frac{\gamma}{\sigma}x\Bigr)dx+\int_\sigma^w\bigl(y_0+(x-\sigma)\tan\theta\bigr)dx
=y_0\sigma+\frac{\gamma\sigma}{2}+y_0b\gamma+\frac{b^2\gamma\sigma}{2}=A_0(y_0).
\]
For a finite family $\cC$ in $\Reg(y)$ we write $E(y):=A_0(y)-|\cC|$, the \emph{end cost}. It is the uncovered area of
the part of $\Sigma$ below the first row (for $y=y_0$), resp.\ above the last row (for $y=y_1$, after applying $\rho$).

\begin{proposition}[waste identity]\label{prop:waste}
In the situation of Lemma~\ref{lem:ends}, with $E(y_0)=A_0(y_0)-|\cC_0|$ and $E(y_1)=A_0(y_1)-|\cC_1|$,
\[
wL_b-bR-|\cC_0|-|\cC_1|=(R-1)\tan\theta+E(y_0)+E(y_1).
\]
In the situation of Corollary~\ref{cor:assembly},
\[
k^2-N=(k^2-S^2)+\sum_{\text{2 bands}}\bigl[(R-1)\tan\theta+E(y_0)+E(y_1)\bigr].
\]
\end{proposition}

\begin{proof}
Split $\Sigma$ into the part below the lower boundary of $B(P_0)$ (area $A_0(y_0)$), the part between the lower boundary
of $B(P_0)$ and the upper boundary of $B(P_{R-1})$, and the part above the upper boundary of $B(P_{R-1})$. The lower
boundary of $B(P_j)$ is that of $B(P_0)$ shifted up by $jd$, and each row has $x$-projection $[0,w]$; so the middle part
has area $(R-1)dw+b$, where $b$ is the area of $B(P_{R-1})$. By \eqref{eq:rot}, $\rho$ maps $B(P_{R-1})$ onto the
rectangle $B(P_0')$ (a rotation by $\pi$ maps a rectangle with sides along $u$ and $v$ onto such a rectangle, and the
highest vertex onto the lowest one), so the top part has area $A_0(y_1)$. Hence
$wL_b=A_0(y_0)+A_0(y_1)+(R-1)dw+b$, and since $dw=b+\tan\theta$,
\[
wL_b-bR-|\cC_0|-|\cC_1|=E(y_0)+E(y_1)+(R-1)(dw-b)=E(y_0)+E(y_1)+(R-1)\tan\theta .
\]
For the second identity, $[0,S]^2$ is the union of $G$, $A$ and $T$ with disjoint interiors, $G$ is covered by the grid,
and $A$, $T$ are the images of the two strips.
\end{proof}

The term $(R-1)\tan\theta$ is the \emph{body cost} of a band. Since $\tan\theta\approx2/b$ and $R\approx L_b$, the two
bands cost about $4k/b$; the four ends cost $4E$, and $k^2-S^2=(b-w)(k+S)$ is tiny when $\theta$ is close to $\theta_0$.
For the packing of Theorem~\ref{thm:A} ($y_0=0$, no end squares) $E(0)=\frac12\gamma\sigma(1+b^2)\approx b$; this is the
end cost $\approx b$ of Remark~\ref{rem:eight}. Any packing of $\Reg(y)$ with end cost $E=o(b)$, for lifts
$y$ in an interval of length $\sec\theta$, makes the total cost $4k/b+o(b)$; Section~\ref{sec:stair} constructs such
packings with $E=O(b^{2/3})$.

\section{The staircase end and the exponent \texorpdfstring{$2/5$}{2/5}}\label{sec:stair}

In this section we prove Theorem~\ref{thm:main}. Section~\ref{ssec:stairlemma} packs a general trapezoid of the form
$\Reg(y)$ by a staircase of columns with one common tilt (Lemma~\ref{lem:stair}); Section~\ref{ssec:b23} specializes
to the band ends, and Section~\ref{ssec:proofmain} assembles the packing. The idea of replacing the lower squares of
tilted columns by horizontal layers is that of [Bui, Section~3] (Section~\ref{ssec:credit}).

\subsection{Columns, layers and the end lemma}\label{ssec:stairlemma}
Fix reals $w>0$, $\mu>0$, $\sigma_0\ge0$ and $y$ with $g_0:=y-\sigma_0\mu>1$, and put
\[
g(X):=y+(X-\sigma_0)\mu,\qquad \Reg:=\{(X,Y):\ 0\le X\le w,\ 0\le Y\le g(X)\},
\]
$\bar g:=g(w)$ and $\Delta:=\bar g-g_0=w\mu$.
$\Reg$ is a convex trapezoid of area $\int_0^wg$, and $g$ is increasing with $g(0)=g_0$. For a packing in $\Reg$ let
$U$ be the area of $\Reg$ minus the number of unit squares of the packing (the \emph{uncovered area}).

\emph{Tilt notation.} For $t\in(0,1)$ let $\alpha:=2\arctan t$, so that
\[
\ca=\frac{1-t^2}{1+t^2},\quad \sa=\frac{2t}{1+t^2},\quad \ta=\frac{2t}{1-t^2},\quad \se=\frac{1+t^2}{1-t^2}
\]
are rational if $t$ is. Put $u_\alpha:=(\ca,\sa)$, $v_\alpha:=(-\sa,\ca)$ and $\lambda(X,Y):=X+Y\ta$ (the \emph{slab
coordinate}). We use $\se\ca=1$, $\se\sa=\ta$, $\sa\le\ta$, $\sa\le2t$, $1-\ca=2t^2/(1+t^2)\le2t^2$, and that
$\ta$ and $\se$ are increasing in $t$.

\emph{Columns.} For a real $\xi$ and integers $J\ge0$, $n\ge1$ put $P:=(\xi-J\ta,\,J)$ and
\[
K(\xi,J,n):=\{P+a\,u_\alpha+e\,v_\alpha:\ a\in[0,1],\ e\in[0,n]\},
\]
the union of the $n$ unit squares $\{P+a\,u_\alpha+e\,v_\alpha:\ a\in[0,1],\ e\in[i,i+1]\}$, $0\le i<n$: a $1\times n$
rectangle tilted by $\alpha$ and leaning to the left, with lowest vertex $P$.

\begin{lemma}[columns]\label{lem:col}
Every point $(X,Y)$ of $K(\xi,J,n)$ satisfies
\begin{enumerate}[label=(\roman*)]
\item $\xi\le\lambda(X,Y)\le\xi+\se$;
\item $J\le Y\le J+n\ca+\sa$;
\item $X\ge\xi-J\ta-n\sa\ge\xi-(J+n)\ta$;
\item $X\le\xi-J\ta+\ca\le\xi+\se$.
\end{enumerate}
\end{lemma}

\begin{proof}
For $z=P+a\,u_\alpha+e\,v_\alpha$ we have $X=\xi-J\ta+a\ca-e\sa$ and $Y=J+a\sa+e\ca$. Hence
$\lambda=\xi-J\ta+a\ca-e\sa+J\ta+a\sa\ta+e\ca\ta=\xi+a(\ca+\sa^2/\ca)=\xi+a\se$, which is (i). (ii), (iii) and (iv)
follow from $a\in[0,1]$, $e\in[0,n]$, $\sa,\ca\ge0$, $\sa\le\ta$ and $\ca\le1\le\se$.
\end{proof}

\begin{lemma}[the tilt]\label{lem:tiltS}
Let $m$ be an integer and $g_0$ a real number with $1<g_0<m$. Put $D_m:=m+g_0$, $\delta:=m-g_0>0$,
$F(t):=D_mt^2-2t-\delta$ and
\[
t^*:=\frac{1+\sqrt{1+D_m\delta}}{D_m}.
\]
Then $t^*\in(0,1)$, and for $t\in(0,1)$: $m\ca+\sa\le g_0$ if and only if $F(t)\ge0$, if and only if $t\ge t^*$.
Moreover, if $g_0\ge\nu>0$ and $\delta\le\bar\delta$, then $t^*\le\bigl(1+\sqrt{1+2\nu\bar\delta}\bigr)/(2\nu)$.
\end{lemma}

\begin{proof}
With $\ca,\sa$ as above, $(1+t^2)(m\ca+\sa-g_0)=m(1-t^2)+2t-g_0(1+t^2)=-F(t)$. $F$ is a convex quadratic with
$F(0)=-\delta<0$, so $F<0$ on $[0,t^*)$ and $F>0$ on $(t^*,\infty)$, where $t^*$ is its positive root; $F(1)=2g_0-2>0$
gives $t^*<1$. For the bound, $f(D',e):=\bigl(1+\sqrt{1+D'e}\bigr)/D'$ is increasing in $e\ge0$ and decreasing in $D'>0$,
because
\[
\frac{\partial f}{\partial D'}=\frac{D'e/(2\sqrt{1+D'e})-1-\sqrt{1+D'e}}{D'^2}\quad\text{and}\quad
\frac{D'e}{2\sqrt{1+D'e}}<\frac{\sqrt{1+D'e}}2 .
\]
Since $D_m=m+g_0>2g_0\ge2\nu$ and $\delta\le\bar\delta$, $t^*=f(D_m,\delta)\le f(2\nu,\bar\delta)$.
\end{proof}

\emph{The construction.} Let $Q\ge1$ be an integer (the tilt grid).
\begin{enumerate}[label=(S\arabic*)]
\item \emph{Column length:} $m:=\lfloor\bar g\rfloor+1$. So $m>\bar g\ge g_0$, and $\delta:=m-g_0\in(\Delta,\Delta+1]$.
\item \emph{Tilt:} $t:=\lceil t^*Q\rceil/Q$ with $t^*=t^*(m,g_0)$ of Lemma~\ref{lem:tiltS}, so $t\in[t^*,t^*+1/Q]$;
we require $t<1$. Let $\alpha:=2\arctan t$; this one tilt is used for all columns.
\item \emph{Slabs:} $\xi_i:=m\ta+i\se$ for $i=0,1,2,\dots$, and $I:=\#\{i\ge0:\ \xi_i+\se\le w\}$; we require $I\ge1$,
that is $m\ta+\se\le w$. Then $w-\se<\xi_I\le w$.
\item \emph{Staircase:} for $0\le i<I$ let $\Gamma_i:=g(\xi_i-m\ta)$, $H(J):=J+(m-J)\ca+\sa$,
\[
J_i:=\Bigl\lfloor\frac{\Gamma_i-m\ca-\sa}{1-\ca}\Bigr\rfloor,\qquad n_i:=m-J_i,
\]
and the column $K_i:=K(\xi_i,J_i,n_i)$ ($n_i$ squares with base at height $J_i$).
\item \emph{Layers:} for integers $j\ge0$ let $p(j):=\#\{i<I:\ J_i\le j\}$ and $a_j:=\xi_{p(j)}-j\ta$, and let $j_1$ be
the largest $j\in\{-1,0,1,\dots,\lfloor\bar g\rfloor\}$ such that $g(a_{j'})\ge j'+1$ for all integers $0\le j'\le j$.
For $0\le j\le j_1$ put $L_j:=\lfloor w-a_j\rfloor$ and, if $L_j\ge1$, the layer
$\Lambda_j:=[a_j,a_j+L_j]\times[j,j+1]$ ($L_j$ axis-parallel unit squares in a row).
\item \emph{Left wedge:} for $0\le r\le\lfloor g_0\rfloor-1$ put $l_r:=\lfloor(m-r-1)\ta\rfloor$ and, if $l_r\ge1$, the
row $\Omega_r:=[0,l_r]\times[r,r+1]$ ($l_r$ axis-parallel unit squares).
\end{enumerate}
All columns have the same tilt and the same total length $m$: column $i$ keeps its top $n_i$ squares tilted, and its
$J_i$ lowest squares are replaced by the layers $j<J_i$. Each layer raises the top of a column by exactly $1-\ca$,
which is how a fixed tilt follows the rising ceiling ($J_0\le J_1\le\cdots$, a staircase). The left wall is handled by
the rows $\Omega_r$ in the wedge left of the first slab, the right wall by the layers, which run from the slab
boundary to the wall $X=w$.

\begin{lemma}[end lemma]\label{lem:stair}
Assume $g_0>1$, $t<1$ and $m\ta+\se\le w$. Then the columns $K_0,\dots,K_{I-1}$, the layers $\Lambda_j$ and the
wedge rows $\Omega_r$ form a packing in $\Reg$, and
\begin{multline*}
U\le B^*:=w\sa+w(1-\ca)+\Delta(\se+m\ta)+\lfloor\bar g\rfloor\bigl(1+\tfrac12\ta\bigr)+(1+\mu\omega)\,\omega\\
+\lfloor g_0\rfloor\bigl(1+\tfrac12\ta\bigr)+(\delta+1)\ta\bigl(1+\mu(\delta+1)\ta\bigr),\qquad
\omega:=\se+\bar g\ta .
\end{multline*}
\end{lemma}

The seven terms are: the triangles beside the columns; the gaps above the columns (the staircase quantum $1-\ca$);
the slope of the ceiling across one column; the ends and fractional parts of the layers; the corner at the right
wall; the wedge rows; and the corner at the left wall. No small-angle expansion is used.

\begin{proof}
(a) \emph{The staircase exists up to the right wall.} For $0\le i<I$ we have
$0\le\xi_i-m\ta=i\se\le\xi_i\le w-\se$, so $g_0\le\Gamma_i\le\bar g$, and $\Gamma_i$ is non-decreasing in $i$. $H$ is
affine in $J$ with slope $1-\ca>0$ (note $t\ge t^*>0$), so $J_i=\max\{J\in\Z:\ H(J)\le\Gamma_i\}$. Moreover
$J_i\ge0$, because $H(0)=m\ca+\sa\le g_0\le\Gamma_i$ by Lemma~\ref{lem:tiltS} ($t\ge t^*$ and $1<g_0<m$);
$J_i\le m-1$, that is $n_i\ge1$, because $H(m)=m+\sa>m>\bar g\ge\Gamma_i$; the gap $\Gamma_i-H(J_i)$ lies in
$[0,1-\ca)$, because $H(J_i+1)=H(J_i)+1-\ca>\Gamma_i$; and $J_0\le J_1\le\dots\le J_{I-1}$. So every slab up to the
right wall receives a column; the columns stop only because the width runs out, and $w-\xi_I<\se$.

(b) \emph{Columns lie in $\Reg$.} By Lemma~\ref{lem:col}(iii),(iv) and $J_i+n_i=m$, every point of $K_i$ has
$X\ge\xi_i-m\ta\ge0$ and $X\le\xi_i+\se\le w$, and by (ii) and (a),
$J_i\le Y\le H(J_i)\le\Gamma_i=g(\xi_i-m\ta)\le g(X)$.

(c) \emph{Layers lie in $\Reg$.} For $0\le j\le\lfloor\bar g\rfloor$ we have $a_j\le\xi_{p(j)}\le\xi_I\le w$ and
$a_j\ge\xi_0-j\ta=(m-j)\ta>0$. If $j\le j_1$, then $g(a_j)\ge j+1$, so every point of $\Lambda_j$ has
$0\le a_j\le X\le a_j+L_j\le w$ and $0\le Y\le j+1\le g(a_j)\le g(X)$. Also $j_1\le\lfloor\bar g\rfloor-1$: for
$j=\lfloor\bar g\rfloor$, $g(a_j)\le g(w)=\bar g<j+1$.

(d) \emph{Wedge rows lie in $\Reg$.} For $r\le\lfloor g_0\rfloor-1$, $m-r-1\ge m-\lfloor g_0\rfloor>0$, so
$0\le l_r\le(m-1)\ta\le w$, and $Y\le r+1\le\lfloor g_0\rfloor\le g_0=g(0)\le g(X)$ on $\Omega_r$.

(e) \emph{Disjointness.} Two closed convex sets on the two sides of a line have disjoint interiors; all separations
below are of this kind.
\begin{itemize}
\item Columns: for $i<i'$, $K_i\subset\{\lambda\le\xi_i+\se\}=\{\lambda\le\xi_{i+1}\}\subset\{\lambda\le\xi_{i'}\}$ and
$K_{i'}\subset\{\lambda\ge\xi_{i'}\}$ (Lemma~\ref{lem:col}(i)). No gap is needed between slabs, because the tilt never
changes.
\item Wedge rows against columns and layers: $\Omega_r\subset\{\lambda\le l_r+(r+1)\ta\le(m-r-1)\ta+(r+1)\ta=\xi_0\}$,
while the columns lie in $\{\lambda\ge\xi_0\}$ and $\Lambda_j\subset\{\lambda\ge a_j+j\ta=\xi_{p(j)}\ge\xi_0\}$ (as
$X\ge a_j$, $Y\ge j$ and $\ta>0$ on $\Lambda_j$).
\item Layer $\Lambda_j$ against column $K_i$: if $i<p(j)$, then $K_i\subset\{\lambda\le\xi_{i+1}\}\subset
\{\lambda\le\xi_{p(j)}\}$ and $\Lambda_j\subset\{\lambda\ge\xi_{p(j)}\}$. If $i\ge p(j)$, then $J_i\ge j+1$ (the $J_i$ are
non-decreasing, so $\{i:\ J_i\le j\}=\{0,\dots,p(j)-1\}$), hence $K_i\subset\{Y\ge j+1\}$ and $\Lambda_j\subset\{Y\le j+1\}$.
\item Layers among themselves, and wedge rows among themselves, lie in different rows $[j,j+1]$.
\end{itemize}

(f) \emph{A partition of $\Reg$.} Put
\[
\mathrm{Wd}:=\Reg\cap\{\lambda<\xi_0\},\quad \mathcal S_i:=\Reg\cap\{\xi_i\le\lambda<\xi_{i+1}\}\ (0\le i<I),\quad
\mathrm{Rt}:=\Reg\cap\{\lambda\ge\xi_I\},
\]
$\mathcal S_i^+:=\mathcal S_i\cap\{Y\ge J_i\}$, $\mathcal S_i^-:=\mathcal S_i\cap\{Y<J_i\}$ and
$\mathrm{Lay}:=\bigcup_i\mathcal S_i^-\cup\mathrm{Rt}$. Then $\mathrm{Wd}$, the $\mathcal S_i^+$ and $\mathrm{Lay}$
partition $\Reg$ (as $\xi_{i+1}=\xi_i+\se$, the slabs tile $[\xi_0,\xi_I)$). The interior of each placed piece lies in
one part: $\operatorname{int}K_i\subset\mathcal S_i^+$ by Lemma~\ref{lem:col}(i),(ii);
$\operatorname{int}\Omega_r\subset\mathrm{Wd}$ by (e); and $\operatorname{int}\Lambda_j\subset\mathrm{Lay}$, because a
point of $\Reg$ with $\lambda\ge\xi_{p(j)}$ and $Y<j+1$ lies either in $\mathrm{Rt}$ or in some $\mathcal S_i$ with
$i\ge p(j)$, where $J_i\ge j+1>Y$. Boundaries have measure $0$ and the pieces have disjoint interiors, so $U$ is the
sum over the parts of the area of the part minus the number of unit squares in the pieces assigned to it.

(g) \emph{The column parts.} Fix $i$ and write $J=J_i$, $n=n_i$, $H_i=H(J)$. The parallelogram
$\Pi_i:=\{\xi_i\le\lambda\le\xi_i+\se,\ J\le Y\le H_i\}$ lies in $\Reg$: its points have
$X\ge\xi_i-H_i\ta\ge\xi_i-m\ta\ge0$ (as $H_i\le\Gamma_i\le\bar g<m$), $X\le\xi_i+\se\le w$, and
$Y\le H_i\le\Gamma_i=g(\xi_i-m\ta)\le g(X)$. Its horizontal sections have length $\se$, so
$\operatorname{area}(\Pi_i)=\se(H_i-J)=\se(n\ca+\sa)=n+\ta$. The rest of $\mathcal S_i^+$ lies in $\{Y>H_i\}$; there
$X\le\lambda<\xi_i+\se$, so $Y\le g(\xi_i+\se)=\Gamma_i+\mu(\se+m\ta)$, and the sections still have length at most
$\se$. Hence, by (a),
\begin{align*}
\operatorname{area}(\mathcal S_i^+)-n_i&\le\ta+\se\bigl(\Gamma_i+\mu(\se+m\ta)-H_i\bigr)\\
&<\ta+\se(1-\ca)+\se\,\mu(\se+m\ta).
\end{align*}
Summing, with $I\se=\xi_I-\xi_0\le w$, $I\ta=I\se\sa\le w\sa$ and $w\mu=\Delta$,
\[
\sum_i\bigl[\operatorname{area}(\mathcal S_i^+)-n_i\bigr]\le w\sa+w(1-\ca)+\Delta(\se+m\ta).
\]

(h) \emph{The layer part.} For an integer $j$ let $\mathrm{Lay}_j:=\mathrm{Lay}\cap\{j\le Y<j+1\}$. By the argument at the
end of (f), $\mathrm{Lay}_j$ contains $\Reg\cap\{j\le Y<j+1,\ \lambda\ge\xi_{p(j)}\}$; conversely, a point of
$\mathcal S_i^-$ at height $Y\in[j,j+1)$ has $J_i>Y\ge j$, hence $J_i\ge j+1$, that is $i\ge p(j)$ and
$\lambda\ge\xi_{p(j)}$; and the points of $\mathrm{Rt}$ have $\lambda\ge\xi_I\ge\xi_{p(j)}$. So $\mathrm{Lay}_j=\Reg\cap\{j\le Y<j+1,\ \lambda\ge\xi_{p(j)}\}$.

\emph{Rows $0\le j\le j_1$.} The section of $\mathrm{Lay}_j$ at height $Y$ lies in $[\xi_{p(j)}-Y\ta,\,w]$, so
$\operatorname{area}(\mathrm{Lay}_j)\le\int_j^{j+1}\bigl(w-a_j+(Y-j)\ta\bigr)dY=w-a_j+\frac12\ta$, while $\Lambda_j$
covers $L_j>w-a_j-1$ (if $L_j=0$, nothing is placed and $w-a_j<1$). So
$\operatorname{area}(\mathrm{Lay}_j)-L_j<1+\frac12\ta$, and there are $j_1+1\le\lfloor\bar g\rfloor$ such rows by (c).

\emph{The rest, $\mathrm{Lay}\cap\{Y\ge j_1+1\}$.} First, $j_1\ge J_{I-1}-1$: if $0\le j\le J_{I-1}-1$ (note
$J_{I-1}-1\le m-2=\lfloor\bar g\rfloor-1$), then $p:=p(j)\le I-1$ and $J_p\ge j+1$, so
$\Gamma_p\ge H(J_p)\ge J_p\ge j+1$ (as $(m-J)\ca+\sa\ge0$); and $a_j=\xi_p-j\ta\ge\xi_p-m\ta$, so
$g(a_j)\ge\Gamma_p\ge j+1$. Hence the points of $\mathcal S_i^-$, which have $Y<J_i\le J_{I-1}\le j_1+1$, are not in
$\{Y\ge j_1+1\}$, and $\mathrm{Lay}\cap\{Y\ge j_1+1\}=\mathrm{Rt}\cap\{Y\ge j_1+1\}$. Next, $j_1+1\le\lfloor\bar g\rfloor$
by (c) and $j_1+1\ge J_{I-1}$, so $p(j_1+1)=I$ and, by the maximality of $j_1$, $g(a_{j_1+1})<j_1+2$, where
$a_{j_1+1}=\xi_I-(j_1+1)\ta\ge\xi_I-\bar g\ta>w-\se-\bar g\ta=w-\omega$. Hence
$\bar g-\mu\omega=g(w-\omega)<g(a_{j_1+1})<j_1+2$, that is $\bar g-(j_1+1)<1+\mu\omega$. The sections of $\mathrm{Rt}$
at heights $Y\le\bar g$ lie in $[\xi_I-Y\ta,\,w]$ and have length less than $\se+\bar g\ta=\omega$, and $\mathrm{Rt}$
lies below $Y=\bar g$. So $\operatorname{area}(\mathrm{Rt}\cap\{Y\ge j_1+1\})<\omega(1+\mu\omega)$. In total,
\[
\operatorname{area}(\mathrm{Lay})-\sum_jL_j<\lfloor\bar g\rfloor\bigl(1+\tfrac12\ta\bigr)+(1+\mu\omega)\,\omega .
\]

(i) \emph{The wedge part.} For $0\le r\le\lfloor g_0\rfloor-1$, the section of $\mathrm{Wd}$ at height $Y\in[r,r+1)$ lies
in $[0,\xi_0-Y\ta]$ (note $\xi_0-Y\ta=(m-Y)\ta>0$), so
$\operatorname{area}(\mathrm{Wd}\cap\{r\le Y<r+1\})\le\xi_0-(r+\frac12)\ta$, while $\Omega_r$ covers
$l_r>(m-r-1)\ta-1=\xi_0-(r+1)\ta-1$ (or nothing, when $(m-r-1)\ta<1$): the waste is less than $1+\frac12\ta$ per row,
in $\lfloor g_0\rfloor$ rows. The points of $\mathrm{Wd}$ with $Y\ge\lfloor g_0\rfloor$ have
$0\le X<\xi_0-Y\ta\le(m-\lfloor g_0\rfloor)\ta<(\delta+1)\ta$ and
$Y\le g(X)\le g_0+\mu(\delta+1)\ta<\lfloor g_0\rfloor+1+\mu(\delta+1)\ta$, so their area is less than
$(\delta+1)\ta\bigl(1+\mu(\delta+1)\ta\bigr)$.

Adding (g), (h) and (i) gives $U\le B^*$.
\end{proof}

The bound holds for every $y$ with $g_0>1$: no tuning of $y$, $m$ or $\alpha$ is needed, and by (a) every slab up to
the right wall gets a column.

\subsection{The band ends}\label{ssec:b23}

\begin{corollary}\label{cor:b23}
Let $b\ge10^4$ be an integer, $x:=b^{-1/3}$, $\kappa:=3/4$, $\tan(\theta/2)=b/(b^2-1)$, and $w,\gamma,\sigma$ as in
Section~\ref{sec:ends}. Let $y$ be any real number with
\[
\kappa b^{2/3}\le y\le\kappa b^{2/3}+2 .
\]
Apply the construction of Section~\ref{ssec:stairlemma} with $\sigma_0=\sigma$ and $\mu=\tan\theta$ (so that
$\Reg=\Reg(y)$ of Section~\ref{sec:ends}), and with any integer $Q\ge16b$. Then the construction gives a packing in
$\Reg(y)$ with
\[
U\ \le\ \Phi(x)\,b^{2/3}\ \le\ \Phi(10^{-4/3})\,b^{2/3}\ \le\ 5.03\,b^{2/3},
\]
where $\Phi$, defined in the proof, is non-decreasing on $(0,10^{-4/3}]$, $\Phi(10^{-4/3})\le5.0163$ and
$\Phi(x)\to\sqrt{A/\kappa}+3/2\approx4.3284$ as $x\to0$, with $A:=6.000002$. Consequently the end cost of
Proposition~\ref{prop:waste} satisfies $E(y)=U+\frac12\tan\theta\le5.03\,b^{2/3}+1.0001/b$.
\end{corollary}

\begin{proof}
(i) \emph{Band facts.} Write $t_\theta:=b/(b^2-1)$ and $u:=b^{1/3}=1/x$. Then $t_\theta\le1.00000002/b$ (since
$bt_\theta=b^2/(b^2-1)$), $\sigma<2t_\theta\le2.0000001/b$ and $\mu=2t_\theta/(1-t_\theta^2)\le2.000001/b$; so
$\Delta=w\mu<b\mu\le2.000001$ and $\sigma\mu\le4.1\cdot10^{-8}$. Also $w<b$ by Lemma~\ref{lem:width}, as $t_\theta>1/b$,
and $w\ge b\gamma=b\bigl(1-2t_\theta^2/(1+t_\theta^2)\bigr)\ge b-2bt_\theta^2\ge b-2.1/b$.

(ii) \emph{Ranges.} $g_0=y-\sigma\mu\in[\kappa u^2-4.1\cdot10^{-8},\ \kappa u^2+2]$. Put
$\nu:=\kappa u^2-5\cdot10^{-8}=p(x)u^2$ with $p(x):=\kappa-5\cdot10^{-8}x^2$; then $g_0\ge\nu>1$. Next
$\bar g=y+(w-\sigma)\mu\le y+\Delta\le\kappa u^2+4.000001$, $m=\lfloor\bar g\rfloor+1\le\kappa u^2+5.000001$,
$\lfloor g_0\rfloor\le\kappa u^2+2$, and $\delta=m-g_0\in(\Delta,\Delta+1]$, so $\delta\le\bar\delta:=3.000001$ and
$\delta+1\le4.000001$.

(iii) \emph{Tilt.} By Lemma~\ref{lem:tiltS} with $g_0\ge\nu$, $\delta\le\bar\delta$ and $A=2\bar\delta$,
\[
u\,t^*\le u\,\frac{1+\sqrt{1+2\nu\bar\delta}}{2\nu}=\frac{x+\sqrt{x^2+Ap(x)}}{2p(x)}\le\frac{x+\sqrt{x^2+A\kappa}}{2p(x)},
\]
and $u/Q\le u/(16u^3)=x^2/16$. So $t\le t^*+1/Q\le x\vartheta(x)=:\bar t(x)$ with
\[
\vartheta(x):=\frac{x+\sqrt{x^2+A\kappa}}{2p(x)}+\frac{x^2}{16}.
\]
Put $\bar\tau:=2\bar t/(1-\bar t^2)$ and $\bar\varsigma:=(1+\bar t^2)/(1-\bar t^2)$. Then $\ta\le\bar\tau$,
$\se\le\bar\varsigma$, $\sa\le2t\le2\bar t$ and $1-\ca\le2t^2\le2\bar t^2$. At $x_0:=10^{-4/3}$ (that is, $b=10^4$),
$\bar t\le0.0672$ and $\bar\tau\le0.1349$.

(iv) \emph{The hypotheses of Lemma~\ref{lem:stair}.} $g_0>1$ by (ii); $t\le\bar t(x)\le\bar t(x_0)\le0.0672<1$ by the
monotonicity in (v); and $m\ta+\se\le(0.75\,b^{2/3}+5.000001)\cdot0.1349+1.01<b-2.1/b\le w$, since $b^{2/3}\le b/21.5$
for $b\ge10^4$. Lemma~\ref{lem:stair} gives the packing and $U\le B^*$. We divide each term of $B^*$ by $u^2=b^{2/3}$,
using $w<b=u^3$, $\Delta\le2.000001$, $\mu\le2.000001\,x^3$ and $\omega=\se+\bar g\ta$:
\begin{align*}
T_1&=w\sa/u^2\le2ut\le2\vartheta(x),\qquad T_2=w(1-\ca)/u^2\le2ut^2\le2x\vartheta(x)^2,\\
T_3&=\Delta(\se+m\ta)/u^2\le2.000001\bigl[x^2\bar\varsigma+(\kappa+5.000001\,x^2)\bar\tau\bigr],\\
T_4&=\lfloor\bar g\rfloor\bigl(1+\tfrac12\ta\bigr)/u^2\le(\kappa+4.000001\,x^2)\bigl(1+\tfrac12\bar\tau\bigr),\\
T_5&=(1+\mu\omega)\,\omega/u^2\le(1+2.000001\,x\,\bar\omega)\,\bar\omega,\qquad
\bar\omega:=x^2\bar\varsigma+(\kappa+4.000001\,x^2)\bar\tau,\\
T_6&=\lfloor g_0\rfloor\bigl(1+\tfrac12\ta\bigr)/u^2\le(\kappa+2x^2)\bigl(1+\tfrac12\bar\tau\bigr),\\
T_7&=(\delta+1)\ta\bigl(1+\mu(\delta+1)\ta\bigr)/u^2\le4.000001\,\bar\tau x^2\bigl(1+2.000001\cdot4.000001\,x^3\bar\tau\bigr).
\end{align*}
For $T_5$ we used $\omega/u^2\le\bar\omega$ and $\mu\omega\le2.000001\,x^3u^2\bar\omega=2.000001\,x\bar\omega$. Let
$\Phi(x)$ be the sum of the seven right-hand sides.

(v) \emph{Monotonicity.} On $0<x\le x_0$, $p(x)$ is positive and decreasing and $x+\sqrt{x^2+A\kappa}$ is increasing,
so $\vartheta$ is increasing; $\bar t=x\vartheta$ is increasing and less than $1$, so $\bar\tau$ and $\bar\varsigma$ are
increasing. Every right-hand side in (iv) is a sum of products of non-negative non-decreasing functions of $x$. So
$\Phi$ is non-decreasing in $x$, that is non-increasing in $b$, and $\Phi(x)\le\Phi(x_0)$ for all $b\ge10^4$.

(vi) \emph{Values.} In interval arithmetic, $\Phi(x_0)\le5.01625$ (the terms are about $2.89126$, $0.19400$,
$0.20946$, $0.80975$, $0.10545$, $0.80515$ and $0.00116$). As $x\to0$, $T_1\to\sqrt{A\kappa}/\kappa=\sqrt{A/\kappa}$,
$T_4,T_6\to\kappa$ and the other terms tend to $0$, so $\Phi\to\sqrt{A/\kappa}+3/2$. Finally $A_0(y)$ exceeds the area
of $\Reg(y)$ by $\frac12\tan\theta$ (the triangle with vertices $(0,y-\sigma\mu)$, $(\sigma,y)$, $(0,y+\gamma)$), and
$\tan\theta\le2.000001/b$.
\end{proof}

The values in (iii) and (vi) were certified in interval arithmetic (\texttt{mpmath.iv}, $40$ digits, with $x_0$
enclosed in an interval) by the author's program and by three independent computations (two reviews and an
independent reimplementation), which all give $\Phi(x_0)\in[5.0162487,\,5.0162488]$; one of them also
bounds $\Phi$ on all of $(0,x_0]$ by a subdivision into $4000$ intervals, without using (v), and finds the maximum
$5.0163$ at $x_0$. For larger $b$ the constant improves: $\Phi(10^{-5/3})\le4.6351$, $\Phi(10^{-2})\le4.4685$,
$\Phi(10^{-8/3})\le4.3583$. The choice $\kappa=3/4$ is close to the minimizer $(\sqrt6/4)^{2/3}\approx0.721$ of the
limit $\sqrt{6/\kappa}+2\kappa$. The window has length $2>\sec\theta$, which is what Section~\ref{ssec:proofmain}
needs: the bound holds for every $y$ in the window, not only at sampled points.

\subsection{Proof of Theorem~\ref{thm:main}}\label{ssec:proofmain}
Let $C_E:=5.03$, $k\ge k_0=1.6\cdot10^7$, $b_*:=(3k/(2C_E))^{3/5}$ (so $k=\frac23C_Eb_*^{5/3}$) and $b:=\lceil b_*\rceil$.
Since $k_0\ge\frac23C_E\cdot10^{20/3}\approx1.5565\cdot10^7$, we have $b\ge b_*\ge10^4$; since $3/(2C_E)<1$,
$b\le k^{3/5}+1$. Take $\tan(\theta/2)=b/(b^2-1)>1/b$; then $w<b$ and $S:=k-b+w<k$. Use the packing of
Corollary~\ref{cor:assembly} with
\[
y_0:=\tfrac14\bigl\lceil3b^{2/3}\bigr\rceil\in\bigl[\kappa b^{2/3},\ \kappa b^{2/3}+\tfrac14\bigr]
\]
and, in each band, $R:=\lfloor(L_b-h-2y_0)/d\rfloor+1$. Here $h=b\sigma+\gamma<3$ (as $b\sigma=2b^2(b^2-1)/((b^2-1)^2+b^2)<2$) and
$y_0\le0.75(k^{3/5}+1)^{2/3}+\frac14\le0.75\,k^{2/5}+1$, so $h+2y_0<1.5\,k^{2/5}+5<k-b\le L_b$; hence $R\ge1$ and
$y_1=L_b-h-(R-1)d-y_0\in[y_0,\,y_0+d)$. As $d<1.0001$, both $y_0$ and $y_1$ lie in the window of
Corollary~\ref{cor:b23}, whatever value $y_1$ takes. Fill the four end regions $\Reg(y_0)$, $\Reg(y_1)$ of the two bands
by Corollary~\ref{cor:b23} with $Q$ the least power of $2$ that is at least $16b$. By Corollary~\ref{cor:assembly} this is
a packing in $[0,S]^2$, and by Proposition~\ref{prop:waste}
\[
k^2-N=(k^2-S^2)+\sum_{\text{2 bands}}\bigl[(R-1)\tan\theta+E(y_0)+E(y_1)\bigr].
\]
\begin{itemize}
\item $(R-1)\tan\theta=(R-1)d\sin\theta=(L_b-h-y_0-y_1)\sin\theta\le L_b\sigma<2L_bt_\theta$, so the two bands give less
than $(S+k-b)\,2t_\theta<4kb/(b^2-1)=4k/b+4k/(b(b^2-1))$.
\item $b-w=2t_\theta(bt_\theta-1)/(1+t_\theta^2)=2t_\theta/((b^2-1)(1+t_\theta^2))<2b/(b^2-1)^2$, so
$k^2-S^2<2k(b-w)<4kb/(b^2-1)^2$.
\item $\sum E\le4\bigl(C_Eb^{2/3}+\frac12\tan\theta\bigr)$ by Corollary~\ref{cor:b23}.
\end{itemize}
Hence $k^2-N<4k/b+4C_Eb^{2/3}+\lambda'$ with $\lambda':=4k/(b(b^2-1))+4kb/(b^2-1)^2+2\tan\theta\le8.001\,k/b^3+4.1/b$.
Since $k\le\frac23C_Eb^{5/3}$, $8.001\,k/b^3\le8.001\cdot\frac23C_E\,b^{-4/3}\le1.25\cdot10^{-4}$, and
$4.1/b\le4.1\cdot10^{-4}$; so $\lambda'<1$, and Lemma~\ref{lem:reduction} gives
$\cs(k)\le k^2-N-1<4k/b+4C_Eb^{2/3}$. Finally $4k/b\le4k/b_*=\frac83C_Eb_*^{2/3}$ and
$b^{2/3}\le(b_*+1)^{2/3}\le b_*^{2/3}+\frac23b_*^{-1/3}\le b_*^{2/3}+\frac23\cdot10^{-4/3}$, so
\[
\cs(k)<\frac{20}3C_Eb_*^{2/3}+\frac83C_E\cdot10^{-4/3}=\frac{20}3\Bigl(\frac32\Bigr)^{2/5}C_E^{3/5}k^{2/5}
+\frac83C_E\cdot10^{-4/3}.
\]
In interval arithmetic,
\[
\tfrac{20}3\bigl(\tfrac32\bigr)^{2/5}5.03^{3/5}\in[20.66740,\,20.66741]\quad\text{and}\quad
\tfrac83\cdot5.03\cdot10^{-4/3}\in[0.62259,\,0.62260],
\]
which gives Theorem~\ref{thm:main}. For every $C>\sqrt{A/\kappa}+3/2\approx4.3284$ we
have $\Phi(b^{-1/3})\le C$ for all large $b$, and the same argument with $C$ in place of $C_E$ gives
$\cs(k)\le(\frac{20}3(\frac32)^{2/5}C^{3/5}+o(1))k^{2/5}$; letting $C$ tend to the limit gives
$\cs(k)\le(18.89+o(1))k^{2/5}$. \qed

\begin{remark}[comparison and exact checks]\label{rem:stairexact}
The construction of Section~\ref{ssec:stairlemma} was implemented in exact rational arithmetic, with $Q$ the least power
of $2$ that is at least $16b$. For $b=400$, $1600$, $6400$, $25600$, $102400$ and five lifts each in the window of
Corollary~\ref{cor:b23}, all $25$ end packings were checked by an exact checker that does not use
Lemma~\ref{lem:stair} (shapes, containment, separating axis tests for all pairs of pieces with overlapping bounding
boxes; for $b=102400$ also by a check of the separating lines of the proof); all pass, $U\le B^*$ holds in each, and
$U/b^{2/3}$ stays between $3.26$ and $4.36$. Deliberately corrupted versions (a column shifted by $2^{-30}$, a layer
lengthened, a column given one more square, and so on) were rejected. The full packings of
Section~\ref{ssec:proofmain}, counted exactly (the four ends checked by exact containment and by the separating lines
of the proof; the rows not checked, their correctness rests on Lemmas~\ref{lem:stack} and~\ref{lem:ends}; the identity
of Proposition~\ref{prop:waste} checked exactly), give:
\begin{center}
\small
\begin{tabular}{rrrrr}
\toprule
$k$ & $b$ & $y_0$ & $\cs(k)\le$ & $20.668\,k^{2/5}+0.623$\\
\midrule
$10^5$ & $484$ & $185/4$ & $1794$ & ($k<k_0$)\\
$10^6$ & $1927$ & $465/4$ & $4467$ & ($k<k_0$)\\
$10^7$ & $7669$ & $1167/4$ & $10579$ & ($k<k_0$)\\
$1.6\cdot10^7$ & $10167$ & $352$ & $12707$ & $15738.5$\\
$3.825\cdot10^7$ & $17152$ & $499$ & $17865$ & $22303.0$\\
$10^8$ & $30530$ & $2931/4$ & $26998$ & $32757.2$\\
\bottomrule
\end{tabular}
\end{center}
For $k<k_0$ these packings are still valid (Lemma~\ref{lem:stair} does not need $b\ge10^4$); only the closed form of
Theorem~\ref{thm:main} is restricted to $k\ge k_0$. An independent reimplementation, written from the text of the proof
only (with its own admissible choices of $Q$), checked the end packings exactly for $b$ up to $4\cdot10^5$ and many
lifts in the window, and the full packings for $k=1.6\cdot10^7$, $3.825\cdot10^7$, $10^8$, $10^9$ and $3\cdot10^9$: no
overlaps, no containment errors, $U\le B^*$ in every case, and $\cs(10^9)\le67065$, $\cs(3\cdot10^9)\le99640$ (against
$82281.4$ and $127687.7$ from Theorem~\ref{thm:main}); $8$ of $8$ deliberately corrupted versions were rejected. It
relies on Lemmas~\ref{lem:stack} and~\ref{lem:ends} for the rows. The tuned packings of Theorem~\ref{thm:cert} are
$4.9$--$11.8\%$ better than these (Section~\ref{sec:cert}).
\end{remark}

\section{Chained wall fillers and the exponent \texorpdfstring{$3/8$}{3/8}}\label{sec:38}

In the end lemma of Section~\ref{sec:stair} the columns cost about $b/\sqrt m$ and the two walls about $m$ each. In this
section the walls are filled more efficiently: each wall is a thin trapezoid, and it is filled by a chain of
staircases (\emph{blocks}) whose tilt is lowered slightly from block to block, so that the horizontal layers opened in
one block close just before the next block. The wall cost becomes $O(m^{3/4})$, the best column length is
$m\asymp b^{4/5}$, the end cost is $O(b^{3/5})$, and Proposition~\ref{prop:waste} gives Theorem~\ref{thm:main38}. The
ideas (horizontal layers under tilted columns, and filling the wall trapezoids by stairs whose tilt is lowered from
block to block) are those of Bui [Bui, Sections~3--4]; the construction and the proof below are self-contained and do not
use any argument of [Bui].

The construction has six real parameters: the threshold $b_0$, the window coefficient $c_y$ of the lifts, the wall
widths $c_D$, $c_R$, the wall tilt factor $a_w$ and the proof parameter $\varpi$ of Lemma~\ref{lem:W}. Version~1.0 of
this paper used the single choice $(10^{16},1/12,4,4,3/2,0.65)$. Version~1.1 keeps all lemmas of version~1.0 and needs
no new lemma for the parameters; it states them with the parameters, adds two short arguments that version~1.0 did not
need (the error of the wall tilts in (E6) for all tilts up to $1/10$, and the room condition in
Section~\ref{ssec:proofthm38}), and certifies the constants for each row of
Table~\ref{tab:certtiers} (Section~\ref{ssec:b35}); this gives the tiers of Theorem~\ref{thm:main38}.
Section~\ref{ssec:refine} adds a variant ZC$'$ of the wall filler and two lemmas (Lemmas~\ref{lem:P2}
and~\ref{lem:P3}) with a smaller bound; only the tiers marked $\dagger$ use them.
Sections~\ref{ssec:sawtooth} and~\ref{ssec:liftT2} add a sawtooth accounting of two groups of wall rows
(Lemmas~\ref{lem:SL}--\ref{lem:W4}) and a choice of the lift of each band (Lemma~\ref{lem:T2}); only the tiers marked
$\ddagger$ use them.

Throughout this section we use the tilt notation of Section~\ref{ssec:stairlemma}, the columns $K(\xi,J,n)$ and
Lemmas~\ref{lem:col} and~\ref{lem:tiltS}, and we put
\[
q:=1-\ca=\frac{\ta^2}{\se(\se+1)}\le\frac{\ta^2}2,\qquad\text{so that}\qquad \se-1=q\se,\quad \se\sa=\ta .
\]
$q$ is increasing in $t$. On $t\in(0,1/10]$ we have $\ta\le2.1t$, and $t\mapsto\ta$ has derivative
$2(1+t^2)/(1-t^2)^2\le2.1$, so neighbouring points of a grid $(1/Q^*)\Z$ differ in $\ta$ by at most
$\delta_Q:=2.1/Q^*$. A tilt with index $0$ (as $t_0$, $\alpha_0$) is the \emph{initial tilt} of a wall. Two closed convex
sets that lie in the two closed half-planes of one line have disjoint interiors; every disjointness statement below is
of this kind.

\emph{Grids.} $Q$ is the grid of the main tilt: any power of $2$ with $Q\ge16b$. The wall tilts and the cut tilt use the
grid $Q^*:=\max(2^{40},Q)$, and the phases the grid $Q_f:=\max(2^{48},Q)$.

\subsection{The wall filler}\label{ssec:wall}
For reals $\chi\ge4$, $\zeta>0$ and $L\ge0$ let $r(x):=\chi+\zeta x$ and
\[
Z=Z(\chi,\zeta,L):=\{(x,y):\ 0\le x\le L,\ 0\le y\le r(x)\},
\]
a right trapezoid of area $L\chi+\zeta L^2/2$. The input is a grid tilt $t_0\in(0,1/10]\cap(1/Q^*)\Z$; put
$N_0:=\lfloor\chi\rfloor-1\ge3$. For a block $\beta$ with tilt $t_\beta$ we write $\lambda_\beta(x,y):=x+y\tan\alpha_\beta$.

\emph{Construction ZC.}
\begin{enumerate}[label=(Z\arabic*)]
\item \emph{Block $0$:} tilt $t_0$, base $F_0:=\lfloor\chi-N_0\caz-\saz\rfloor\ge0$, total length $M_0:=F_0+N_0$, offset
$s_0:=M_0\taz$.
\item \emph{Block $\beta$} has tilt $t_\beta$ (with $\ca,\sa,\ta,\se,q$), offset $s_\beta$, base $F_\beta$ and total
length $M_\beta=F_\beta+N_0$. If $s_\beta+\se>L$ the chain stops (block $\beta$ does not exist). Otherwise, for
$i=0,1,2,\dots$ let $\xi_i:=s_\beta+i\se$; stop (reason `$L$') if $\xi_i+\se>L$; else let
$\Gamma_i:=r(\xi_i-M_\beta\ta)$ and $\tilde J_i:=\lfloor(\Gamma_i-M_\beta\ca-\sa)/q\rfloor$; stop (reason
`saturation') if $\tilde J_i>M_\beta-1$; else place the column $K(\xi_i,J_i,n_i)$ with $J_i:=\tilde J_i$,
$n_i:=M_\beta-J_i$. Let $P_\beta$ be the number of columns and $\hat\xi_\beta:=s_\beta+P_\beta\se$.
\item \emph{Next block.} If block $\beta$ stopped by saturation (it is \emph{complete}), let $\kappa_\beta:=q^2/\zeta$,
let $t_{\beta+1}$ be the largest grid point $t'\le t_\beta$ with $\tan\alpha'\le\ta-\kappa_\beta$, and put
$\hat\kappa_\beta:=\ta-\tan\alpha_{\beta+1}\ge\kappa_\beta$ and $\eta_\beta:=\hat\kappa_\beta-\kappa_\beta\in[0,\delta_Q)$.
\emph{Phase:} let $J_{\rm lo}:=J_0$ and $J_{\rm L}:=J_{P_\beta-1}$ (indices of block $\beta$). If $J_{\rm L}>J_{\rm lo}$, put
$A:=(\Gamma_0-M_\beta\ca-\sa)/q$, $B:=q/(\zeta\se)$, $p:=\#\{i:\ J_i\le J_{\rm lo}\}$, $f:=p-B(J_{\rm lo}+1-A)$,
\[
\psi:=\hat\xi_\beta-\tan\alpha_{\beta+1}-(s_\beta+p\se)+J_{\rm lo}\hat\kappa_\beta,\qquad
\omega_0:=\operatorname{frac}\bigl(-(\psi+f(\se-1)-(\se-1))\bigr),
\]
and $\omega:=\lceil\omega_0Q_f\rceil/Q_f\in[0,1]$; otherwise $\omega:=0$. Then $s_{\beta+1}:=\hat\xi_\beta+\omega$ and,
with primes for the tilt $t_{\beta+1}$,
\begin{align*}
F_{\beta+1}&:=\max\{J\in\Z:\ J+N_0\ca'+\sa'\le r(s_{\beta+1}-(J+N_0)\ta')\}\\
&=\Bigl\lfloor\frac{\chi+\zeta s_{\beta+1}-N_0(\ca'+\zeta\ta')-\sa'}{1+\zeta\ta'}\Bigr\rfloor,
\end{align*}
$M_{\beta+1}:=F_{\beta+1}+N_0$.
\item \emph{Rows.} Block $\beta$ \emph{blocks} row $j$ if $J_{\beta,0}\le j$ ($J_{\beta,i}$ is $J_i$ of block $\beta$).
For each integer $0\le j<\lfloor r(L)\rfloor$, walk from $x=0$ to $x=L$: a \emph{stretch} starts at $x=0$ or just after
a blocking block $\beta$, at the abscissa $\xi_{\beta,p}-j\tan\alpha_\beta$ with $p:=\#\{i:\ J_{\beta,i}\le j\}$
($p=P_\beta$ allowed, $\xi_{\beta,P_\beta}=\hat\xi_\beta$), and ends at the next blocking block $\beta_2$, at the
abscissa $s_{\beta_2}-(j+1)\tan\alpha_{\beta_2}$, or at $x=L$. For a stretch from $a$ to $a^+$ with $a^+-a\ge1$, $a\ge0$
and $j+1\le r(a)$, place the row $[a,\,a+\lfloor a^+-a\rfloor]\times[j,j+1]$.
\end{enumerate}
Each block is one staircase of Section~\ref{ssec:stairlemma} (one tilt, total length $M_\beta$, $J$ non-decreasing); a
block ends when its staircase is used up, and the next block starts after a phase gap $\omega<1$ with the slightly
smaller tilt, chosen so that every layer opened in the block ends just before the next block with a small residual
(see (R) below).

\emph{Parameters.} Let $\varpi\in(0,1)$ (a free proof parameter; the tiers of Section~\ref{ssec:b35} use $\varpi=0.65$
or $\varpi=0.681818$), let $q_{\min}$ be the value of $q$ at the slope
$\varpi\taz$, that is $q_{\min}:=(\varpi\taz)^2/(s_m(s_m+1))$ with $s_m:=\sqrt{1+\varpi^2\taz^2}$, and put
$M_{\max}:=r(L)+N_0q_0$ ($q_0$ is $q$ at $t_0$),
\[
\Lambda:=\frac{q_{\min}N_0-1-\zeta\taz}{\zeta},\qquad \Kb:=\frac L\Lambda,\qquad
\varepsilon_W:=(\sez-1)+N_0\delta_Q+\frac1{Q_f}.
\]
\emph{Hypotheses.}
\begin{enumerate}[label=(W\arabic*)]
\item $\Lambda\ge N_0\taz$ (in particular $\Lambda>0$);
\item $\Delta_{\rm dr}:=\dfrac{q_0}{N_0}\Bigl(L+\Kb\,\dfrac{1+\zeta\taz}{\zeta}\Bigr)+\Kb\delta_Q<(1-\varpi)\taz$;
\item $N_0\caz\ge\dfrac{\zeta\sez}{q_{\min}}+\taz+\zeta(2+\sez+M_{\max}\taz)+\dfrac{1+\zeta\taz}{q_{\min}}$;
\item $\varepsilon_W<1$;
\item $\delta_Q\le\varpi\taz$ and $t_0\le1/10$.
\end{enumerate}

\begin{lemma}[wall filler]\label{lem:W}
Under (W1)--(W5) the construction ZC does not fail, its pieces form a packing in $Z$, and the uncovered area
$U_Z:=\operatorname{area}(Z)-\#\text{squares}$ satisfies
\begin{multline*}
U_Z\le B_W:=Lc_{\rm col}+\Kb\Bigl[(N_0-1)(\varepsilon_W+\taz)+\frac{1+\zeta\taz}{q_{\min}}(1+\taz)+A_T\Bigr]\\
+\lfloor\chi\rfloor\bigl(1+\tfrac12\taz\bigr)+w_0(1+\zeta w_0)+r(L)\bigl(1+\tfrac12\taz\bigr)+w_e(4+2\zeta w_e),
\end{multline*}
where $c_{\rm col}:=\taz+\sez q_0+\sez\zeta(\sez+N_0q_0\taz)$, $A_T:=h_Tw_T$,
\[
h_T:=1+\frac{\zeta\sez}{q_{\min}}+\taz+\zeta(1+\sez+M_{\max}\taz),\qquad
w_T:=1+(r(L)+1)\Bigl(\frac{q_0^2}\zeta+\delta_Q\Bigr),
\]
$w_0:=(1+N_0q_0)\taz$ and $w_e:=1+\sez+(r(L)+N_0q_0)\taz$.
\end{lemma}

The terms are: column triangles, gaps and ceiling slope; aligned closings; unaligned closings; transition slivers; the
left end; the right end.

\begin{proof}
(a) \emph{One block.} Fix an existing block $\beta$ with tilt $t$ ($\ca,\sa,\ta,\se,q$), offset $s$, base $F$ and total
$M=F+N_0$, and assume the invariants
\begin{enumerate}[label=(I\arabic*)]
\item $\ta\ge\varpi\taz$ (hence $q\ge q_{\min}$);
\item $s\ge M\ta$;
\item $F\ge0$ and $F+N_0\ca+\sa\le r(s-M\ta)<F+1+N_0\ca+\sa+\zeta\ta$.
\end{enumerate}
(I3) holds for $\beta=0$ ($\Gamma_0=r(0)=\chi$ and the definition of $F_0$) and for $\beta+1$ by the definition of
$F_{\beta+1}$ as a maximum ($J=F_{\beta+1}+1$ violates the defining inequality, and
$r(s'-(F'+1+N_0)\ta')=r(s'-M'\ta')-\zeta\ta'$).
$\Gamma_i=r(s-M\ta)+i\zeta\se$ is affine in $i$, so $\tilde J_i=\lfloor A+i/B\rfloor$, and the $J_i$ are non-decreasing.
$H(J):=J+(M-J)\ca+\sa$ is affine in $J$ with slope $q$, so $J_i=\max\{J:\ H(J)\le\Gamma_i\}$ and the gap
$\Gamma_i-H(J_i)$ lies in $[0,q)$. $J_0\ge F$ because $H(F)=F+N_0\ca+\sa\le\Gamma_0$ by (I3); so $J_i\ge F$,
$n_i\le N_0$, and $n_i\ge1$ (no saturation). By (I3), $A<F+(1+\zeta\ta)/q$, so $J_0<F+(1+\zeta\ta)/q$.
The block is not empty: it exists, so $s+\se\le L$, and $J_0<F+(1+\zeta\taz)/q_{\min}<F+N_0$ by (W1); so
$\tilde J_0\le M-1$ and $P\ge1$.
\emph{Containment of columns} (Lemma~\ref{lem:col}, $J_i+n_i=M$): $X\ge\xi_i-M\ta\ge s-M\ta\ge0$ by (I2),
$X\le\xi_i+\se\le L$, and $J_i\le Y\le H(J_i)\le\Gamma_i=r(\xi_i-M\ta)\le r(X)$.
\emph{Length of a complete block.} Saturation at $i=P$ means $A+P/B\ge M$, so $P>B(N_0-(1+\zeta\ta)/q)$ and
\begin{equation}\label{eq:L1}
\Lambda_\beta:=P\se>\frac{qN_0-1-\zeta\ta}{\zeta}\ge\Lambda,\qquad\text{equivalently}\qquad
q\le\frac{\zeta\Lambda_\beta+1+\zeta\ta}{N_0}
\end{equation}
(the middle expression is increasing in $q$ and decreasing in $\ta$; use $q\ge q_{\min}$, $\ta\le\taz$).
\emph{Last column of a complete block:} $J_{\rm L}:=J_{P-1}\ge\tilde J_P-1/B-1\ge M-1/B-1$, so
$n_{\rm L}:=M-J_{\rm L}\le1+1/B=1+\zeta\se/q$.

(b) \emph{Passing to block $\beta+1$} ($\beta$ complete; primes refer to $t_{\beta+1}$).
\emph{Tilt drift.} Let $C_\beta$ be the sum of $\kappa_\gamma+\delta_Q$ over the complete blocks $\gamma\le\beta$, and
suppose (I1) for the blocks $0,\dots,\beta$. By \eqref{eq:L1} these blocks have total length at least
$(\beta+1)\Lambda$ inside $[0,L]$, so $\beta+1\le\Kb$; and by \eqref{eq:L1},
$\kappa_\gamma=q_\gamma^2/\zeta\le(q_0/N_0)\bigl(\Lambda_\gamma+(1+\zeta\taz)/\zeta\bigr)$. Summing ($\sum\Lambda_\gamma\le L$)
gives $C_\beta\le\Delta_{\rm dr}<(1-\varpi)\taz$ by (W2). By induction $\tan\alpha_\beta\ge\taz-C_{\beta-1}$, so
$\ta-\kappa_\beta\ge\taz-C_\beta+\delta_Q>\varpi\taz+\delta_Q\ge\tan\alpha(1/Q^*)$ by (W5): a positive grid point $t'$
exists, and $\ta'>\ta-\kappa_\beta-\delta_Q\ge\taz-C_\beta>\varpi\taz$. So (I1) holds for $\beta+1$, and $t'\le t$,
$\ca'\ge\ca$.
\emph{$F'\ge F$.} $J=F$ is admissible in the definition of $F'$:
$r(s'-M\ta')\ge r(\hat\xi-M\ta)=\Gamma_P$, saturation gives $\Gamma_P\ge M\ca+\sa+qM=M+\sa$, and
$F+N_0\ca'+\sa'\le M+\sa'\le M+\sa$.
\emph{$J'_0\le J_{\rm L}$.} By maximality of $J_{\rm L}$,
$\Gamma_{P-1}<H(J_{\rm L}+1)\le J_{\rm L}+n_{\rm L}+\sa$. As $s'\le\hat\xi+1=\xi_{P-1}+\se+1$,
$r(s')\le\Gamma_{P-1}+\zeta(1+\se+M\ta)$. By (I3) for $\beta+1$,
$F'+N_0\ca'+\sa'\le r(s'-M'\ta')\le r(s')$, so $F'<J_{\rm L}+n_{\rm L}+\sa-N_0\ca'+\zeta(1+\se+M\ta)$; and by (a) for
block $\beta+1$, $J'_0<F'+(1+\zeta\ta')/q'$. With $n_{\rm L}\le1+\zeta\se/q$, $M\le M_{\max}$ (from $F+N_0\ca+\sa\le r(L)$),
$\ca'\ge\caz$, $q,q'\ge q_{\min}$, $\se\le\sez$, $\sa\le\taz$, hypothesis (W3) gives $J'_0<J_{\rm L}+1$.
\emph{(I2) for $\beta+1$:} $M'=F'+N_0\le J'_0+N_0\le J_{\rm L}+N_0\le M-1+N_0$, so
$M'\ta'\le M\ta+N_0\taz\le s+\Lambda_\beta=\hat\xi\le s'$ by (I2), \eqref{eq:L1} and (W1).
\emph{Aligned closings.} For $J_{\rm lo}\le j<J_{\rm L}$ put $p(j):=\#\{i:\ J_i\le j\}$; then $1\le p(j)\le P-1$ and,
since $j\ge\lfloor A\rfloor=J_{\rm lo}$, $p(j)=\lceil B(j+1-A)\rceil=B(j+1-A)+f_j$ with $f_j\in[0,1)$. The stretch of row
$j$ from $a_j=\xi_{p(j)}-j\ta$ to $s'-(j+1)\ta'$ has length $\ell_j=s'-\ta'-s-p(j)\se+j\hat\kappa$. Modulo $1$ (as
$p(j)$ is an integer and $\se=1+(\se-1)$, and $B(\se-1)=q(\se-1)/(\zeta\se)=q^2/\zeta=\kappa$),
\[
\ell_j\equiv\omega+C+j\eta-f_j(\se-1),\qquad C:=\hat\xi-\ta'-s-B(1-A)(\se-1).
\]
By (Z3), $\psi+f(\se-1)\equiv C+J_{\rm lo}\eta$, so $\omega_0+C+J_{\rm lo}\eta-(\se-1)$ is an integer and
\[
\ell_j\equiv(\omega-\omega_0)+(j-J_{\rm lo})\eta+(1-f_j)(\se-1)\pmod1,
\]
where $0\le\omega-\omega_0<1/Q_f$, $0\le(j-J_{\rm lo})\eta<N_0\delta_Q$ (as $J_{\rm L}-J_{\rm lo}\le N_0-1$) and
$0<(1-f_j)(\se-1)\le\sez-1$. The right-hand side lies in $(0,\varepsilon_W]$ with $\varepsilon_W<1$ (W4). Also
$\ell_j\ge\hat\xi-\xi_{p(j)}-\ta'\ge\se-\ta'>0$. Hence
\[
\ell_j\ge0\quad\text{and}\quad\operatorname{frac}(\ell_j)\le\varepsilon_W .\tag{R}
\]

(c) \emph{Global structure.} By (a), (b) and induction every block satisfies (I1)--(I3), ZC never fails, the bases
$F_\beta$ are non-decreasing, the tilts non-increasing, $s_{\beta+1}\ge\hat\xi_\beta$, $J_{\beta+1,0}\le J_{\beta,\rm L}$, and
there are at most $\Kb$ complete blocks. For $Y\ge0$ and $\beta<\beta''$ we have $\lambda_{\beta''}\le\lambda_\beta$.
A stretch of row $j$ runs from just after a blocking block $\beta$ (left line $\lambda_\beta=\xi_{\beta,p(j)}$) to the
next blocking block $\beta_2>\beta$ (line $\lambda_{\beta_2}=s_{\beta_2}$), passing below all blocks in between, whose
columns have $J\ge J_{\beta'',0}\ge j+1$. If $\beta$ is not the last block and $\beta+1$ does not block row $j$, then
$j<J_{\beta+1,0}\le J_{\beta,\rm L}$, so $p(j)\le P_\beta-1$. Hence $p(j)=P_\beta$ (left end on
$\lambda_\beta=\hat\xi_\beta$) occurs only for stretches that end at the next block's line (\emph{transition rows},
$j\ge J_{\beta,\rm L}$) or for the last block (\emph{end rows}).

(d) \emph{Containment of rows.} A placed row has $a\ge0$, $a+\lfloor a^+-a\rfloor\le a^+\le L$ (as $a^+$ is $L$ or
$s_\beta-(j+1)\tan\alpha_\beta\le s_\beta\le L$), and $j+1\le r(a)\le r(X)$ on it. So every piece lies in $Z$.

(e) \emph{Disjointness.} Columns of one block lie in the slabs $\xi_i\le\lambda_\beta\le\xi_{i+1}$
(Lemma~\ref{lem:col}(i)). For $\beta<\beta''$, a column of $\beta$ lies in
$\{\lambda_\beta\le\hat\xi_\beta\}\subset\{\lambda_{\beta''}\le\hat\xi_\beta\}\subset\{\lambda_{\beta''}\le s_{\beta''}\}$
(as $Y\ge0$), and a column of $\beta''$ in $\{\lambda_{\beta''}\ge s_{\beta''}\}$. A row piece of a stretch of row $j$
from block $\beta$ to block $\beta_2$ (or $x=L$): all columns of blocks $\gamma<\beta$ and the columns $i<p(j)$ of $\beta$
lie in $\{\lambda_\beta\le\xi_{\beta,p(j)}\}$, and the row piece in $\{\lambda_\beta\ge\xi_{\beta,p(j)}\}$ (its lower left
corner is on that line and $\tan\alpha_\beta>0$); the columns $i\ge p(j)$ of $\beta$ and all columns of the blocks
strictly between $\beta$ and $\beta_2$ have $J\ge j+1$, so they lie in $\{Y\ge j+1\}$ and the row piece in
$\{Y\le j+1\}$; the columns of $\beta_2$ and of later blocks lie in $\{\lambda_{\beta_2}\ge s_{\beta_2}\}$, and the row
piece in $\{\lambda_{\beta_2}\le s_{\beta_2}\}$ (its top right corner is on or left of that line). For the first stretch
of a row (left end $x=0$) only the last two cases occur. Different rows are separated by the lines $Y=j+1$, and the
stretches of one row by the lines above. So ZC is a packing in $Z$.

(f) \emph{Accounting.} For each column $K_{\beta,i}$ let its \emph{cell} be
$\{z\in Z:\ \xi_i\le\lambda_\beta<\xi_{i+1},\ Y\ge J_i\}$. The cells are pairwise disjoint: slabs of one block are
disjoint, and with $E_\beta:=\{z\in Z:\ \lambda_\beta(z)<s_\beta\}$ we have $E_0\subset E_1\subset\cdots$, and the cells
of block $\beta$ lie in $E_{\beta+1}\setminus E_\beta$. For $Y\in[j,j+1)$ and a slab of a blocking block, $Y\ge J_i$ iff
$J_i\le j$; so the complement of the cells in $Z\cap\{j\le Y<j+1\}$ is the disjoint union of the \emph{stretch regions}
$\{j\le Y<j+1,\ \lambda_\beta\ge\xi_{\beta,p(j)}\ (\text{or }X\ge0),\ \lambda_{\beta_2}<s_{\beta_2}\ (\text{or }X\le L)\}\cap Z$.
The interior of every placed piece lies in its cell or stretch region, and boundaries have measure $0$; so
$U_Z$ is the sum over the cells of (area $-\,n$) plus the sum over the stretch regions of (area $-$ placed squares).

(f1) \emph{Cells.} The parallelogram $\{\xi_i\le\lambda_\beta\le\xi_i+\se,\ J_i\le Y\le H(J_i)\}$ has area
$\se(n\ca+\sa)=n+\ta$, so the part of the cell in $\{Y\le H(J_i)\}$ has area at most $n+\ta$. The rest of the cell lies in
$\{Y>H(J_i)\}$, where $X\le\xi_i+\se-H(J_i)\ta$, so $Y\le r(\xi_i+\se-H(J_i)\ta)=\Gamma_i+\zeta(\se+(M-H(J_i))\ta)$ with
$M-H(J_i)=nq-\sa\le N_0q$; its sections have length at most $\se$. So (area of the cell) $-\,n\le
\ta+\se q+\se\zeta(\se+N_0q\ta)\le c_{\rm col}$ (all terms increase with $t\le t_0$, and the gap is less than $q$). The
slab of column $(\beta,i)$ meets the line $Y=0$ in $[\xi_i,\xi_i+\se]\subset[0,L]$; these intervals have disjoint
interiors (consecutive inside a block, and $s_{\beta+1}\ge\hat\xi_\beta$) and length $\se\ge1$, so there are at most $L$
columns, and the cells cost at most $Lc_{\rm col}$.

(f2) \emph{Stretch areas.} A stretch region between the lines $\lambda_\beta=c$ and $\lambda_{\beta_2}=s'$ has horizontal
sections of length $(s'-Y\tan\alpha_{\beta_2})-(c-Y\tan\alpha_\beta)$, affine in $Y$ with slope
$\tan\alpha_\beta-\tan\alpha_{\beta_2}\ge0$; so its area is at most $\ell+\taz$, where
$\ell:=s'-(j+1)\tan\alpha_{\beta_2}-(c-j\tan\alpha_\beta)$ is the stretch length of (Z4) (with $\frac12\taz$ for an end at
$x=0$ or $x=L$). A placed piece leaves at most $\operatorname{frac}(\ell)+\taz$; if $0\le\ell<1$ the waste is at most
$\ell+\taz$ as well.

(f3) \emph{The stretches}, each counted by its right end.
\begin{itemize}
\item Right end at $s_{\beta+1}$, left end in block $\beta$ with $J_{\beta,0}\le j<J_{\beta,\rm L}$ (\emph{aligned
closings}): by (R), $\ell\ge0$ and $\operatorname{frac}(\ell)\le\varepsilon_W$; the left end has
$a_j\ge\xi_p-M\ta\ge0$ and $j+1\le J_{p(j)}\le H(J_{p(j)})\le\Gamma_{p(j)}\le r(a_j)$, so the piece is placed when
$\ell\ge1$. Waste at most $\varepsilon_W+\taz$, in at most $N_0-1$ rows per complete block.
\item Right end at $s_{\beta+1}$, left end before block $\beta$ ($j<J_{\beta,0}$; \emph{unaligned closings}): these rows
have $J_{\beta+1,0}\le j<J_{\beta,0}$, at most $J_{\beta,0}-F_{\beta}<(1+\zeta\taz)/q_{\min}$ per complete block; the
left end is $x=0$ or a slab line with $p\le P-1$, and the waste is at most $1+\taz$.
\item Right end at $s_{\beta+1}$, left end on $\hat\xi_\beta$ ($j\ge J_{\beta,\rm L}$; \emph{transition rows}): their
union lies in $\{\lambda_\beta\ge\hat\xi_\beta,\ \lambda_{\beta+1}<s_{\beta+1},\ Y\ge J_{\beta,\rm L}\}$. Its points have
$X<s_{\beta+1}$, so $Y\le r(s_{\beta+1})<J_{\rm L}+n_{\rm L}+\sa+\zeta(1+\se+M\ta)$ (see (b)): height at most $h_T$.
At height $Y$ its section has length $\omega+Y\hat\kappa\le w_T$ (as $\hat\kappa\le q_0^2/\zeta+\delta_Q$). Area at most
$A_T$ per complete block.
\item Right end at $s_0$, left end $x=0$: for $j\le\lfloor\chi\rfloor-1$ the piece is placed ($j+1\le\chi=r(0)$), waste
at most $1+\frac12\taz$, in at most $\lfloor\chi\rfloor$ rows. The part of $E_0$ with $Y\ge\lfloor\chi\rfloor$ has
$X<s_0-\lfloor\chi\rfloor\taz\le(M_0-\lfloor\chi\rfloor)\taz<w_0$ and $Y\le r(w_0)\le\lfloor\chi\rfloor+1+\zeta w_0$:
area at most $w_0(1+\zeta w_0)$.
\item Right end at $x=L$ (\emph{final stretches}), at most one per row $j<\lfloor r(L)\rfloor$. If the left end is $x=0$ or
a slab line with $p\le P-1$, the piece is placed when $\ell\ge1$ and the waste is at most $1+\frac12\taz$. Otherwise, by
(c), the left end is on $\hat\xi$ of the last block, and these end rows lie in $E_{\rm end}:=\{\lambda_{\rm last}\ge
\hat\xi_{\rm last}\}$, where $L-\hat\xi_{\rm last}<1+\sez$ (the last block stopped with reason `$L$', or the next block
could not start). An end row with $a_j=\hat\xi-j\ta\ge0$ and $j+1\le r(a_j)$ is placed with waste at most
$1+\frac12\taz$ (counted in $r(L)(1+\frac12\taz)$). The other points of $E_{\rm end}$: (i) rows with $a_j\ge0$ and
$j+1>r(a_j)$; as $a_j\ge L-w_e$, these rows have $j>r(L)-\zeta w_e-1$, and their points have $X\in[L-w_e,L]$,
$Y\le r(L)$: area at most $w_e(1+\zeta w_e)$; (ii) rows with $a_j<0$: then $j>\hat\xi/\ta\ge(s+\se)/\ta\ge M+1/\ta>N_0$
by (I2), so there are at most $r(L)-\lfloor\chi\rfloor+1\le\zeta L+2$ such rows, and $\hat\xi<j\ta\le r(L)\taz$ gives
$L<w_e$, so each such row meets $Z$ in area at most $L<w_e$: area at most $(\zeta w_e+2)w_e$. In total at most
$w_e(4+2\zeta w_e)$.
\end{itemize}
(f4) Summing (the cells; per complete block, of which there are at most $\Kb$, the aligned, unaligned and transition
rows; the start; the final stretches) gives $U_Z\le B_W$. If there is no block at all ($s_0+\sez>L$), all rows run from
$0$ to $L$, rows $j\le\lfloor\chi\rfloor-1$ are placed with waste at most $1$, and the rest has area at most
$L(1+\zeta L)\le w_e(1+\zeta w_e)$, as $L<s_0+\sez\le w_e$.
\end{proof}

\subsection{The variant ZC\texorpdfstring{$'$}{'}, transition rows and the end region}\label{ssec:refine}
This subsection is new in version~1.1. Only the tiers marked $\dagger$ in Tables~\ref{tab:tiers}
and~\ref{tab:certtiers} use it. It changes one rule of ZC and replaces two crude terms of $B_W$, the transition term
$A_T$ and the end term $w_e(4+2\zeta w_e)$.

\emph{Construction ZC$'$.} It is identical to ZC except for the last sentence of (Z4), which becomes:
\begin{enumerate}[label=(Z4$'$)]
\item For a stretch of row $j$ from $a$ to $a^+$ put
\[
a^*:=\max\bigl(a,\ 0,\ (j+1-\chi)/\zeta\bigr)
\]
(the least $x\ge\max(a,0)$ with $r(x)\ge j+1$); if $a^+-a^*\ge1$, place the row $[a^*,\,a^*+\lfloor a^+-a^*\rfloor]\times[j,j+1]$.
\end{enumerate}
(Z1)--(Z3) (blocks, columns, tilts, phases) and the stretches (their left and right ends) are unchanged.

\emph{Observation (O).} If $a\ge0$ and $j+1\le r(a)$, then $a^*=a$, and ZC$'$ places exactly the piece that ZC places
(if $a^+-a<1$, neither places a piece). So ZC$'$ differs from ZC only on stretches where ZC places nothing because $a<0$
or $j+1>r(a)$; there ZC$'$ may place a row that starts further to the right.

\begin{lemma}\label{lem:P3a}
Under (W1)--(W5) the construction ZC$'$ does not fail, its pieces form a packing in $Z$, and $U_Z\le B_W$.
\end{lemma}

\begin{proof}
(Z1)--(Z3) are those of ZC, so parts (a)--(c) of the proof of Lemma~\ref{lem:W} apply verbatim (they do not involve the
rows). \emph{Containment:} a placed row of ZC$'$ has $a^*\ge0$, $a^*+\lfloor a^+-a^*\rfloor\le a^+\le L$, and
$r(X)\ge r(a^*)\ge j+1$ on it, as $r$ is increasing and $a^*\ge(j+1-\chi)/\zeta$. \emph{Disjointness:} the piece lies in
$[a,a^+]\times[j,j+1]$. Every separation statement of Lemma~\ref{lem:W}(e) about a row piece uses only that (1) its lower
left corner $(a',j)$ satisfies $\lambda_\beta(a',j)\ge\xi_{\beta,p(j)}$ (or $a'\ge0$ for a first stretch), (2) its top
right corner $(a'',j+1)$ satisfies $\lambda_{\beta_2}(a'',j+1)\le s_{\beta_2}$ (or $a''\le L$ for a last stretch), and
(3) the piece lies in $\{j\le Y\le j+1\}$. As $a^*\ge a$, $a^*+\lfloor a^+-a^*\rfloor\le a^+$ and $\lambda_\beta$,
$\lambda_{\beta_2}$ are increasing in $X$, (1)--(3) hold for the piece of ZC$'$ because they hold for the segment
$[a,a^+]\times[j,j+1]$; so (e) holds verbatim. \emph{Accounting:} the partition of (f) into cells and stretch regions
does not depend on the placement rule, and the pieces that ZC$'$ places in a stretch lie in its region. Each bound in
(f1)--(f4) is either of the form ``waste at most area'', which holds for any placement, or of the form ``the piece is
placed by ZC and the waste is at most \dots''; in the second case ZC$'$ places the same piece by (O).
\end{proof}

\begin{lemma}[transition rows]\label{lem:P2}
Assume (W1)--(W5), and let $\beta$ be a complete block whose successor $\beta+1$ exists ($s_{\beta+1}+\sec\alpha_{\beta+1}
\le L$). For ZC and for ZC$'$ the total waste of the transition rows of $\beta$ (the stretches with left end on the line
$\lambda_\beta=\hat\xi_\beta$ and right end on the line $\lambda_{\beta+1}=s_{\beta+1}$) is at most
\[
A_T':=(h_T+1)(1+\taz)+\bigl(2+\zeta(w_T+\taz+\zeta\taz)\bigr)(w_T+\taz),
\]
with $h_T$, $w_T$ as in Lemma~\ref{lem:W}. Hence Lemmas~\ref{lem:W} and~\ref{lem:P3a} hold with $A_T$ replaced by $A_T'$.
\end{lemma}

\begin{proof}
Write $\ta$, $\se$, $q$, $M$, $s$, $\hat\xi:=\hat\xi_\beta$, $J_{\rm L}:=J_{\beta,\rm L}$ and $n_{\rm L}$ for block $\beta$, and
$\ta':=\tan\alpha_{\beta+1}\le\ta$, $s':=s_{\beta+1}=\hat\xi+\omega$ with $\omega\in[0,1]$, and
$\hat\kappa:=\ta-\ta'=\kappa_\beta+\eta_\beta\le q_0^2/\zeta+\delta_Q$. By Lemma~\ref{lem:W}(c) the transition rows are
the rows $j\ge J_{\rm L}$. The stretch of row $j$ has left end $a_j:=\hat\xi-j\ta$, right end $a_j^+:=s'-(j+1)\ta'$,
length $\ell_j:=a_j^+-a_j=\omega-\ta'+j\hat\kappa$ and region
$\mathrm{Str}_j:=\{z\in Z:\ j\le Y<j+1,\ \lambda_\beta\ge\hat\xi,\ \lambda_{\beta+1}<s'\}$. The rows $j\ge\lfloor r(L)\rfloor$
are not treated by (Z4); their regions are counted in the same way below, and they are bad rows in the sense of (3)
(as $j+1>r(L)\ge r(a_j)$), so no part of $\{\lambda_\beta\ge\hat\xi,\ \lambda_{\beta+1}<s',\ Y\ge J_{\rm L}\}\cap Z$ is missed.

(1) \emph{Number of rows.} By Lemma~\ref{lem:W}(f3) every point of
$\{\lambda_\beta\ge\hat\xi,\ \lambda_{\beta+1}<s',\ Y\ge J_{\rm L}\}\cap Z$ has $Y<J_{\rm L}+h_T$. So $\mathrm{Str}_j$ is
empty unless $j$ is an integer in $[J_{\rm L},J_{\rm L}+h_T)$: at most $h_T+1$ rows.

(2) \emph{The left ends are positive.} $h_T-1$ equals the right-hand side of (W3) minus $\zeta+(1+\zeta\taz)/q_{\min}$, so
(W3) gives $h_T-1<N_0\caz\le N_0$. For a row of (1), $j<J_{\rm L}+h_T<(M-1)+1+N_0=M+N_0$. By (I2), \eqref{eq:L1} and (W1),
$\hat\xi=s+\Lambda_\beta>M\ta+\Lambda\ge M\ta+N_0\taz\ge(M+N_0)\ta$. Hence $a_j>0$.

(3) \emph{Waste of one row.} The horizontal sections of $\{j\le Y<j+1,\ \lambda_\beta\ge\hat\xi,\ \lambda_{\beta+1}<s'\}$
have length $(s'-Y\ta')-(\hat\xi-Y\ta)=\omega+Y\hat\kappa$, affine in $Y$, equal to $\ell_j+\ta'$ at $Y=j$ and to
$\ell_j+\ta$ at $Y=j+1$; so $\operatorname{area}(\mathrm{Str}_j)\le\max(0,\ell_j+\taz)$. Call row $j$ \emph{good} if
$j+1\le r(a_j)$ and \emph{bad} otherwise. In a good row, by (2) and (O), ZC and ZC$'$ both place
$[a_j,a_j+\lfloor\ell_j\rfloor]\times[j,j+1]$ if $\ell_j\ge1$, and the waste is at most
$\ell_j+\taz-\lfloor\ell_j\rfloor<1+\taz$; if $\ell_j<1$ it is at most $\max(0,\ell_j+\taz)<1+\taz$. In a bad row the waste
is at most $\operatorname{area}(\mathrm{Str}_j)\le\max(0,\ell_j+\taz)$; a point of $\mathrm{Str}_j$ has
$j\le Y\le r(X)\le r(L)$, so if $\mathrm{Str}_j\ne\emptyset$ then $j\le r(L)$ and
$\ell_j\le1+j\hat\kappa\le1+r(L)(q_0^2/\zeta+\delta_Q)\le w_T$; the waste is at most $w_T+\taz$.

(4) \emph{Number of bad rows.} Let $j$ be bad with $\mathrm{Str}_j\ne\emptyset$, and $(X,Y)\in\mathrm{Str}_j$. Then
$j\le Y\le r(X)$ and $X<s'-Y\ta'\le s'-j\ta'$, so $j\le r(s'-j\ta')$, that is $j(1+\zeta\ta')\le H+\zeta\omega$ with
$H:=r(\hat\xi)$. Bad means $j+1>r(\hat\xi-j\ta)$, that is $j(1+\zeta\ta)>H-1$. So $j$ lies in the half-open interval
$\bigl((H-1)/(1+\zeta\ta),\ (H+\zeta\omega)/(1+\zeta\ta')\bigr]$, whose length is
\begin{align*}
\frac{H\zeta(\ta-\ta')+\zeta\omega(1+\zeta\ta)+1+\zeta\ta'}{(1+\zeta\ta)(1+\zeta\ta')}
&\le1+\zeta\bigl(H\hat\kappa+1+\taz+\zeta\taz\bigr)\\
&\le1+\zeta\bigl(w_T+\taz+\zeta\taz\bigr),
\end{align*}
using $\omega\le1$, $\ta,\ta'\le\taz$, a denominator at least $1$, and $H\le r(L)$ (as $\hat\xi\le s'\le L$, since the
successor exists), so that $H\hat\kappa+1\le w_T$. A half-open interval of length $\ell$ contains at most $\ell+1$
integers: there are at most $2+\zeta(w_T+\taz+\zeta\taz)$ bad rows.

Counting every row with the bound $1+\taz$ and the bad rows once more with $w_T+\taz$ gives at most $A_T'$. The other
parts of the proof of Lemma~\ref{lem:W} are unchanged.
\end{proof}

\begin{lemma}[end region of ZC$'$]\label{lem:P3}
Assume (W1)--(W5). For ZC$'$ the last term $w_e(4+2\zeta w_e)$ of $B_W$ may be replaced by $w_e$.
\end{lemma}

\begin{proof}
Let $P_z:=\{(X,Y):\ 0\le X\le L,\ \lfloor r(X)\rfloor\le Y\le r(X)\}$ (the \emph{partial layer}). Its vertical sections
have length $\operatorname{frac}(r(X))<1$, so $\operatorname{area}(P_z\cap\{X\in I\})\le|I|$ for every interval $I$.

\emph{No block.} Every row $j<\lfloor r(L)\rfloor$ is one stretch from $0$ to $L$. The part of its region with
$X\ge a_j^*$ is $[a_j^*,L]\times[j,j+1)$ (inside $Z$, as $r(X)\ge r(a_j^*)\ge j+1$ there), of which ZC$'$ leaves
uncovered less than $1$; the part with $X<a_j^*$ lies in $P_z$ (as in (iii) below). The points with
$Y\ge\lfloor r(L)\rfloor$ lie in $P_z$. So $U_Z\le\lfloor r(L)\rfloor+\operatorname{area}(P_z)\le r(L)+L$, and
$L<s_0+\sez\le w_e$ as in Lemma~\ref{lem:W}(f4); this is at most $r(L)(1+\frac12\taz)+w_e$.

\emph{Blocks exist.} Let $\beta$ be the last block, $\ta:=\tan\alpha_\beta$, $\hat\xi:=\hat\xi_\beta$ and
$E_{\rm end}:=\{z\in Z:\ \lambda_\beta(z)\ge\hat\xi\}$. As in Lemma~\ref{lem:W}(f3), $L-\hat\xi<1+\sez$, so every point
of $E_{\rm end}$ has $X\ge\hat\xi-Y\ta\ge\hat\xi-r(L)\taz>L-w_e$: $E_{\rm end}\subset\{L-w_e<X\le L\}$. No cell meets
$E_{\rm end}$: the cells of $\beta$ lie in $\{\lambda_\beta<\hat\xi\}$, and the cells of an earlier block $\beta''$ lie
in $\{\lambda_{\beta''}<\hat\xi_{\beta''}\}$, while $\lambda_\beta\le\lambda_{\beta''}$ for $Y\ge0$ and
$\hat\xi_{\beta''}\le s_\beta\le\hat\xi$. The \emph{end rows} are the stretches with left end $a_j:=\hat\xi-j\ta$ on the
line $\lambda_\beta=\hat\xi$ and right end $x=L$ (rows $j<\lfloor r(L)\rfloor$); their regions are
$\mathrm{Str}_j=\{z\in Z:\ j\le Y<j+1,\ \lambda_\beta\ge\hat\xi\}$. Split $\mathrm{Str}_j$ by $X$:
\begin{enumerate}[label=(\roman*)]
\item $X\ge a_j^*$: if $a_j^*\le L$ this is the rectangle $[a_j^*,L]\times[j,j+1)$ (its points satisfy $X\ge0$,
$Y<j+1\le r(a_j^*)\le r(X)$ and $\lambda_\beta\ge a_j+j\ta=\hat\xi$); ZC$'$ places $\lfloor L-a_j^*\rfloor$ squares in it
if $L-a_j^*\ge1$, so the waste of this part is less than $1$;
\item $X<a_j$: this part lies in the triangle $\{\hat\xi-Y\ta\le X<\hat\xi-j\ta,\ j\le Y<j+1\}$ of area at most
$\frac12\ta\le\frac12\taz$;
\item $a_j\le X<a_j^*$ (only if $a_j^*>a_j$): points of $Z$ have $X\ge0$, so this part is empty if $a_j^*=0$; otherwise
$a_j^*=(j+1-\chi)/\zeta$, its points have $r(X)<j+1$ and $j\le Y\le r(X)$, so $\lfloor r(X)\rfloor=j$ and they lie in
$P_z\cap E_{\rm end}$.
\end{enumerate}
The points of $E_{\rm end}$ with $Y\ge\lfloor r(L)\rfloor$ (rows not treated by (Z4)) lie in $P_z$, as
$\lfloor r(X)\rfloor\le\lfloor r(L)\rfloor\le Y$. The parts (iii) of different rows and this part are disjoint subsets of
$P_z\cap\{L-w_e<X\le L\}$, of total area at most $w_e$. There is at most one final stretch per row
$j<\lfloor r(L)\rfloor$. A final stretch is either an end row, whose waste outside $P_z$ is less than $1+\frac12\taz$
by (i) and (ii), or, by Lemma~\ref{lem:W}(c), a stretch whose left end is $x=0$ or a slab line with $p\le P_{\beta'}-1$
of some block $\beta'$; such a stretch is placed as by ZC (by (O)) and has waste at most $1+\frac12\taz$ as in
Lemma~\ref{lem:W}(f3). This covers the rows $j<J_{\beta,\rm L}$ of $E_{\rm end}$: if $J_{\beta,0}\le j<J_{\beta,\rm L}$,
then $p(j)\le P_\beta-1$ and $E_{\rm end}\cap\{j\le Y<j+1\}$ lies in the region of the final stretch of row $j$, which
starts on the slab line $\xi_{\beta,p(j)}$ of $\beta$; if $j<J_{\beta,0}$, the last block does not block row $j$, and
$E_{\rm end}\cap\{j\le Y<j+1\}$ lies in the region of the final stretch of row $j$, which starts at $x=0$ or after an
earlier block (with $p\le P_{\beta'}-1$ by Lemma~\ref{lem:W}(c)). So the final stretches together with
$E_{\rm end}\cap\{Y\ge\lfloor r(L)\rfloor\}$ waste at most $r(L)(1+\frac12\taz)+w_e$, which replaces
$r(L)(1+\frac12\taz)+w_e(4+2\zeta w_e)$ in (f4).
\end{proof}

\begin{lemma}[wall filler ZC$'$]\label{lem:Wpp}
Under (W1)--(W5) the construction ZC$'$ does not fail, its pieces form a packing in $Z$, and
\begin{multline*}
U_Z\le B_W'':=Lc_{\rm col}+\Kb\Bigl[(N_0-1)(\varepsilon_W+\taz)+\frac{1+\zeta\taz}{q_{\min}}(1+\taz)+A_T'\Bigr]\\
+\lfloor\chi\rfloor\bigl(1+\tfrac12\taz\bigr)+w_0(1+\zeta w_0)+r(L)\bigl(1+\tfrac12\taz\bigr)+w_e+2(1+\taz),
\end{multline*}
with $A_T'$ of Lemma~\ref{lem:P2} and the other symbols as in Lemma~\ref{lem:W}.
\end{lemma}

\begin{proof}
Lemmas~\ref{lem:P3a}, \ref{lem:P2} and~\ref{lem:P3}. Lemma~\ref{lem:P2} changes the transition rows of the complete
blocks that have a successor, and Lemma~\ref{lem:P3} the final stretches and $E_{\rm end}$; these terms of (f4) are
disjoint. The term $2(1+\taz)$ is not needed (Lemma~\ref{lem:P3} gives $w_e$); it is kept because the certification of
Section~\ref{ssec:b35} uses the larger bound.
\end{proof}

\subsection{Sawtooth accounting of the start rows and the final stretches}\label{ssec:sawtooth}
This subsection is new in version~1.1 and is used only by the tiers marked $\ddagger$ (which also use
Section~\ref{ssec:refine}). The construction ZC$'$ is not changed; only the accounting of two groups of rows is
sharpened: the lengths of these rows form arithmetic progressions, and the sum of the fractional parts of such a
progression is about half its length.

\emph{Notation.} $\operatorname{frac}(x):=x-\lfloor x\rfloor\in[0,1)$. For a block $\beta$ of ZC (equally of ZC$'$) we use
the quantities of Lemma~\ref{lem:W}(a) for that block: $A_\beta:=(\Gamma_{\beta,0}-M_\beta\ca_\beta-\sa_\beta)/q_\beta$ and
$B_\beta:=q_\beta/(\zeta\se_\beta)$, so that $J_{\beta,i}=\lfloor A_\beta+i/B_\beta\rfloor$ for $0\le i\le P_\beta-1$ for
every block, complete or not; $J_{\beta,\rm L}:=J_{\beta,P_\beta-1}$; $\kappa_\beta:=q_\beta^2/\zeta$ (for a complete
block this is $\kappa_\beta$ of (Z3)); and $p(j):=\#\{i:\ J_{\beta,i}\le j\}$. The new hypothesis is
\begin{enumerate}[label=(W6)]
\item $\nu_*:=\varpi\taz-q_0^2/\zeta>0$.
\end{enumerate}

\begin{lemma}[sawtooth sums]\label{lem:SL}
Let $0<\nu<1$, $\epsilon\ge0$, $n\ge0$ an integer, $c$ real, $\pm\in\{+1,-1\}$, and $x_j=c\pm j\nu+d_j$ ($0\le j\le n-1$)
with $d_j\in[0,\epsilon]$. Then
\[
\sum_{j=0}^{n-1}\operatorname{frac}(x_j)\le n\,\frac{1+\nu}2+n\epsilon+\frac1{8\nu}.
\]
\end{lemma}

\begin{proof}
(1) For $d\ge0$, $\operatorname{frac}(y+d)\le\operatorname{frac}(y)+d$: if $\operatorname{frac}(y)+d<1$ the two sides are
equal; otherwise $\operatorname{frac}(y+d)=\operatorname{frac}(y)+d-\lfloor\operatorname{frac}(y)+d\rfloor\le
\operatorname{frac}(y)+d-1$. So it suffices to treat $d_j=0$ and to add $n\epsilon$. (2) The sign $-1$ is the sign $+1$
read backwards: $x_{n-1-i}=c'+i\nu$ with $c':=c-(n-1)\nu$. (3) Let $y_j:=\operatorname{frac}(c+j\nu)$. As $0<\nu<1$,
$y_{j+1}=y_j+\nu$ if $y_j+\nu<1$, and $y_{j+1}=y_j+\nu-1$ otherwise (a \emph{wrap}). Cut $0,\dots,n-1$ before each index
that follows a wrap; this gives maximal runs of consecutive indices. In a run of $r\ge1$ indices the values are
$f,f+\nu,\dots,f+(r-1)\nu=:l$ with $l<1$, and their sum is $r(f+l)/2$. A run that starts right after a wrap has
$f=y_{\rm prev}+\nu-1\in[0,\nu)$ (because $y_{\rm prev}<1\le y_{\rm prev}+\nu$), so its sum is less than $r(1+\nu)/2$. Only
the first run may start higher; its sum is $r(2l-(r-1)\nu)/2<r(2-(r-1)\nu)/2=r(1+\nu)/2+(r-r^2\nu)/2$, and
$(r-r^2\nu)/2\le1/(8\nu)$ for every real $r$. The run lengths add up to $n$.
\end{proof}

The constant $1/(8\nu)$ cannot be lowered: it is attained in the limit by one run of length $1/(2\nu)$ that starts just
below $1-(r-1)\nu$.

\emph{The rows treated below} (ZC$'$ on $Z(\chi,\zeta,L)$; the waste of a stretch is the area of its stretch region minus
the number of squares placed in it, as in Lemmas~\ref{lem:W}(f) and~\ref{lem:P3a}). The \emph{start rows}: if block $0$
exists, the rows $J_{0,0}\le j\le\lfloor\chi\rfloor-1$; the first stretch of such a row runs from $x=0$ to the line
$\lambda_0=s_0$. The \emph{final stretches}: if a block exists, the stretch of row $j$ ($0\le j<\lfloor r(L)\rfloor$)
whose right end is $x=L$, together with the set $E_{\rm end}\cap\{Y\ge\lfloor r(L)\rfloor\}$ of Lemma~\ref{lem:P3}. These
are exactly the stretches bounded in Lemma~\ref{lem:W}(f3) by $\lfloor\chi\rfloor(1+\frac12\taz)$ and in
Lemma~\ref{lem:P3} by $r(L)(1+\frac12\taz)+w_e$.

\begin{lemma}[sawtooth accounting]\label{lem:T1}
Assume (W1)--(W6) and consider ZC$'$ on $Z(\chi,\zeta,L)$.
\begin{enumerate}[label=(\alph*)]
\item If block $0$ exists, the start rows have total waste at most $\lfloor\chi\rfloor(\frac12+\taz)+1/(8\taz)$.
\item If at least one block exists, the final stretches together with $E_{\rm end}\cap\{Y\ge\lfloor r(L)\rfloor\}$ have total
waste at most
\[
r(L)\bigl(\tfrac12+\taz+(\sez-1)\bigr)+\frac{\Kb+2}{8\nu_*}+\frac12\Bigl(2+N_0q_0+\frac{1+\zeta\taz}{q_{\min}}\Bigr)
+\frac12(1+\zeta L+\zeta w_e)+w_e .
\]
\item If no block exists, $U_Z<\lfloor\chi\rfloor+\zeta L+1+L$.
\end{enumerate}
\end{lemma}

All three parts use ZC$'$ and the facts of Lemmas~\ref{lem:W}, \ref{lem:P3a} and~\ref{lem:P3} about it, hence (W1)--(W5).
Beyond that, (a) uses (W5) ($0<\taz\le0.21$); (b) uses (W5), (W6) and (I1), (I2) of Lemma~\ref{lem:W}; (c) uses only
the definitions and $\operatorname{area}(P_z)\le L$ from the proof of Lemma~\ref{lem:P3}.

\begin{proof}
(a) Let $j$ be a start row and $\ell_j:=s_0-(j+1)\taz$. Its stretch region is
$\{j\le Y<j+1,\ X\ge0,\ X+Y\taz<s_0\}$; it lies in $Z$, as $X\le s_0\le L$ and $Y<j+1\le\lfloor\chi\rfloor\le r(X)$. In
(Z4$'$) $a=0$ and $j+1\le\chi$, so $a^*=0$, and ZC$'$ places $[0,\lfloor\ell_j\rfloor]\times[j,j+1]$ if $\ell_j\ge1$. If
$\ell_j\ge0$ the sections of the region have length $s_0-Y\taz$, so its area is $\ell_j+\frac12\taz$ and the waste is
$\operatorname{frac}(\ell_j)+\frac12\taz$. If $\ell_j<0$: $s_0=M_0\taz$ and $M_0=F_0+\lfloor\chi\rfloor-1$, so
$\ell_j\ge(F_0-1)\taz\ge-\taz$, the section at height $Y$ has length at most $\taz(1-(Y-j))$, and the waste (the area) is
at most $\frac12\taz$. In all cases the waste is at most $\operatorname{frac}(\ell_j)+\frac12\taz$. The start rows form
an integer interval of $n\le\lfloor\chi\rfloor$ rows, and $\ell_j=(s_0-\taz)-j\taz$ is an arithmetic progression; Lemma~\ref{lem:SL}
with the sign $-1$, $\nu=\taz\in(0,0.21]$ and $\epsilon=0$ gives the total
$n(1+\taz)/2+1/(8\taz)+n\taz/2\le\lfloor\chi\rfloor(\frac12+\taz)+1/(8\taz)$.

(b) Let the blocks be $0,1,\dots,\beta_e$. Every block except the last is complete, and there are at most $\Kb$ complete
blocks, so $\beta_e+1\le\Kb+1$. For $0\le j<\lfloor r(L)\rfloor$ the final stretch of row $j$ starts after the last block
that blocks row $j$, and at $x=0$ if there is none. Let $R_\beta$ be the set of rows whose last blocking block is $\beta$,
and $R_{\rm none}$ the set of rows blocked by no block; these are integer intervals and partition
$\{0,\dots,\lfloor r(L)\rfloor-1\}$. Four classes:

(F0) $j\in R_{\rm none}$. The final stretch is the whole row: $j+1\le J_{0,0}\le H(J_{0,0})\le\Gamma_{0,0}=r(0)=\chi$, so the
region is the rectangle $[0,L]\times[j,j+1)$ inside $Z$, $a^*=0$, and the waste is less than $1$. By
Lemma~\ref{lem:W}(a) for block $0$, $|R_{\rm none}|\le J_{0,0}<F_0+(1+\zeta\taz)/q_{\min}$, and
$F_0\le\chi-N_0\caz-\saz<\chi-N_0+N_0q_0<2+N_0q_0$; so $|R_{\rm none}|<2+N_0q_0+(1+\zeta\taz)/q_{\min}$.

(F1) $j\in R_\beta$ with $j<J_{\beta,\rm L}$ (for $\beta<\beta_e$ every $j\in R_\beta$, by Lemma~\ref{lem:W}(c)). Then
$p:=p(j)\le P_\beta-1$, the left end is $a_j=\xi_{\beta,p}-j\tan\alpha_\beta\ge0$, the length is
$\ell_j=L-a_j\ge\sec\alpha_\beta>0$, the region has area at most $\ell_j+\frac12\tan\alpha_\beta$, and, as in
Lemma~\ref{lem:W}(f3), $j+1\le r(a_j)$; so by (O) the waste is at most
$\operatorname{frac}(\ell_j)+\frac12\taz$. As in the derivation of (R) in Lemma~\ref{lem:W}(b) (which uses only
$J_{\beta,i}=\lfloor A_\beta+i/B_\beta\rfloor$ for $i\le P_\beta-1$ and $j<J_{\beta,\rm L}$, so it holds for the last block as
well), $p(j)=B_\beta(j+1-A_\beta)+f_j$ with $f_j\in[0,1)$, and modulo $1$
\[
\ell_j\equiv c_\beta'+j\nu_\beta+d_j,\qquad \nu_\beta:=\tan\alpha_\beta-\kappa_\beta,\quad d_j:=(1-f_j)(\sec\alpha_\beta-1)\in[0,\sez-1],
\]
with a constant $c_\beta'$. By (I1) and $q_\beta\le q_0$, $\nu_\beta\ge\varpi\taz-q_0^2/\zeta=\nu_*>0$, and
$\nu_\beta\le\taz\le0.21$. Lemma~\ref{lem:SL} (sign $+1$, $\epsilon=\sez-1$) gives, for the $n_\beta$ rows of (F1) with last
blocking block $\beta$, a total waste at most $n_\beta(\frac12+\taz+(\sez-1))+1/(8\nu_*)$.

(F2) $j\in R_{\beta_e}$ with $j\ge J_{\rm L}:=J_{\beta_e,\rm L}$: the end rows of Lemma~\ref{lem:P3}, with
$a_j=\hat\xi-j\ta$ ($\ta:=\tan\alpha_{\beta_e}$, $\hat\xi:=\hat\xi_{\beta_e}$). The conditions $a_j\ge0$ and
$j+1\le r(a_j)$ are both upper bounds on $j$; let $j_1$ be the least $j\ge J_{\rm L}$ where one of them fails (or
$\lfloor r(L)\rfloor$). (F2a) $J_{\rm L}\le j<j_1$ ($n'$ rows): $a_j^*=a_j$, the part (iii) of the proof of
Lemma~\ref{lem:P3} is empty, the part (i) is $[a_j,L]\times[j,j+1)$ with waste $\operatorname{frac}(L-a_j)$, and the part (ii)
has area at most $\frac12\taz$; $L-a_j=(L-\hat\xi)+j\ta$ is an arithmetic progression with step $\ta\in[\varpi\taz,\taz]$
and $\varpi\taz\ge\nu_*$, so Lemma~\ref{lem:SL} gives at most $n'(\frac12+\taz)+1/(8\nu_*)$. (F2b) $j_1\le j<\lfloor r(L)\rfloor$
($n_{\rm tail}$ rows): the parts (i) and (ii) waste less than $1+\frac12\taz$, and the parts (iii) lie in
$P_z\cap E_{\rm end}$. If $a_{j_1}\ge0$ and $j_1+1>r(a_{j_1})$, then, as $L-\hat\xi<1+\sez$ and $j_1<r(L)$,
$a_{j_1}>\hat\xi-r(L)\taz>L-w_e$, so $j_1+1>r(L)-\zeta w_e$ and $n_{\rm tail}<\zeta w_e+1$. If $a_{j_1}<0$, then as in
Lemma~\ref{lem:W}(f3)(ii) $j_1\ge\lfloor\chi\rfloor$ and $n_{\rm tail}<\zeta L+1$. So $n_{\rm tail}<1+\zeta L+\zeta w_e$. The parts
(iii) of all end rows and $E_{\rm end}\cap\{Y\ge\lfloor r(L)\rfloor\}$ are disjoint subsets of $P_z\cap E_{\rm end}$, of area
at most $w_e$ (proof of Lemma~\ref{lem:P3}).

\emph{Summation.} Put $\bar q:=\frac12+\taz+(\sez-1)$. A row of (F0) wastes less than $1\le\bar q+\frac12$, and a row of
(F2b) less than $1+\frac12\taz\le\bar q+\frac12$; the bounds of (F1) and (F2a) are at most (rows)$\,\bar q+1/(8\nu_*)$ per
chunk, and there are at most $\beta_e+2\le\Kb+2$ chunks. As
$|R_{\rm none}|+\sum_\beta n_\beta+n'+n_{\rm tail}=\lfloor r(L)\rfloor\le r(L)$, the total is at most
$r(L)\bar q+(\Kb+2)/(8\nu_*)+\frac12|R_{\rm none}|+\frac12n_{\rm tail}+w_e$, which gives (b).

(c) Every row $j<\lfloor r(L)\rfloor$ is one stretch from $0$ to $L$. Rows $j\le\lfloor\chi\rfloor-1$: the region is
$[0,L]\times[j,j+1)$ and the waste is less than $1$. Rows $\lfloor\chi\rfloor\le j<\lfloor r(L)\rfloor$ (fewer than
$\zeta L+1$): the part $X\ge a_j^*$ is a rectangle with waste less than $1$, and the part $X<a_j^*$ lies in $P_z$. The
points with $Y\ge\lfloor r(L)\rfloor$ lie in $P_z$, and $\operatorname{area}(P_z)\le L$.
\end{proof}

\begin{lemma}[ZC$'$ with sawtooth accounting]\label{lem:W4}
Under (W1)--(W6) the construction ZC$'$ does not fail, its pieces form a packing in $Z$, and
\begin{multline*}
U_Z\le B_{W4}:=Lc_{\rm col}+\Kb\Bigl[(N_0-1)(\varepsilon_W+\taz)+\frac{1+\zeta\taz}{q_{\min}}(1+\taz)+A_T'\Bigr]\\
+\lfloor\chi\rfloor\bigl(\tfrac12+\taz\bigr)+\frac1{8\taz}+w_0(1+\zeta w_0)+r(L)\bigl(\tfrac12+\taz+(\sez-1)\bigr)+\frac{\Kb+2}{8\nu_*}\\
+\frac12\Bigl(2+N_0q_0+\frac{1+\zeta\taz}{q_{\min}}\Bigr)+\frac12(1+\zeta L+\zeta w_e)+w_e+2(1+\taz).
\end{multline*}
\end{lemma}

\begin{proof}
Packing: Lemma~\ref{lem:P3a}. Bound, if a block exists: the partition (f) of Lemma~\ref{lem:W} for ZC$'$; the cells,
aligned and unaligned closings as in (f1), (f3); the transition rows by Lemma~\ref{lem:P2}; the part of $E_0$ with
$Y\ge\lfloor\chi\rfloor$ by $w_0(1+\zeta w_0)$; the start rows by Lemma~\ref{lem:T1}(a) in place of
$\lfloor\chi\rfloor(1+\frac12\taz)$; the final stretches and $E_{\rm end}\cap\{Y\ge\lfloor r(L)\rfloor\}$ by
Lemma~\ref{lem:T1}(b) in place of $r(L)(1+\frac12\taz)+w_e$. Every stretch is counted once, by its right end. The term
$2(1+\taz)$ is not needed and is kept as slack. If no block exists, Lemma~\ref{lem:T1}(c) gives
$U_Z<\lfloor\chi\rfloor+\zeta L+1+L$, while $B_{W4}\ge\lfloor\chi\rfloor+\zeta L+\frac32+w_e$ and $L<s_0+\sez\le w_e$.
\end{proof}

\subsection{The end packing}\label{ssec:E38}
We use the band end $\Reg(y)$ of Section~\ref{sec:ends} with $\tan(\theta/2)=b/(b^2-1)$, and, as in
Corollary~\ref{cor:b23}, $\mu:=\tan\theta$, $g(X):=y+(X-\sigma)\mu$, $g_0:=g(0)$, $\bar g:=g(w)$ and $\Delta:=w\mu$.

\emph{Construction E38$(y)$.} It depends on three positive parameters $c_D$, $c_R$ (the wall widths) and $a_w$ (the wall
tilt factor); version~1.0 used $c_D=c_R=4$ and $a_w=3/2$.
\begin{enumerate}[label=(E\arabic*)]
\item $m:=\lfloor\bar g\rfloor+1$, $D_\ell:=\lfloor c_Dm^{3/4}\rfloor$, $D_r:=\lfloor c_Rm^{3/4}\rfloor$ and $g_D:=g(D_\ell)$.
\item \emph{Main tilt:} $t:=\lceil t^*Q\rceil/Q$ with $t^*=t^*(m,g_D)$ of Lemma~\ref{lem:tiltS} ($g_0$ replaced by
$g_D$).
\item \emph{Main columns:} $\xi_i:=D_\ell+m\ta+i\se$ for $0\le i<N:=\#\{i\ge0:\ \xi_i+\se\le w-D_r\}$,
$\Gamma_i:=g(\xi_i-m\ta)=g(D_\ell+i\se)$, $J_i:=\lfloor(\Gamma_i-m\ca-\sa)/q\rfloor$ and the columns
$K(\xi_i,J_i,m-J_i)$. Put $J_{\rm L}:=J_{N-1}$ and $\xi_N:=\xi_0+N\se$.
\item \emph{Cut:} $\kappa_m:=q^2/\mu$; $\tau:=\tan\alpha_c$, where $t_c$ is the largest point of $(1/Q^*)\Z$ with
$\tan\alpha_c\le\ta-\kappa_m$; $\hat\kappa:=\ta-\tau$. The phase $\omega_m$ is defined as in (Z3), with
$(s_\beta,\hat\xi_\beta,\tan\alpha_{\beta+1},B,A)$ replaced by $(\xi_0,\xi_{N-1},\tau,q/(\mu\se),(g_D-m\ca-\sa)/q)$ and
the rows $J_0\le j<J_{\rm L}$; put $\ell_0:=\xi_{N-1}+\omega_m$.
\item \emph{Main rows:} for $0\le j<J_{\rm L}$ let $p(j):=\#\{i:\ J_i\le j\}$, $a_j:=\xi_{p(j)}-j\ta$ and
$a_j^+:=\ell_0-(j+1)\tau$; if $a_j^+-a_j\ge1$ place $[a_j,\,a_j+\lfloor a_j^+-a_j\rfloor]\times[j,j+1]$.
\item \emph{Walls.} Each of the three regions below is filled by ZC with $t_0:=\lceil Q^*\tilde t\rceil/Q^*$, where
$\tilde t$ is within relative error $10^{-12}$ of the tilt $t_{\rm tgt}$ with $\tan\alpha(t_{\rm tgt})=a_w\sqrt\zeta$, and
mapped by an exact rotation by $\pm90^\circ$ and a translation:
\begin{itemize}
\item $Z_{\rm L}$: $\chi=\xi_0-g_0\ta$, $\zeta=\ta$, $L=g_0$, $(x,y)\mapsto(y,\,g_0-x)$; region
$\{\lambda\le\xi_0,\ Y\le g_0\}$;
\item $Z_{\rm LR}$: $\chi=w-\ell_0$, $\zeta=\tau$, $L=J_{\rm L}$, $(x,y)\mapsto(w-y,\,x)$; region
$\{X+Y\tau\ge\ell_0,\ Y\le J_{\rm L}\}$;
\item $Z_{\rm UR}$: $\chi=w-\xi_N+J_{\rm L}\ta$, $\zeta=\ta$, $L=Y_u-J_{\rm L}$, $(x,y)\mapsto(w-y,\,J_{\rm L}+x)$, with
$Y_u:=(g_0+\mu\xi_N)/(1+\mu\ta)$; region $\{\lambda\ge\xi_N,\ J_{\rm L}\le Y\le Y_u\}$.
\end{itemize}
\end{enumerate}
\emph{The error of the wall tilts.} Put $\epsilon':=1.03\cdot10^{-12}$, $T_w:=a_w\sqrt\zeta\,(1+\epsilon')+\delta_Q$ and let
$t_{\rm up}:=T_w/(1+\sqrt{1+T_w^2})$ be the tilt with $\tan\alpha(t_{\rm up})=T_w$ (the half-angle formula
$t=\ta/(1+\se)$). Assume $t_{\rm up}<1/10$ (Corollary~\ref{cor:b35} checks this on every box, with a certified margin
of at least $0.055$). As $a_w\sqrt\zeta\le T_w$ and $t\mapsto\ta$ is increasing, $t_{\rm tgt}<t_{\rm up}<1/10$. On
$t\in(0,1/10]$ we have $d\log\ta/d\log t=(1+t^2)/(1-t^2)\le1.01/0.99<1.020203$; since
$|\log(\tilde t/t_{\rm tgt})|\le-\log(1-10^{-12})\le1.0000000000006\cdot10^{-12}$ and the mean value theorem applies on
the interval between $t_{\rm tgt}$ and $\tilde t$ (which lies in $(0,1/10]$, as $\tilde t\le t_{\rm tgt}(1+10^{-12})$), the
relative error of $\tan\alpha(\tilde t)$ is at most $\exp(1.020203\cdot1.0000000000006\cdot10^{-12})-1\le1.0203\cdot10^{-12}
\le\epsilon'$. Finally $\tan\alpha(t_0)\le\tan\alpha(\tilde t)+\delta_Q$, so
\[
\taz\in\bigl[a_w\sqrt\zeta\,(1-\epsilon'),\ a_w\sqrt\zeta\,(1+\epsilon')+\delta_Q\bigr],\qquad t_0\le t_{\rm up}<1/10 .
\]
(Version~1.0 used $t\le0.0014$ and $\epsilon'=1.00001\cdot10^{-12}$, which is valid only for very large $b$.)

\emph{Construction E38$'(y)$} is E38$(y)$ with every wall filled by ZC$'$ (Section~\ref{ssec:refine}) instead of ZC.

Here $\lambda=\lambda(X,Y)=X+Y\ta$ is the slab coordinate of the main tilt. Hypotheses:
(E-i) $1<g_D<m$ and $t\le1/10$; (E-ii) $N\ge1$; (E-iii) $\ta-\kappa_m\ge\tan\alpha(1/Q^*)$, so that $\tau>0$;
(E-iv) $\varepsilon_m+\tau<1$, where $\varepsilon_m:=(\se-1)+m\delta_Q+1/Q_f$; (E-v) each of the three walls has
$\chi\ge4$ and $t_{\rm up}<1/10$ and satisfies (W1)--(W5), for $Z_{\rm LR}$ and $Z_{\rm UR}$ with $L$ replaced by the
upper bound $m$ (which only strengthens (W2) and (W3)). If (E-iii) fails, E38$(y)$ is not defined (the wall
$Z_{\rm LR}$ would have $\zeta=\tau\le0$, outside Lemma~\ref{lem:W}); Lemma~\ref{lem:E38} and Corollary~\ref{cor:b35}
assume (E-iii) and never use that case.

\begin{lemma}[end lemma for the exponent $3/8$]\label{lem:E38}
Under (E-i)--(E-v), E38$(y)$ is a packing in $\Reg(y)$, and its uncovered area satisfies
\begin{multline*}
U(y)\le w\sa+wq+\Delta(\se+m\ta)+J_0(1+\ta)+J_{\rm L}(\varepsilon_m+\ta)+\tfrac12\mu h_{\rm L}^2+\mu w_R^2\\
+U_{Z_{\rm L}}+U_{Z_{\rm LR}}+U_{Z_{\rm UR}},
\end{multline*}
where $h_{\rm L}:=\xi_0-g_0\ta$, $w_R:=w-\xi_N+\bar g\ta$, and each $U_Z\le B_W$ by Lemma~\ref{lem:W}.
The same holds for E38$'(y)$, with each $U_Z\le B_W''$ by Lemma~\ref{lem:Wpp}.
\end{lemma}

The proof uses Lemma~\ref{lem:W} only through ``the pieces of a wall form a packing in the wall region and
$U_Z\le B_W$'', and it uses the parameters only through $D_\ell,D_r\ge0$ (integers) and $t_0\in(0,1/10]$.
So it applies verbatim to E38$'(y)$ with Lemma~\ref{lem:Wpp} in place of Lemma~\ref{lem:W}.

\begin{proof}
(a) \emph{Main staircase.} As in Lemma~\ref{lem:stair}(a), with $g_0$ replaced by $g_D$ and the shift $D_\ell$:
$\Gamma_i=g_D+i\mu\se$ is affine in $i$; $J_i\in[0,m-1]$, because $H(0)=m\ca+\sa\le g_D\le\Gamma_i$ by
Lemma~\ref{lem:tiltS} and $H(m)=m+\sa>\bar g\ge\Gamma_i$; the $J_i$ are non-decreasing; the gaps lie in $[0,q)$; and the
columns lie in $\Reg(y)$ ($X\ge\xi_i-m\ta=D_\ell+i\se\ge0$, $X\le\xi_i+\se\le w-D_r$). The derivation of (R) in the proof
of Lemma~\ref{lem:W}(b) applies verbatim to the rows $J_0\le j<J_{\rm L}$ (with $p(j)=\lceil B(j+1-A)\rceil\in[1,N-1]$,
$B(\se-1)=q^2/\mu=\kappa_m$, $\eta:=\hat\kappa-\kappa_m\in[0,\delta_Q)$, $J_{\rm L}-J_0\le m-1$, and $\hat\xi$ replaced
by $\xi_{N-1}$): the length $\ell_j:=a_j^+-a_j$ satisfies $\operatorname{frac}(\ell_j)\le\varepsilon_m$. Moreover
$\ell_j\ge\ell_0-\xi_{p(j)}-\tau\ge-\tau$ (as $p(j)\le N-1$); together with $\operatorname{frac}(\ell_j)\le\varepsilon_m<
1-\tau$ (E-iv) this forces $\ell_j\ge0$. The rows $j<J_0$ ($p(j)=0$) are not aligned and waste at most $1+\ta$ each. Row
ceiling: for $j<J_{\rm L}$, $p(j)\le N-1$ and $J_{p(j)}\ge j+1$, so
$j+1\le H(J_{p(j)})\le\Gamma_{p(j)}=g(\xi_{p(j)}-m\ta)\le g(a_j)$ (as $j<m$). Also $a_j\ge\xi_0-m\ta=D_\ell\ge0$ and
$a_j^+\le\ell_0\le w$.

(b) \emph{The wall regions lie in $\Reg(y)$.} $Z_{\rm L}=\{0\le Y\le g_0,\ 0\le X\le\xi_0-Y\ta\}$ and $Y\le g_0\le g(X)$.
$Z_{\rm LR}=\{0\le Y\le J_{\rm L},\ \ell_0-Y\tau\le X\le w\}$: its points have $X\ge\ell_0-J_{\rm L}\tau\ge\xi_{N-1}-m\ta$,
so $g(X)\ge\Gamma_{N-1}\ge H(J_{\rm L})\ge J_{\rm L}\ge Y$. $Z_{\rm UR}=\{J_{\rm L}\le Y\le Y_u,\ \xi_N-Y\ta\le X\le w\}$:
by the definition of $Y_u$, $Y_u=g(\xi_N-Y_u\ta)$, so $Y\le Y_u\le g(\xi_N-Y\ta)\le g(X)$; and
$J_{\rm L}\le\Gamma_{N-1}\le g(\xi_N-J_{\rm L}\ta)$ gives $J_{\rm L}\le Y_u$ (the map $Y\mapsto g(\xi_N-Y\ta)-Y$ is
decreasing), that is $L\ge0$. The maps are isometries, so by Lemma~\ref{lem:W} every wall piece lies in its region.

(c) \emph{Disjointness.} Main columns: slabs. Main rows against columns: as in Lemma~\ref{lem:stair}(e) on the left (the
line $\lambda=\xi_{p(j)}$, and columns with $J\ge j+1$ above $Y=j+1$); on the right the rows lie in
$\{X+Y\tau\le\ell_0\}$, the columns $i\le N-2$ in $\{\lambda\le\xi_{N-1}\}\subset\{X+Y\tau\le\ell_0\}$ (as $\tau\le\ta$,
$Y\ge0$ and $\ell_0\ge\xi_{N-1}$), and column $N-1$ lies in $\{Y\ge J_{\rm L}\}$ while the rows lie in $\{Y\le J_{\rm L}\}$.
$Z_{\rm L}\subset\{\lambda\le\xi_0\}$, and all main pieces lie in $\{\lambda\ge\xi_0\}$. $Z_{\rm LR}\subset\{X+Y\tau\ge
\ell_0\}$, against the main rows and the columns $i\le N-2$; against column $N-1$, $Z_{\rm UR}$ and the main rows the line
$Y=J_{\rm L}$ is used where needed. $Z_{\rm UR}\subset\{\lambda\ge\xi_N\}$, against the main columns (in
$\{\lambda\le\xi_N\}$), and against the rows by $Y=J_{\rm L}$. $Z_{\rm L}$ against $Z_{\rm LR}$: $Z_{\rm LR}$ lies in
$\{\lambda\ge\xi_{N-1}\ge\xi_0\}$. $Z_{\rm L}$ against $Z_{\rm UR}$: $\lambda=\xi_0\le\xi_N$. $Z_{\rm LR}$ against
$Z_{\rm UR}$: $Y=J_{\rm L}$.

(d) \emph{Accounting.} Partition $\Reg(y)$ into $\{\lambda<\xi_0\}$ ($=Z_{\rm L}$ plus
$S_1:=\{\lambda<\xi_0,\ Y>g_0\}$); $\{\lambda\ge\xi_N,\ Y\ge J_{\rm L}\}$ ($=Z_{\rm UR}$ plus $S_2$, its part with
$Y>Y_u$); $\{X+Y\tau>\ell_0,\ Y<J_{\rm L}\}$ ($=Z_{\rm LR}$ up to the boundary); the column cells
$\{\xi_i\le\lambda<\xi_{i+1},\ Y\ge J_i\}$ minus the previous set (for $i\le N-2$ the cell does not meet it; for $i=N-1$,
$Y\ge J_{\rm L}$); and the row regions $\{j\le Y<j+1,\ \lambda\ge\xi_{p(j)},\ X+Y\tau\le\ell_0\}$, $0\le j<J_{\rm L}$. The
cells cost at most $w\sa+wq+\Delta(\se+m\ta)$, exactly as in Lemma~\ref{lem:stair}(g) ($N\se\le w$, $N\ta\le w\sa$). A row
region has sections of length $\ell_0-\xi_{p(j)}+Y\hat\kappa$, so its area is at most $\ell_j+\frac12(\ta+\tau)$ and its
waste at most $\varepsilon_m+\ta$ for $J_0\le j<J_{\rm L}$, and at most $1+\ta$ for $j<J_0$. $S_1$ consists of points with
$0\le X<\xi_0-g_0\ta=h_{\rm L}$ and $g_0<Y\le g_0+\mu X$: area at most $\frac12\mu h_{\rm L}^2$. $S_2$ consists of points with
$X\in[\xi_N-\bar g\ta,w]$ and $Y_u<Y\le g(X)\le Y_u+\mu w_R$: area at most $\mu w_R^2$. The three wall regions contribute
$U_{Z_{\rm L}}+U_{Z_{\rm LR}}+U_{Z_{\rm UR}}$ (the maps are isometries).
\end{proof}

\subsection{The band ends}\label{ssec:b35}
A \emph{parameter row} consists of a threshold $b_0\ge10^4$, a window coefficient $c_y\in(0,3/4]$, wall widths
$c_D,c_R>0$, a wall tilt factor $a_w>0$ and the proof parameter $\varpi\in(0,1)$ of Lemma~\ref{lem:W}.
In the rows marked $\dagger$ the walls are filled by ZC$'$ (Lemma~\ref{lem:Wpp}) instead of ZC.
The rows used in this paper are those of Table~\ref{tab:certtiers}; the row ``v1.0'' is the choice of version~1.0.

\begin{corollary}\label{cor:b35}
Let one row of Table~\ref{tab:certtiers} be given, with its constant $C_E$. Let $b\ge b_0$ be an integer,
$\tan(\theta/2)=b/(b^2-1)$, and let $y$ be any real number with
\[
c_yb^{4/5}\le y\le c_yb^{4/5}+2 .
\]
Let $Q$ be a power of $2$ with $Q\ge16b$. Then E38$(y)$
(E38$'(y)$ in the rows marked $\dagger$)
with the parameters $c_D$, $c_R$, $a_w$ of the row is a packing in $\Reg(y)$ with
\[
U(y)\le C_Eb^{3/5},
\]
and the end cost of Proposition~\ref{prop:waste} satisfies $E(y)=U(y)+\frac12\tan\theta\le C_Eb^{3/5}+1.0001/b$.
\end{corollary}

\begin{proof}
Put $z:=m^{-1/4}$.

(i) \emph{Ranges.} By the band facts of Corollary~\ref{cor:b23}(i) (they need $b\ge10^4$; every $b_0$ of
Table~\ref{tab:certtiers} is at least $10^{7.7}$), $2/b<\mu\le2.000001/b$, $\Delta\le2.000001$,
$\sigma\mu\le4.1\cdot10^{-8}$ and $b-2.1/b\le w<b$. As $\bar g=y-\sigma\mu+\Delta\in(y,y+2.000001]$ and
$m=\lfloor\bar g\rfloor+1$, we get $c_yb^{4/5}<m\le c_yb^{4/5}+5.000001$, $\delta:=m-g_0\le3.000001$ and
$m-g_D\le\delta$. From $m>c_yb^{4/5}\ge m-5.000001=m(1-5.000001z^4)$,
\[
\hat b:=bz^5\in\bigl[c_y^{-5/4}(1-5.000001z^4)^{5/4},\ c_y^{-5/4}\bigr],\qquad
z<c_y^{-1/4}b^{-1/5}\le z_0:=c_y^{-1/4}b_0^{-1/5},
\]
and $D_\ell z^3\in(c_D-z^3,c_D]$, $D_rz^3\in(c_R-z^3,c_R]$. \emph{Main tilt.} Let $D'=m+g_D$ and $\delta'=m-g_D$, so that
$t^*=(1+\sqrt{1+D'\delta'})/D'$ (Lemma~\ref{lem:tiltS} with $g_D$ in place of $g_0$). From below,
$t^*\ge\sqrt{\delta'/D'}$ with $D'\le2m$ and $\delta'>\bar g-g_D=(w-D_\ell)\mu\ge2(1-D_\ell/b-2.1/b^2)$, where
$D_\ell/b\le c_Dz^2/\hat b$ and $2.1/b^2=2.1z^{10}/\hat b^2$. From above, $t^*$ is increasing in $\delta'$ and decreasing
in $D'$ (proof of Lemma~\ref{lem:tiltS}); we substitute $D'\ge2m-3.000001$ and $\delta'\le3.000001$ directly into the
formula for $t^*$ (not into the second bound of Lemma~\ref{lem:tiltS}) and use $t\le t^*+1/Q$ and $1/Q\le1/(16b)$. This
gives
\[
\frac t{z^2}\in\Bigl[\sqrt{1-c_Dz^2/\hat b-2.1z^{10}/\hat b^2},\ \frac{z^2+\sqrt{z^4+3.000001(2-3.000001z^4)}}{2-3.000001z^4}
+\frac{z^3}{16\hat b}\Bigr].
\]
Further $\kappa_m=q^2/\mu\le q^2b/2$ and $\tau\in[\ta-\kappa_m-\delta_Q,\ \ta-\kappa_m]$. Wall inputs: $Z_{\rm L}$:
$\zeta=\ta$, $\chi\in[D_\ell,D_\ell+3.000001\ta]$, $L\le m$; $Z_{\rm LR}$: $\zeta=\tau$, $\chi\in[D_r+\se-1,D_r+2\se]$,
$L=J_{\rm L}$; $Z_{\rm UR}$: $\zeta=\ta$, $\chi\in[D_r,D_r+\se+m\ta]$, $L=Y_u-J_{\rm L}$; and
$J_{\rm L}+(Y_u-J_{\rm L})=Y_u\le\bar g<m$. Wall tilts: by (E6),
$\taz\in[a_w\sqrt\zeta(1-\epsilon'),\,a_w\sqrt\zeta(1+\epsilon')+\delta_Q]$ with $\epsilon'=1.03\cdot10^{-12}$, provided
$t_{\rm up}<1/10$, which (iii) checks. (E-i): $m-3.000001\le g_0\le g_D\le\bar g<m$, and $g_D>1$ because
$m>c_yb_0^{4/5}>10^5$ in every row.

(ii) \emph{Normalisation.} Every term of Lemma~\ref{lem:E38} and of $B_W$ (with the $\varpi$ of the row)
(of $B_W''$ in the rows marked $\dagger$)
is multiplied by $z^3=m^{-3/4}$ and rewritten in terms of $\hat b$, $t/z^2$, $\ta/z^2$,
$\zeta/z^2$, $\chi z^3$, $Lz^4\in[0,1]$, $\taz/z$, $q_0/z^2$, $q_{\min}/z^2$, $N_0z^3$ and $\Lambda z^3$, so that only
non-negative powers of $z$ occur. The complete list of the normalised inequalities is Appendix~\ref{app:enc}. For example
$w\sa z^3\le\hat b\,\sa/z^2$, $\Kb=(Lz^4)/(\Lambda z^3)\cdot z^{-1}$, and $J_0\le(4mt+2)/(Qq)$, because
$H(0)=m\ca+\sa$ has $|dH/dt|\le4mt+2$ and $t-t^*\le1/Q$. For the two right walls the terms linear in $L$ are bounded
with the larger of the two coefficients (as $L_{\rm LR}+L_{\rm UR}\le m$) and the other terms are evaluated at $L=m$
(the bounds of Lemmas~\ref{lem:W}
and~\ref{lem:Wpp}
and the hypotheses (W2), (W3) are non-decreasing in $L$).

(iii) \emph{Interval evaluation.} $(0,z_0]$ is split into $N_Z$ equal intervals ($z=0$ included) and, on each, the range
of $t/z^2$ from (i) into $N_T$ equal pieces; on each box all enclosures of Appendix~\ref{app:enc} are evaluated in
interval arithmetic. On every box of the row the lower bounds of the following margins are positive: $1/10-t$; (E-ii);
(E-iii); (E-iv) in the form $1-\varepsilon_m-\tau$; and, for each of the three walls, (W1)--(W5), $\chi-4$ and
$1/10-t_{\rm up}$. The smallest margins and $\Psi_{\sup}$, the supremum over the boxes of the upper bound of
$\Psi(z):=z^3\times$(the bound of Lemma~\ref{lem:E38}), are in Table~\ref{tab:certtiers}.

(iv) By Lemma~\ref{lem:E38},
\begin{align*}
U(y)&\le\Psi_{\sup}m^{3/4}\le\Psi_{\sup}\bigl(c_yb^{4/5}+5.000001\bigr)^{3/4}\\
&\le\Psi_{\sup}c_y^{3/4}\Bigl(1+\frac{5.000001}{c_yb_0^{4/5}}\Bigr)^{3/4}b^{3/5}=:C_E^{\rm cert}b^{3/5}\le C_Eb^{3/5}.
\end{align*}
For $c_y=1/12$ this is the factor $12^{-3/4}(1+60.000012\,b_0^{-4/5})^{3/4}$ of version~1.0. The relation between $E$ and
$U$ is as in Corollary~\ref{cor:b23}(vi).
\end{proof}

\begin{table}[ht]
\centering
\footnotesize
\setlength{\tabcolsep}{2.9pt}
\begin{tabular}{lllllllrrrrrr}
\toprule
row & $b_0$ & $c_y$ & $c_D=c_R$ & $a_w$ & $\varpi$ & boxes & $\Psi_{\sup}$ & $C_E^{\rm cert}$ & $C_E$ & (E-iii) & (W2) &
$t_{\rm up}$\\
\midrule
B1$^\dagger$ & $10^{7.7}$ & $0.170157$ & $2.92683$ & $1.25781$ & $0.65$ & $320\times128$ & $65.9031$ & $17.46022$ & $17.461$ & $0.592$ & $0.0173$ & $0.0558$\\
B1s$^\dagger$ & $10^{7.9}$ & $0.16$ & $2.92683$ & $1.26$ & $0.65$ & $320\times128$ & $67.0681$ & $16.96722$ & $16.968$ & $0.593$ & $0.0167$ & $0.0590$\\
B2$^\dagger$ & $10^{8}$ & $0.154681$ & $2.92683$ & $1.26562$ & $0.65$ & $320\times128$ & $67.7915$ & $16.72079$ & $16.721$ & $0.584$ & $0.0141$ & $0.0603$\\
A1 & $10^{8.5}$ & $0.133923$ & $2.90865$ & $1.13636$ & $0.681818$ & $320\times128$ & $85.6481$ & $18.96099$ & $18.961$ & $0.611$ & $0.0144$ & $0.0706$\\
A1s & $10^{8.58}$ & $0.131625$ & $2.90865$ & $1.13636$ & $0.681818$ & $320\times128$ & $85.6691$ & $18.72102$ & $18.722$ & $0.631$ & $0.0146$ & $0.0716$\\
A2 & $10^{9}$ & $0.116939$ & $2.90909$ & $1.13636$ & $0.681818$ & $320\times128$ & $87.6231$ & $17.52222$ & $17.523$ & $0.661$ & $0.0156$ & $0.0759$\\
A3 & $10^{10}$ & $0.0957429$ & $2.61999$ & $1.13$ & $0.65$ & $320\times128$ & $92.1363$ & $15.85845$ & $15.859$ & $0.926$ & $0.0157$ & $0.0841$\\
A4 & $10^{12}$ & $0.0769401$ & $2.60016$ & $1.13$ & $0.65$ & $320\times128$ & $101.1582$ & $14.77799$ & $14.778$ & $1.606$ & $0.0145$ & $0.0933$\\
A5 & $10^{16}$ & $0.0838149$ & $2.86364$ & $1.25$ & $0.65$ & $320\times128$ & $90.8528$ & $14.15237$ & $14.153$ & $1.947$ & $0.0156$ & $0.0988$\\
\midrule
v1.0 & $10^{16}$ & $1/12$ & $4$ & $3/2$ & $0.65$ & $160\times64$ & $96.4769$ & $14.96365$ & $15.54$ & $1.946$ & $0.1314$ & $0.0986$\\
\bottomrule
\end{tabular}
\medskip
\caption{The parameter rows of Corollary~\ref{cor:b35} and the values of the author's covering. $\Psi_{\sup}$ and
$C_E^{\rm cert}$ are upper bounds, rounded up; $C_E$ is the constant used in Corollary~\ref{cor:b35} and
Theorem~\ref{thm:main38}. The last three columns are lower bounds over all boxes, rounded down, of the margins (E-iii)
(normalised by $z^2$), (W2) (normalised by $z$, smallest of the three walls) and $1/10-t_{\rm up}$ (smallest of the three
walls). The margin $1-\varepsilon_m-\tau$ is at least $0.9969$, and all other margins are positive on every box.
Rows marked $\dagger$ use ZC$'$ and Lemma~\ref{lem:Wpp}.
The row v1.0 is the parameter choice of version~1.0; its $C_E=15.54$ is that of [U].}\label{tab:certtiers}
\end{table}

\emph{Certification.} The values of Table~\ref{tab:certtiers} were computed by the author's program
\texttt{cert\_v3.py} (\texttt{mpmath.iv}, $120$ bits; for each row $320\times128$ boxes, run in two slices of $160$
$z$-intervals), and merged by \texttt{finalize\_v3.py}, which checks that the slices cover every box exactly once and
checks the stated constants of Theorem~\ref{thm:main38} in interval arithmetic. With the row v1.0 ($160\times64$ boxes)
the program reproduces $\Psi_{\sup}=96.47689\ldots$ and $C_E^{\rm cert}\le14.96365$ of the program \texttt{cert3e.py} of
version~1.0, which was written independently. The constants are certified by the author's covering and by three
coverings written independently from the text: two by verification agents (\texttt{mpmath.iv} and \texttt{python-flint}
arb) and one by a reviewer (\texttt{mpmath.iv}), and all three implement the same normalised inequalities (the list of Appendix~\ref{app:enc}), so they check the arithmetic and the coverage but are only a weak check of the transcription of these inequalities from the lemmas; the point checks of the un-normalised formulas (at \(90\), \(120\) and \(5850\) points) partly address that. The reviewer's covering uses the same boxes and reproduces $\Psi_{\sup}$,
$C_E^{\rm cert}$ and the margins of every row to the printed digits, and the worst-box split to within one unit in the
last printed digit; the two other coverings use
adaptive subdivisions and give slightly smaller values of $C_E^{\rm cert}$ (see ``Verification of version~1.1'' at the
end of the paper). All programs were written by AI systems of the same family. In every row the worst box is the last
$z$-interval ($b$ near $b_0$) at the top of the $t$-range; for the row A1s, for example, the bound splits there as column
triangles $30.909$, other main terms $0.287$, left wall $23.608$ and right walls
$30.867$ (in units of $m^{3/4}$). Exponent bookkeeping: the triangles cost about
$bt\asymp bm^{-1/2}\asymp m^{3/4}$ (as $m\asymp b^{4/5}$), and each wall about
$L\taz+\Kb N_0\taz+\Kb/q_{\min}+2\chi\asymp m^{3/4}$, with $\taz\asymp m^{-1/4}$, $\chi\asymp m^{3/4}$ and
$\Kb\asymp m^{1/4}$.

\emph{Why these thresholds.} The threshold $k_i$ enters the proof only through $b_*\ge b_0$
(Section~\ref{ssec:proofthm38}); the other size conditions (the band facts, the room for the lifts, $\lambda'<1$,
$Q\ge16b$) hold far below. With the parameters of version~1.0 (window coefficient $1/12$, $c_D=c_R=4$, $a_w=3/2$) the
hypothesis that fails for smaller $b_0$ is (E-iii), $\tau>0$: $\kappa_m/\ta$ is about $1.84\,c_y^{-3/2}b^{-1/5}$ (about
$76\,b^{-1/5}$ for $c_y=1/12$), and the covering fails for $b$ below about $2.7\cdot10^9$. A larger window coefficient
removes this obstruction (the threshold of (E-iii) on the covering scales like $20.9\,c_y^{-15/2}$), at the price of
larger column triangles; this is why the rows of Table~\ref{tab:certtiers} have $c_y$ between $0.0769$ and $0.1702$.
This threshold is only a rough scale for (E-iii) to hold on the whole window, as the covering requires; for a single
lift $y$, whether $\tau>0$ depends on $y$ (through $m-g_D$) and is not monotone in $b$, and aligned cuts inside the
window occur far below this scale (Remark~\ref{rem:38exact}).
For smaller $b_0$ the remaining cost is the growth of $C_E$ with $z_0$: the terms of $B_W$ of order $m^{1/2}$, in units
of $m^{3/4}$, are of order $z$ with large coefficients.
Lemmas~\ref{lem:P2} and~\ref{lem:P3} reduce these coefficients: for the row B2 with $B_W$ in place of $B_W''$ an
independent covering gives $C_E^{\rm cert}\approx22.50$ instead of $16.72$.
Hypothesis (W2) is not the obstruction to small $b_0$, but it is the binding constraint of the parameter choice: the
rows were found by a local search under the side condition that the normalised margin of (W2) be at least $0.012$ on a
coarser cover, which costs about $1.4\%$ in $C_E$ or less compared with the best parameters found without it. The
parameters are not optimised further, and the per-row bounds of Lemma~\ref{lem:W} remain crude. With the parameters of
version~1.0 the certified bound tends to about $14.6$ as $b_0\to\infty$; this reflects the estimates of the proof, not
$\cs(k)$.

\subsection{Lifts with \texorpdfstring{$\operatorname{frac}(\bar g)\ge1/2$}{frac(g)>=1/2} and the tiers marked \texorpdfstring{$\ddagger$}{double dagger}}\label{ssec:liftT2}
This subsection is new in version~1.1 and is used only by the tiers marked $\ddagger$. Choosing the lift of each band so
that $\operatorname{frac}(\bar g)\ge1/2$ at both of its ends lowers $\delta=m-g_0$ from $3.000001$ to $2.500001$ and
the main tilt accordingly.

\emph{Band notation.} We use the strip frame of Section~\ref{sec:ends}: a band of length $L_b$, with $h$ and $d=\sec\theta$,
a lift $y_0\ge0$ with $L_b\ge h+2y_0$, $R:=\lfloor(L_b-h-2y_0)/d\rfloor+1$ rows and $y_1:=L_b-h-(R-1)d-y_0\in[y_0,y_0+d)$;
the bottom end is $\Reg(y_0)$ and the top end is the image of $\Reg(y_1)$. Both ends have the same function
$\bar g(y)=y+c_0$ with $c_0:=(w-\sigma)\mu$, and $1<d<1.0001$ for $b\ge10^4$.

\emph{Real lifts.} Up to here the lift was a quarter-integer. Lemma~\ref{lem:T2} below gives a real $y_0$ (rational for
integer $b$ and $k$, but with no bounded denominator). Nothing in the results used depends on $y_0$ being a multiple of
$1/4$ or lying on a grid: the strip frame of Section~\ref{sec:ends} takes any lift $y_0\ge0$ and any integer $R\ge1$ with
$y_1\ge0$, and $\Reg(y)$ for any real $y\ge0$; Lemma~\ref{lem:ends} uses $L_b=y_0+y_1+(R-1)d+h$ and separating lines that
are affine in $y_0$; Lemmas~\ref{lem:proj} and~\ref{lem:stack} hold for every base point; Corollary~\ref{cor:assembly} asks
only $y_0\ge0$; the area $A_0(y)$ in Proposition~\ref{prop:waste} is affine in $y$, and the waste identity is an area
decomposition, valid for real $y_0$, $y_1$. Lemma~\ref{lem:E38} and Corollaries~\ref{cor:b35}, \ref{cor:b35c} take any
real $y$ in their windows; the grids $Q$, $Q^*$, $Q_f$ are grids of tilts and of phase offsets, not of heights. Two
independent AI readings (one of them made without this text or the earlier reviews) checked the cited results line by
line: they hold for every real lift and use no lattice, rationality or integrality; they need only (1) $y_0,y_1\ge0$,
(2) $L_b\ge h+2y_0$, (3) $y\ge\sigma\mu$ at both ends for the area identity of Proposition~\ref{prop:waste} (below that
only $E(y)\le U(y)+\frac12\tan\theta$ holds, which is the direction used), and (4) $\bar g(y)\le L_b$ at both ends. For the
lift of Lemma~\ref{lem:T2} with $\alpha_0=c_yb^{4/5}$, as used in Section~\ref{ssec:proofthm38}, these hold: (1)
$y_1\ge y_0\ge\alpha_0>0$; (2) is step (1) of the proof of Lemma~\ref{lem:T2}; (3) $y\ge\alpha_0\ge c_yb_0^{4/5}>8\cdot10^4$, while
$\sigma\mu\le4.1\cdot10^{-8}$ by Corollary~\ref{cor:b23}(i); (4) $\bar g(y_0)\le\bar g(y_1)\le y_1+\Delta\le y_0+3.000001
<h+2y_0+3<k-b\le L_b$ by the room argument of Section~\ref{ssec:proofthm38}. Moreover, the lemma is applied to each band
separately: Lemma~\ref{lem:ends} and Proposition~\ref{prop:waste} are statements about one band, and
Corollary~\ref{cor:assembly}, written there with one common $y_0$, uses of each band only that its squares lie in its
region; so the two bands may have different lifts, and $k^2-N=(k^2-S^2)+\sum_{\text{2 bands}}[(R-1)\tan\theta+E(y_0)+E(y_1)]$
holds with each band's own $y_0$, $y_1$, $R$.

\begin{lemma}[lift choice]\label{lem:T2}
Let $b\ge10^4$, a band $(L_b,h,d)$ as above, and $\alpha_0\ge0$ real with $L_b-h-1\ge2\alpha_0$. Then there is a real $y_0$
with $\alpha_0\le y_0<\alpha_0+(d+1)/2$ such that, with $R$ and $y_1$ defined from $y_0$: $R\ge1$, $y_1-y_0\in\{0,1\}$, and
\[
\operatorname{frac}(\bar g(y_0))=\operatorname{frac}(\bar g(y_1))\ge\tfrac12 .
\]
In particular $y_0$ and $y_1$ lie in $[\alpha_0,\alpha_0+2.0001]$.
\end{lemma}

\begin{proof}
For an integer $R\ge1$ put $S_R:=L_b-h-(R-1)d$, $K_R:=S_R+2c_0$, $\epsilon_R:=0$ if $\lfloor K_R\rfloor$ is odd and
$\epsilon_R:=1$ if it is even, and $Y_R:=(S_R-\epsilon_R)/2$. (1) If $y_0:=Y_R$, then $L_b-h-2y_0=(R-1)d+\epsilon_R\in
[(R-1)d,Rd)$ (as $0\le\epsilon_R\le1<d$), so the row count is $R$, $y_1=S_R-y_0=y_0+\epsilon_R$, and $L_b-h-2y_0\ge0$.
(2) $\bar g(Y_R)=(K_R-\epsilon_R)/2$. The integer $n:=\lfloor K_R-\epsilon_R\rfloor$ is odd, and
$\phi_R:=\operatorname{frac}(K_R-\epsilon_R)=\operatorname{frac}(K_R)$; so $\bar g(Y_R)=(n-1)/2+(1+\phi_R)/2$ with
$(n-1)/2$ an integer and $(1+\phi_R)/2\in[1/2,1)$, that is $\operatorname{frac}(\bar g(y_0))=(1+\phi_R)/2\ge1/2$; and
$\bar g(y_1)=\bar g(y_0)+\epsilon_R$. (3) $Y_R-Y_{R+1}=(d-\epsilon_R+\epsilon_{R+1})/2\in[(d-1)/2,(d+1)/2]$ with
$(d-1)/2>0$, so $Y_R$ decreases strictly to $-\infty$ in steps of at most $(d+1)/2$. (4) $Y_1\ge(L_b-h-1)/2\ge\alpha_0$.
Let $R^*:=\max\{R\ge1:\ Y_R\ge\alpha_0\}$ and $y_0:=Y_{R^*}$. Then $\alpha_0\le y_0\le Y_{R^*+1}+(d+1)/2<\alpha_0+(d+1)/2$, and
$y_1\le y_0+1<\alpha_0+(d+3)/2<\alpha_0+2.00005$.
\end{proof}

The case $y_1=y_0+1$ is rare but covered by the proof. The bounds are nearly attained:
$\operatorname{frac}(\bar g)=1/2$ occurs.

\begin{corollary}\label{cor:b35c}
Let one row of Table~\ref{tab:certC} be given: $b_0$, $c_y$, $c_D$, $c_R$, wall tilt factors $a_{w,\rm L}$,
$a_{w,\rm LR}$, $a_{w,\rm UR}$, $\varpi$ and $C_E$. Let $b\ge b_0$ be an integer and $y$ real with
\[
c_yb^{4/5}\le y\le c_yb^{4/5}+2.0001\qquad\text{and}\qquad\operatorname{frac}(\bar g(y))\ge\tfrac12 .
\]
Let $Q$ be a power of $2$ with $Q\ge16b$. Then E38$'(y)$, with $D_\ell=\lfloor c_Dm^{3/4}\rfloor$,
$D_r=\lfloor c_Rm^{3/4}\rfloor$ and in (E6) the factor $a_{w,\rm L}$ for $Z_{\rm L}$, $a_{w,\rm LR}$ for $Z_{\rm LR}$ and
$a_{w,\rm UR}$ for $Z_{\rm UR}$, is a packing in $\Reg(y)$ with $U(y)\le C_Eb^{3/5}$, and
$E(y)\le C_Eb^{3/5}+1.0001/b$.
\end{corollary}

Lemma~\ref{lem:E38} holds verbatim with a different factor $a_w$ for each wall (its proof does not use the value of
$a_w$), and, if each wall also satisfies (W6), with each $U_Z\le B_{W4}$ of Lemma~\ref{lem:W4}, for the same reason as for
$B_W''$.

\begin{proof}
As Corollary~\ref{cor:b35} (rows with $\dagger$), with Lemma~\ref{lem:W4} in place of Lemma~\ref{lem:Wpp} and these
changes. (i) $\bar g=y-\sigma\mu+\Delta\in(y,y+2.000001]$; since $\operatorname{frac}(\bar g)\ge1/2$,
$m=\lfloor\bar g\rfloor+1=\bar g+1-\operatorname{frac}(\bar g)\le\bar g+\frac12\le c_yb^{4/5}+4.500101$, and
$m>\bar g>y\ge c_yb^{4/5}$; the covering uses $m\le c_yb^{4/5}+4.5002$, that is
$\hat b\in[c_y^{-5/4}(1-4.5002z^4)^{5/4},c_y^{-5/4}]$. Further $\delta=m-g_0=1-\operatorname{frac}(\bar g)+\Delta\le2.500001$
and $m-g_D\le\delta$; the formula for $t^*$ in Corollary~\ref{cor:b35}(i) is increasing in $m-g_D$, and with $2.500002$:
\[
\frac t{z^2}\le\frac{z^2+\sqrt{z^4+2.500002(2-2.500002z^4)}}{2-2.500002z^4}+\frac{z^3}{16\hat b}.
\]
(E-i) holds as $m-2.500001\le g_0\le g_D\le\bar g<m$ and $g_D>1$, since $m>c_yb_0^{4/5}>8\cdot10^4$ in every row of
Table~\ref{tab:certC}; the band facts need $b\ge10^4$, and $b_0\ge10^{7.15}$ there. The lower end of the $t$-range and the
wall inputs are as in Corollary~\ref{cor:b35} ($Z_{\rm L}$ keeps the wider enclosure
$\chi\in[D_\ell,D_\ell+3.000001\ta]$); (E6) is applied to each wall with its own $a_w$. (ii) The normalised inequalities
are those of Appendix~\ref{app:enc} with the replacements of Appendix~\ref{app:encC}. (iii) The covering checks, on every
box, all margins of Corollary~\ref{cor:b35}(iii) and (W6) for each wall. (iv)
$U(y)\le\Psi_{\sup}c_y^{3/4}(1+4.5002/(c_yb_0^{4/5}))^{3/4}b^{3/5}=:C_E^{\rm cert}b^{3/5}\le C_Eb^{3/5}$.
\end{proof}

\begin{table}[ht]
\centering
\footnotesize
\begin{tabular}{llrrrrrrrr}
\toprule
row & $b_0$ & $c_y$ & $\Psi_{\sup}$ & $C_E^{\rm cert}$ & $C_E$ & (E-iii) & (W2) & (W6) & $t_{\rm up}$\\
\midrule
C1$^\ddagger$ & $10^{7.3617}$ & $0.153846$ & $58.9138$ & $14.47251$ & $14.473$ & $0.482$ & $0.0162$ & $0.591$ & $0.0456$\\
C2$^\ddagger$ & $10^{7.2}$ & $0.16$ & $58.5182$ & $14.80459$ & $14.805$ & $0.454$ & $0.0157$ & $0.572$ & $0.0420$\\
C3$^\ddagger$ & $10^{7.15}$ & $0.165$ & $57.8716$ & $14.98288$ & $14.983$ & $0.494$ & $0.0156$ & $0.596$ & $0.0411$\\
\bottomrule
\end{tabular}
\medskip
\caption{The rows of Corollary~\ref{cor:b35c} (tiers marked $\ddagger$) and the values of the author's covering
($320\times128$ boxes). All rows have $c_D=2.81426$, $c_R=3.00931$, $a_{w,\rm L}=1.26$, $a_{w,\rm LR}=1.3608$, $a_{w,\rm UR}=1.35403$ and $\varpi=0.65$. $\Psi_{\sup}$ and $C_E^{\rm cert}$ are upper bounds, rounded up; the
last four columns are lower bounds of the margins (normalised as in Table~\ref{tab:certtiers}; (W6) normalised by $z$),
rounded down. The margin $1-\varepsilon_m-\tau$ is at least $0.9961$, and all other margins are positive on every
box.}\label{tab:certC}
\end{table}

\emph{Certification of the rows $\ddagger$.} The values of Table~\ref{tab:certC} were computed by the author's program
\texttt{cert\_v4.py} (\texttt{mpmath.iv}, $120$ bits, $320\times128$ boxes per row in two slices); its merge step checks the
cover of the boxes and, in interval arithmetic, the stated constants of Theorem~\ref{thm:main38}. They are certified by
the author's covering and by three coverings written independently from the text: two by verification agents
(\texttt{python-flint} arb: $C_E^{\rm cert}\le14.46634$, $14.79786$, $14.97606$; \texttt{mpmath.iv}: $14.469520$,
$14.801342$, $14.979582$, for the rows C1, C2, C3) and one by a reviewer (\texttt{mpmath.iv}, which reproduces $14.47251$,
$14.80459$, $14.982873$), all passing; all three implement the same normalised inequalities (Appendix~\ref{app:encC}),
so they are only a weak check of the transcription of these inequalities from the lemmas. For $b_0\le b\le10^{16}$ this
limitation is addressed by a direct covering (\texttt{python-flint} arb) of the un-normalised bound of Lemma~\ref{lem:E38}
with $B_{W4}$, in the variables $(\log_2b,\operatorname{frac}(\bar g))$, with the integer quantities enclosed exactly on each
box and the actual wall lengths $L$, without Appendix~\ref{app:encC} (only $J_0$ uses the inequality obtained from the mean
value theorem in Corollary~\ref{cor:b35}(ii)): the certified suprema of $U/b^{3/5}$ are $13.910$ (C1), $14.235$ (C2) and
$14.319$ (C3), all margins are positive, and at about $7.6\cdot10^4$ points per row the direct value is at most the
normalised bound. For $b>10^{16}$ the normalised coverings apply. All programs were written by AI systems of the same
family.
The window coefficients $c_y$ were chosen by coarse coverings, and the wall parameters by a local search; they are not
optimised further.

\subsection{Proof of Theorem~\ref{thm:main38}}\label{ssec:proofthm38}
Fix a tier $i$ of Table~\ref{tab:tiers}, its parameter row (Table~\ref{tab:certtiers}) with $b_0$, $c_y$ and
$C_E:=C_{E,i}$, and an integer $k\ge k_i$. Put $b_*:=(5k/(3C_E))^{5/8}$ (so $k=\frac35C_Eb_*^{8/5}$) and
$b:=\lceil b_*\rceil$. As $k\ge k_i\ge\frac35C_Eb_0^{8/5}$, $b\ge b_*\ge b_0$. Take $\tan(\theta/2)=b/(b^2-1)$, so $w<b$
and $S:=k-b+w<k$, and use the packing of Corollary~\ref{cor:assembly} with
\[
y_0:=\tfrac14\bigl\lceil4c_yb^{4/5}\bigr\rceil\in\bigl[c_yb^{4/5},\ c_yb^{4/5}+\tfrac14\bigr]
\]
and, in each band, $R:=\lfloor(L_b-h-2y_0)/d\rfloor+1$. \emph{Room for the lifts:} $b\le b_*+1$ gives
$b^{4/5}\le b_*^{4/5}+1=(5k/(3C_E))^{1/2}+1$, so $y_0\le c_y(5/(3C_E))^{1/2}k^{1/2}+c_y+\frac14\le k^{1/2}+1$, because
$c_y(5/(3C_E))^{1/2}\le1$ and $c_y\le3/4$ (both hold in every row). As $h<3$ (Section~\ref{ssec:proofmain}),
$h+2y_0<2k^{1/2}+5<k-b\le L_b$ (as $k\ge10^{13}$ and $b\le b_*+1\le k^{5/8}+1$, since $5/(3C_E)<1$); so $R\ge1$ and
$y_1\in[y_0,y_0+d)$ with $d<1.0001$, and both lifts lie in the window of Corollary~\ref{cor:b35}. Fill the four end
regions by Corollary~\ref{cor:b35} with $Q$ the least power of $2$ that is at least $16b$. As in
Section~\ref{ssec:proofmain}, by Corollary~\ref{cor:assembly} and Proposition~\ref{prop:waste},
\[
k^2-N<\frac{4k}b+4C_Eb^{3/5}+\lambda',\qquad \lambda':=\frac{4k}{b(b^2-1)}+\frac{4kb}{(b^2-1)^2}+2\tan\theta\le
\frac{8.001\,k}{b^3}+\frac{4.1}b .
\]
With $k\le\frac35C_Eb^{8/5}$ and $b\ge b_0$ we get $\lambda'\le8.001\cdot\frac35C_Eb_0^{-7/5}+4.1/b_0$; in interval
arithmetic this bound is less than $1$ in every tier without $\ddagger$ (its largest value is about $8.32\cdot10^{-8}$, for the row with the
smallest $b_0$); so $\lambda'<1$, and Lemma~\ref{lem:reduction} gives $\cs(k)\le k^2-N-1<4k/b+4C_Eb^{3/5}$. Finally
$4k/b\le4k/b_*=\frac{12}5C_Eb_*^{3/5}$ and, by the mean value theorem, $b^{3/5}\le b_*^{3/5}+\frac35b_*^{-2/5}\le
b_*^{3/5}+\frac35b_0^{-2/5}$, so
\[
\cs(k)<\frac{32}5C_Eb_*^{3/5}+\frac{12}5C_Eb_0^{-2/5}
=\frac{32}5\Bigl(\frac53\Bigr)^{3/8}C_E^{5/8}k^{3/8}+\frac{12}5C_Eb_0^{-2/5}\le C_ik^{3/8}+a_i .
\]
Here $\frac{32}5(\frac53)^{3/8}\approx7.751277$. The inequalities $C_i\ge\frac{32}5(\frac53)^{3/8}C_E^{5/8}$,
$k_i\ge\frac35C_Eb_0^{8/5}$, $a_i\ge\frac{12}5C_Eb_0^{-2/5}$, the inequality ``bound on $\lambda'$ less than $1$'' and
the two room conditions were checked in interval arithmetic for every tier without $\ddagger$ (\texttt{finalize\_v3.py}, which also prints
the value of the bound on $\lambda'$). \qed

\emph{Version~1.0.} The row v1.0 with $C_E=15.54$ gives, by the same argument, the theorem of version~1.0: in interval
arithmetic $\frac{32}5(\frac53)^{3/8}15.54^{5/8}\in[43.05564,\,43.05565]$, $\frac35\cdot15.54\cdot10^{128/5}\le
3.712\cdot10^{26}\le4\cdot10^{26}$ and $\frac{12}5\cdot15.54\cdot10^{-32/5}\le1.49\cdot10^{-5}$, so
$\cs(k)<43.06\,k^{3/8}+2\cdot10^{-5}$ for every integer $k\ge4\cdot10^{26}$. It is superseded by the tier A5.

\emph{The tiers marked $\ddagger$.} For these tiers the proof above is changed as follows. In each of the two bands
(lengths $L_b=S$ and $L_b=k-b$) take the lift $y_0$ of Lemma~\ref{lem:T2} with $\alpha_0:=c_yb^{4/5}$, instead of
$\frac14\lceil4c_yb^{4/5}\rceil$; the two bands may get different lifts (Section~\ref{ssec:liftT2}). Then $y_0$ and $y_1$
lie in $[c_yb^{4/5},c_yb^{4/5}+2.0001]$ and $\operatorname{frac}(\bar g(y_0))=\operatorname{frac}(\bar g(y_1))\ge1/2$,
so Corollary~\ref{cor:b35c} applies to all four ends. \emph{Room:} $\alpha_0\le c_y(5/(3C_E))^{1/2}k^{1/2}+c_y\le k^{1/2}+\frac34$,
and Lemma~\ref{lem:T2} needs $L_b-h-1\ge2\alpha_0$, which holds as $L_b\ge k-b$, $h<3$, $b\le k^{5/8}+1$ and $k\ge10^{12}$
(every $k_i$ of a tier $\ddagger$ is at least $10^{12}$, checked). Then $y_0<\alpha_0+(d+1)/2\le k^{1/2}+1.76$,
$h+2y_0+3<2k^{1/2}+9.6<k-b\le L_b$, and each end region lies in its band. The count, the bound $\lambda'<1$ (checked in
interval arithmetic by the merge step of \texttt{cert\_v4.py}, together with the stated constants, the room conditions and
$k_i\ge10^{12}$; its largest value is about $2.97\cdot10^{-7}$) and the final inequality are as above, with the $C_E$
of the tier.

\begin{remark}[what the explicit computations test]\label{rem:38exact}
The range of Corollary~\ref{cor:b35} ($b\ge b_0\ge10^{7.7}$) has not been tested by explicit computation at the
parameters of a tier and $b\ge b_0$; for Corollary~\ref{cor:b35c} only the $15$ end packings described at the end of the
paper (b from $1.41\cdot10^7$ to $2.30\cdot10^7$, near $b_0$) were built. The constants are certified by interval
arithmetic.
The explicit checks test
the validity of the construction and of the lemmas at smaller $b$. The construction E38 is defined only when (E-iii)
holds; in some runs below (marked ``with $\tau:=0$'') we used the \emph{convention} $\tau:=0$, with the lower right wall
filled by rows only, which is not part of the construction. Whether the cut is aligned ($\tau>0$) depends on $y$ and is
not monotone in $b$.

\emph{Version~1.0 checks.} The construction was implemented in exact rational arithmetic and checked by an exact checker
that does not use the lemmas (shapes, containment, separating axis tests for all pairs with overlapping bounding boxes):
for $b=1600,\dots,409600$ and five lifts each in the window of version~1.0 (all $25$ pass, with $\tau:=0$, so these runs
test validity only), at lifts outside the window where the cut is aligned (all pass, with the residuals of (R) attained
to $4$--$5$ digits), and for the wall filler alone on canonical regions for $m$ up to $10^4$ (certified) and
$2.56\cdot10^6$ (exact internal checks of Lemma~\ref{lem:W}), with $U_Z\le B_W$ in every case. Deliberately corrupted
packings were rejected. Two blank-slate AI reviews and an independent reimplementation of version~1.0 found no
counterexample. The program \texttt{run\_z.py} builds the wall filler for $m=10^4$, $\zeta=0.5/\sqrt m$,
$\chi=4m^{3/4}$, $L=m$ and certifies it: $3$ blocks, $25289$ pieces, $28402$ pairs decided exactly, no overlap, (W1)--(W5)
hold, $7611$ aligned closings with largest residual $0.0056087\le\varepsilon_W=0.0056093$, and $U_Z=6712\le B_W\le19615$.
Full exact counts with the construction of version~1.0 give $\cs(10^8)\le44299$, $\cs(10^{10})\le236005$ and
$\cs(10^{11})\le547881$ (ends checked exactly); at $k=10^8$ this is worse than Theorem~\ref{thm:cert} ($24790$), and
below $k_i$ the closed forms of Theorem~\ref{thm:main38} must not be used.

\emph{Version~1.1 checks} (programs in \texttt{code/zc\_tests}). (a) E38 with the wall parameters of the tiers A1, A5 and
A3 in windows $[c_yb^{4/5},c_yb^{4/5}+2]$ with $c_y=1/3$ or $1/2$ (not a tier's $c_y$) and an aligned cut, for
$b=102400$, $204800$ and $409600$: five end packings with up to $4.8\cdot10^5$ pieces, all checked exactly (no overlap, no
containment failure, the checker's $U$ equal to the construction's).
(b) ZC$'$ on canonical walls, $m=3000,\dots,2.56\cdot10^6$: $U_Z\le B_W''$ in all eight cases (and $U_Z\le B_W$ for ZC), of which (W1)--(W5) hold in six ((W5) fails in two);
exact certificates for $m\le10^4$; ZC$'$ places $1$--$25$ rows that ZC leaves empty. (c) E38$'$ with the wall parameters
of the tiers B2 and B1s, $b=25600,\dots,204800$, aligned cut: four end packings with ZC$'$ walls and one with ZC walls
for comparison, all checked exactly; at $b=102400$ the
rule (Z4$'$) lowers $U$ from $16643$ to $12322$. (d) A blank-slate review implemented ZC$'$ from the text in exact
arithmetic ($36$ walls, $m$ up to $1.6\cdot10^5$) and checked, besides validity, the exact partition identity of
Lemma~\ref{lem:W}(f) and the per-row and per-block bounds of Lemmas~\ref{lem:P2} and~\ref{lem:P3}: no violation; three
deliberately wrong placement rules were caught.
An independent reimplementation of ZC, ZC$'$, E38, E38$'$ and the full L-shaped packing, written from the text alone,
is described under ``Verification of version~1.1'' at the end of the paper.
It includes aligned-cut tests inside the windows of the tiers B1, B1s and B2 at $b\approx1.5\cdot10^6$--$4\cdot10^6$,
below $b_0$.
\end{remark}

\section{The certified packings}\label{sec:cert}

\emph{Generation.} The end packings of Theorem~\ref{thm:cert} were produced by a modification of the construction of
Section~\ref{ssec:stairlemma} with tuned parameters: $b$ was chosen near the minimizer of a model of the total cost, and the lift $y_0$ (about
$0.9\,b^{2/3}$) by a floating-point pre-scan; for each end the column length $m$ is $\lfloor\bar g\rfloor$ or $\lfloor\bar g\rfloor+1$,
whichever leaves less uncovered area; the tilt is the least multiple $t$ of $2^{-40}$ with
$m\ca+\sa\le g(\ca)$ (the top vertex of the first column, which starts at the left wall, has abscissa $\ca$); and each $J_i$ is the largest value, at least $J_{i-1}$, for which the column (of $m-J_i$ squares) fits under the ceiling.
These rules play no role in the verification, which uses only the resulting explicit pieces.

\emph{Proof of Theorem~\ref{thm:cert}.} Let $(k,b,y_0)$ be a row of Table~\ref{tab:cert} and $\tan(\theta/2)=b/(b^2-1)$.
Then $\tan(\theta/2)>1/b$, so $\theta\in(\theta_0,\pi/2)$ and $w<b$ (Lemma~\ref{lem:width}), and $\gamma$, $\sigma$
are rational. For both bands let $R:=\lfloor(L_b-h-2y_0)/d\rfloor+1$, so that $y_0\le y_1<y_0+d$. The data file for $k$
lists, for each of the four ends, its lift ($y_0$ for the lower ends, $y_1$ for the upper ends) and its pieces, with
exact rational coordinates in the frame of $\Reg(y)$: \emph{columns} (a tilted $1\times n$ rectangle, given by its
lowest vertex, $n$ and the rational cosine and sine of its tilt), \emph{layers} and \emph{wedge rows} (axis-parallel
$1\times L$ rectangles). By Corollary~\ref{cor:assembly} and Lemma~\ref{lem:reduction} it suffices to check, in exact
rational arithmetic:
\begin{enumerate}[label=(C\arabic*),start=0]
\item the lift given in the file for each end equals $y_0$ (lower ends) or the computed $y_1$ (upper ends), and every
$n$ and $L$ is a positive integer;
\item $w<b$, $R\ge1$, $y_0\ge0$ and $y_1\ge0$ for both bands;
\item every piece is a rectangle with side lengths $1$ and $n$ (hence a union of $n$ unit squares with pairwise
disjoint interiors); for a column this means $\cos^2+\sin^2=1$ exactly;
\item the four vertices of every piece lie in $\Reg(y)$ (which is convex);
\item two pieces of the same end have disjoint interiors.
\end{enumerate}
For (C4) it suffices to test pairs whose bounding boxes overlap, and for such pairs the separating axis test along the
edge normals of the two rectangles is exact (touching is allowed). The count
$N=(k-b)^2+b(R_{\rm right}+R_{\rm top})+(\text{end squares})$ gives the values $c=k^2-N$ of Table~\ref{tab:cert}, and
Table~\ref{tab:certdata} lists the resulting data. All checks pass (see below). \qed

\begin{table}[ht]
\centering
\small
\begin{tabular}{rrrrrr}
\toprule
$k$ & $R_{\rm right}$ & $R_{\rm top}$ & pieces (4 ends) & body cost & end cost\\
\midrule
$10^5$ & $99868$ & $99245$ & $2887$ & $639.20$ & $944.79$\\
$10^6$ & $999671$ & $997190$ & $10111$ & $1609.72$ & $2363.28$\\
$10^7$ & $9999166$ & $9989291$ & $42550$ & $4048.29$ & $5991.71$\\
$3.825\cdot10^7$ & $38248577$ & $38226491$ & $93680$ & $6925.21$ & $10063.79$\\
$10^8$ & $99997915$ & $99958601$ & $165152$ & $10172.28$ & $14618.72$\\
\bottomrule
\end{tabular}
\medskip
\caption{Data of the certificates. Body cost is $\sum(R-1)\tan\theta$ over the two bands and end cost is the sum of the
four $E(y)$ of Proposition~\ref{prop:waste}; their sum is $c$ up to $k^2-S^2<0.002$ (rounded to two decimals). The
average of $E/b^{2/3}$ over the four ends is between $3.1$ and $3.3$.}
\label{tab:certdata}
\end{table}

\emph{Verification.} Three programs check the certificates in exact rational arithmetic.
\begin{enumerate}[label=(\alph*)]
\item The generating program checks, for each end, that all vertices lie in $\Reg(y)$, that each column lies in its
slab $\{\xi_i\le\lambda\le\xi_i+\se\}$ and above height $J_i$, that the slabs do not overlap, that the $J_i$ are
non-decreasing, that each layer starts to the right of the slabs of the columns below it and each wedge row to the left
of the first slab (the separating lines of Lemma~\ref{lem:stair}(e)), and the exact uncovered area; it also checks the
identity of Proposition~\ref{prop:waste} exactly.
\item A checker written from a description of the construction, without reading the generating program, takes only the
parameters of the pieces from it, builds the vertices itself and checks (C0)--(C4) with the separating axis test for
all pairs of pieces with overlapping bounding boxes, the exact uncovered area, and the count. It reproduced $N$ and $c$
for all five rows (for $k\ge10^7$ on a remote machine), and rejected deliberately corrupted end packings for
$b=400$ and $6400$ (a column shifted, lowered or lengthened, a layer raised, an extra wedge square, a lowered
ceiling).
\item The checker \texttt{stair\_check.py}, included with this text, reads only the published data files.
It recomputes the band, the rows $R$ and the upper lifts $y_1$ from $(k,b,y_0)$, checks (C0)--(C4)
with the separating axis test for all pairs of pieces with overlapping bounding boxes, the exact uncovered areas, the
count and the identity of Proposition~\ref{prop:waste}. It was adapted from checker (b) and does not import the
generating program. Its results are in Table~\ref{tab:verif}; seven deliberately corrupted versions of the data (a
column shifted, lowered or given an extra square, a layer raised or lengthened, a wedge row lengthened, the ceiling
lowered by $1/100$) were all rejected at $k=10^5$ and at $k=10^8$.
\end{enumerate}
\begin{table}[ht]
\centering
\small
\begin{tabular}{rrrrl}
\toprule
$k$ & pieces & pairs decided exactly & overlaps & result (time)\\
\midrule
$10^5$ & $2887$ & $22274$ & $0$ & pass ($0.3$\,s)\\
$10^6$ & $10111$ & $119522$ & $0$ & pass ($1.2$\,s)\\
$10^7$ & $42550$ & $969664$ & $0$ & pass ($8.2$\,s)\\
$3.825\cdot10^7$ & $93680$ & $2730501$ & $0$ & pass ($29.6$\,s)\\
$10^8$ & $165152$ & $5748521$ & $0$ & pass ($59.2$\,s)\\
\bottomrule
\end{tabular}
\medskip
\caption{Results of the checker \texttt{stair\_check.py} (c), summed over the four ends, with Python
\texttt{fractions} on a laptop. For $k\ge10^7$ it was also run with \texttt{gmpy2} rationals on a remote Linux
machine, with the same output apart from the arithmetic and timing lines.}\label{tab:verif}
\end{table}

All checkers were written by AI systems; none has been reviewed by a human. The rows of the bands are not listed in
the data; their correctness rests on Lemmas~\ref{lem:stack} and~\ref{lem:ends}.

The tuned packings are $4.9$--$11.8\%$ better than the packings of the proof of Theorem~\ref{thm:main}
(Remark~\ref{rem:stairexact}); for $k=10^5,\dots,10^8$ the ratio $c/k^{2/5}$ is between $15.6$ and $16.0$, and the
local slopes of $\log c$ against $\log k$ between $10^5$, $10^6$, $10^7$ and $10^8$ are $0.399$, $0.403$ and $0.393$.

\section{Comparison with the literature}\label{sec:lit}

\subsection{From the wasted area to the deficiency}
By Lemma~\ref{lem:translation}, an upper bound $W(x)\le F(x)$ that holds for all $x$ in an interval $(k-\delta,k)$ gives
$\cs(k)\le\limsup_{x\to k^-}F(x)-1$, that is $\cs(k)\le F(k)-1$ if $F$ is continuous. A bound that holds only for some
$x$ (for instance for $x$ whose fractional part is bounded away from $0$) gives nothing for $\cs(k)$. Conversely, a
packing that shows $\cs(k)\le c-1$ shows $W(x)\le x^2-k^2+c$ for $x$ slightly below $k$.

\subsection{Printed bounds}
The bounds in the literature that we found are listed below (literature search of 2026-10-08; the last column is what
they give for $\cs(k)$ by the translation above).

\begin{center}
\small
\begin{tabular}{>{\raggedright}p{4.4cm}>{\raggedright}p{4.6cm}>{\raggedright\arraybackslash}p{5.6cm}}
\toprule
source & statement & consequence for $\cs(k)$\\
\midrule
Erd\H{o}s--Graham 1975 [EG] & $W(x)=O(x^{7/11})$, all real $x$ & $O(k^{7/11})$\\
Chung--Graham 2009 [CG] & $W(x)\le Cx^{(3+\sqrt2)/7}\log x$ & $O(k^{(3+\sqrt2)/7}\log k)$, exponent $\approx0.6306$\\
Wang--Dong--Li 2016 [WDL] (arXiv) & $W(x)\le(16\sqrt2+38)x^{5/8}$ & $\le60.63\,k^{5/8}-1$\\
Arslanov--Mustafin--Shangitbayev 2021 [AMS] & $s(n^2-n)<n$ for $n>12$ & $\le k-1$ for $k>12$\\
Bui 2025 [Bui], McClenagan 2026 [McC] (preprints) & $W(x)=O(x^{3/5})$, all real $x$ & $O(k^{3/5})$, implicit constant\\
\bottomrule
\end{tabular}
\end{center}

The best printed bound for $\cs(k)$ is therefore $O(k^{3/5})$. The statement $W(x)=O(x^{3/5})$ was first claimed by
Chung and Graham [CG2]; an error in that argument is pointed out in [Bui] and [McC]. We have checked the statement of
[WDL, Theorem~1] but not whether every step of its proof applies to $x$ just below an integer. The bound
$\cs(k)=O(k^{3/5})$ is stated explicitly in the online problem collections [Ev] and [Le]. [AMS] is the only refereed
source we found that states a bound in the form of $\cs(k)$; it also notes that [EG] gives $s(n^2-O(n^{7/11}))<n$.
G\"obel [G\"o] remarked that an improvement of the trivial bound is possible for large $n$ in each range
$k^2+k\le n<(k+1)^2$, without an order.

For $k=10^5$ the explicit bound of [WDL] gives about $8.1\cdot10^4$, Theorem~\ref{thm:sqrt} gives $2535$ and
Theorem~\ref{thm:cert} gives $1583$.

For small $k$, the Squares-in-Squares page maintained by D.~Ellsworth [El] (web resource, version of 2026-09-24) lists
packings showing $s(n^2-n+1)<n$ for $n=16,17,18$, that is $\cs(k)\le k-2$ for these $k$.

To the best of our knowledge, no upper bound of order $o(k^{3/5})$ for $\cs(k)$, or for $W(x)$ just below an integer,
has been stated in the searched sources; in particular we found neither $\cs(k)=O(\sqrt k)$ nor any statement about
the ends of tilted bands. The search covered [EG], [RV], [CG], [CG2], [WDL], [Bui] (versions 1 and 2), [McC] and the
thesis [McC2], [AMS], Friedman's survey [Fr] and the Squares-in-Squares page [El], the repositories [Le] and [Ev],
MathOverflow and Mathematics Stack Exchange, and arXiv up to 2026-10-08. Neither [Bui] nor [McC] makes a special
statement for sides just below an integer. A search cannot exclude unpublished or oral results. The ingredients of
Theorem~\ref{thm:A} are standard, and Theorem~\ref{thm:sqrt} may well be regarded as an easy consequence of known
techniques; the staircase of Section~\ref{sec:stair} is the primitive construction of [Bui, Section~3]
(Section~\ref{ssec:credit}).

\subsection{Why uniform bounds cannot go below \texorpdfstring{$\sqrt k$}{sqrt k}}\label{ssec:structural}
Roth and Vaughan [RV, Theorem] proved that $W(x)\gg(x\|x\|)^{1/2}$ whenever $x(x-\lfloor x\rfloor)>1/6$, where $\|x\|$ is
the distance from $x$ to the nearest integer.\footnote{The theorem is sometimes quoted with $x-\lfloor x\rfloor$ in
place of $\|x\|$, for example on p.~1 of [McC]. That form would give $W(x)\gg\sqrt x$ for $x$ just below an integer, and
hence $\cs(k)\gg\sqrt k$, contradicting Theorem~\ref{thm:main}. The original statement [RV, p.~170] has $\|x\|$ and
gives nothing for $x$ just below an integer.} In particular $W(n+\frac12)\gg n^{1/2}$. Hence an upper bound
$W(x)\le Cx^{\alpha}$ valid for all large real $x$ must have $\alpha\ge1/2$, and through \eqref{eq:translation} such a
bound can never give $\cs(k)=O(k^{\alpha})$ with $\alpha<1/2$. A bound of order $o(\sqrt k)$ for $\cs(k)$ must use the
fact that the side is just below an integer, where the theorem of Roth and Vaughan gives nothing ($\|x\|\to0$). The
constructions of this paper do so: the width $b-\varepsilon$ of the bands is just below the integer $b$. Erd\H{o}s and
Graham remark that perhaps the correct bound is $W(x)=O(x^{1/2})$ [EG, p.~122]; Theorem~\ref{thm:sqrt} gives the
corresponding order for $\cs(k)$, and Theorems~\ref{thm:main} and~\ref{thm:main38} go below it, just below integers.

\section{Lower bounds and the order of \texorpdfstring{$\cs(k)$}{c*(k)}}\label{sec:lower}

In the notation of Section~\ref{sec:prelim}, $\cs(k)=\Wmin(k)-1$. The following lower bounds are in preprints by the
author that have not been refereed.
\begin{itemize}
\item {[R]}: $\Wmin(k)\ge0.0353\log k$ for every integer $k\ge2$, where $\log$ is the natural logarithm (version~1.2;
the constant was $0.033$ in earlier versions). Thus $\cs(k)\to\infty$. The constant uses a computer-assisted lemma
([R, Lemma~4.10]).
\item {[Q]}: $\Wmin(k)\ge0.1\,k^{1/4}$ for every integer $k\ge4.62\cdot10^{12}$.
\item {[T]}: $\Wmin(k)>0.06\,k^{1/3}$ for every integer $k\ge3.5\cdot10^{17}$, and
$\liminf_{k\to\infty}\Wmin(k)/k^{1/3}\ge0.06285$. Without [R, Lemma~4.10], $\Wmin(k)>0.05\,k^{1/3}$ for
$k\ge3.39\cdot10^{17}$.
\end{itemize}
For small $k$ see [R, Section~1]: $\Wmin(k)\ge3$ for $k\ge2$, that is $\cs(k)\ge2$, and a computer-assisted result of
E.~Daniel that has not been refereed gives $s(k^2-3)=k$ for $k\ge6$, that is $\cs(k)\ge3$ ([R, Corollary~1.3];
[Le, entry T-064]).

Together with Theorems~\ref{thm:main} and~\ref{thm:main38},
\begin{align*}
0.06\,k^{1/3}-1\ <\ \cs(k)\ &<\ 20.668\,k^{2/5}+0.623\qquad(k\ge3.5\cdot10^{17}),\\
\cs(k)\ &<\ 40.62\,k^{3/8}+1.4\cdot10^{-5}\qquad(k\ge3.4\cdot10^{26}),
\end{align*}
and the other tiers of Theorem~\ref{thm:main38}; so $k^{1/3}\ll\cs(k)\ll k^{3/8}$ is what is proved: the upper bound by Theorem~\ref{thm:main38} of this paper and the
lower bound in [T]. Neither exponent is claimed to be optimal; the true exponent is open
(Section~\ref{sec:open}). The gap is large for moderate $k$: for $k=10^8$ the lower bounds above give only
$\cs(k)\ge2$ ($\cs(k)\ge3$ with the result of Daniel), while Theorem~\ref{thm:cert} gives $\cs(k)\le24790$.

\section{Other end constructions and open problems}\label{sec:open}

Nothing in this section is used elsewhere. By Proposition~\ref{prop:waste}, an end packing with uncovered area
$O(b^{\beta})$, uniformly for lifts in an interval of length $\sec\theta$, gives
$\cs(k)\le\min_b\bigl(4k/b+O(b^{\beta})\bigr)=O(k^{\beta/(1+\beta)})$, with $b\asymp k^{1/(1+\beta)}$. The staircase
has $\beta=2/3$, and the chained wall fillers of Section~\ref{sec:38} have $\beta=3/5$.

\begin{remark}[an intermediate construction: tilted columns]\label{rem:37}
Before the staircase, we filled the end regions by columns of $n$ squares whose tilt is chosen separately for each
column from the fractional part of the ceiling height, so that $n\cos\beta+\sin\beta$ matches the ceiling; neighbouring
columns are separated by slabs as in Lemma~\ref{lem:col}(i), but with an extra gap whenever the tilt increases. The
loss under the ceiling is then about $\sqrt{2/M}$ per unit width for columns of height $M$, and each change of tilt
costs a gap of order $M\tan\beta$; with $M\asymp\sqrt b$ this gives uncovered area $O(b^{3/4})$ and $\cs(k)=O(k^{3/7})$.
The staircase avoids the changes of tilt altogether: one tilt serves all columns, and the rising ceiling is followed
by moving squares into the layers, each move raising a column by $1-\cos\alpha$. Since Theorem~\ref{thm:main} has a smaller
exponent, and its proof does not use this construction, we do not include the details or a proof of the $k^{3/7}$
bound.
\end{remark}

\begin{remark}[Type~3 regions of Chung and Graham]\label{rem:cg3}
The end regions $\Reg(y)$ are right trapezoids of width $w\approx b$ whose two parallel sides differ by about $2$.
Cut by two vertical lines into three trapezoids of width $w/3$, and with lifts $y$ of order $b$, each piece appears to
satisfy the hypotheses of the ``Type~3'' trapezoids of [CG, Theorem~1(3)], for which a packing with uncovered area
$O(x^{1/\sqrt2}\log x)$ is asserted there. Combining Proposition~\ref{prop:waste} with [CG, Theorem~1(3)] therefore
appears to give $\cs(k)=O\bigl(k^{\sqrt2-1}(\log k)^{2-\sqrt2}\bigr)$. The exponent $\sqrt2-1\approx0.414$ is larger
than $2/5$ and than $3/8=0.375$, so this route is weaker than Theorems~\ref{thm:main} and~\ref{thm:main38}. We do not
verify the details, and the proof in [CG] is written in outline, with lower-order terms suppressed.
\end{remark}

\begin{remark}[the wall wedges]\label{rem:bui4}
In the bound of Lemma~\ref{lem:stair}, the column triangles cost about $w\sin\alpha\asymp b/\sqrt m$ and the two walls
(the wedge rows and the ends of the layers) cost about $m$ each, with $m\asymp b^{2/3}$. Section~\ref{sec:38} reduces the
cost of a wall of height $m$ from order $m$ to order $m^{3/4}$ (following ideas of [Bui, Sections~3--4]); then the best
column length is $m\asymp b^{4/5}$, the end cost is $O(b^{3/5})$, and $\cs(k)=O(k^{3/8})$. The constants of
Section~\ref{sec:38} are crude: the parameter rows come from a local search, and several per-row bounds of
Lemma~\ref{lem:W} were not optimised.
\end{remark}

\emph{Open problems.} The true exponent of $\cs(k)$ is open: it lies between $1/3$ ([T]) and $3/8$
(Theorem~\ref{thm:main38}). We do not claim the true exponent, and we make no claim that $3/8$ is the best exponent
obtainable from this frame or from known techniques. Better constants for Theorem~\ref{thm:main38}, and smaller
thresholds $k_i$, would also be of interest. It would also be interesting to have a matching structure for the lower bound, or a direct argument that
some end cost of order $b^{\beta}$ with $\beta>0$ is unavoidable in this frame.

\appendix
\section{The normalised inequalities of Corollary~\ref{cor:b35}}\label{app:enc}
This appendix lists all inequalities evaluated on each box in step (iii) of the proof of Corollary~\ref{cor:b35}, for a
parameter row $(b_0,c_y,c_D,c_R,a_w,\varpi)$; they are those of the program \texttt{cert\_v3.py} (functions
\texttt{box}, \texttt{wall} and \texttt{taurange}). For the row v1.0 they reduce to the list of version~1.0, except for
the forms of (E-iv) and of the tilt condition of the walls, which are stated below. All quantities are intervals; an
inequality ``$>0$'' or ``$\ge0$'' means that the lower end of the enclosure is positive or non-negative, and an upper
bound means the upper end of the enclosure.

\emph{Box and inputs.} $z:=m^{-1/4}\in[z_\ell,z_u]\subset[0,z_0]$ with $z_0:=c_y^{-1/4}b_0^{-1/5}$, and
$\tau_n:=t/z^2\in[\tau_\ell,\tau_u]$, a piece of the range
\[
\Bigl[\sqrt{1-c_Dz^2/\hat b-2.1z^{10}/\hat b^2},\ \frac{z^2+\sqrt{z^4+3.000001(2-3.000001z^4)}}{2-3.000001z^4}+\frac{z^3}{16\hat b}\Bigr]
\]
evaluated on $[z_\ell,z_u]$ (lower end of the first, upper end of the second). Enclosures:
\begin{gather*}
\hat b\in\bigl[c_y^{-5/4}(1-5.000001z^4)^{5/4},\ c_y^{-5/4}\bigr],\quad \mu b\in[2,2.000001],\quad \Delta\le2.000001,\\
D_\ell z^3\in[c_D-z^3,c_D],\quad D_rz^3\in[c_R-z^3,c_R],\\
\delta_Q\le g_qz^5,\ g_q\in[0,\,2.1/(16\hat b)],\qquad
1/Q_f\le f_qz^5,\ f_q\in[0,\,1/(16\hat b)].
\end{gather*}
With $t=\tau_nz^2$:
\begin{gather*}
\frac{\ta}{z^2}=\frac{2\tau_n}{1-t^2},\quad \frac{\sa}{z^2}=\frac{2\tau_n}{1+t^2},\quad q_n:=\frac q{z^4}=\frac{2\tau_n^2}{1+t^2},\\
\se=\frac{1+t^2}{1-t^2},\quad \frac{\se-1}{z^4}=\frac{2\tau_n^2}{1-t^2},\quad
\frac{\kappa_m}{z^2}\le\frac{q_n^2\hat bz}{\mu b},
\end{gather*}
and $\tau/z^2\in[\ta/z^2-\kappa_m/z^2-g_qz^3,\ \ta/z^2-\kappa_m/z^2]$.

\emph{Main part.}
\begin{itemize}
\item (E-i) $1/10-t>0$.
\item (E-ii) $\hat b(1-2.1z^{10}/\hat b^2)-(D_\ell z^3+D_rz^3)z^2-(\ta/z^2)z^3-\se z^5>0$ (this is
$z^5(w-D_\ell-D_r-m\ta-\se)$ from below).
\item (E-iii) $\ta/z^2-\kappa_m/z^2-g_qz^3>0$.
\item (E-iv) $1-\varepsilon_m-(\tau/z^2)z^2>0$ with $\varepsilon_m\le((\se-1)/z^4)z^4+g_qz+f_qz^5$.
\item Main terms of Lemma~\ref{lem:E38}, times $z^3$: column triangles $w\sa z^3\le\hat b\,\sa/z^2$; and the other main
terms: $wqz^3\le\hat bq_nz^2$; $\Delta(\se+m\ta)z^3\le\Delta(\se z^3+(\ta/z^2)z)$;
$J_0(1+\ta)z^3\le\frac{4\tau_nz^2+2z^4}{16\hat bq_n}\bigl(1+(\ta/z^2)z^2\bigr)$;
$J_{\rm L}(\varepsilon_m+\ta)z^3\le((\se-1)/z^4)z^3+g_q+f_qz^4+(\ta/z^2)z$;
$\frac12\mu h_{\rm L}^2z^3\le\mu b\,(h_{\rm L}z^3)^2z^2/(2\hat b)$ with
$h_{\rm L}z^3\in[D_\ell z^3,\ D_\ell z^3+3.000001(\ta/z^2)z^5]$; and
$\mu w_R^2z^3\le\mu b\,(D_rz^3+\se z^3+(\ta/z^2)z)^2z^2/\hat b$.
\end{itemize}

\emph{One wall.} A wall is given by $\zeta_n:=\zeta/z^2$, $\chi_n:=\chi z^3$ and $L_n:=Lz^4\le1$ (the value $L_n=1$ is
used). Then
\begin{gather*}
a_n:=\taz/z\in\bigl[a_w\sqrt{\zeta_n}(1-\epsilon'),\ a_w\sqrt{\zeta_n}(1+\epsilon')+g_qz^4\bigr],\quad\epsilon'=1.03\cdot10^{-12},\\
\sez=\sqrt{1+a_n^2z^2},\quad \caz=1/\sez,\quad \frac{\sez-1}{z^2}=\frac{a_n^2}{\sez+1},\\
q_{0,n}:=\frac{q_0}{z^2}=\frac{a_n^2}{\sez(\sez+1)},\\
q_{\min,n}:=\frac{q_{\min}}{z^2}=\frac{(\varpi a_n)^2}{s_m(s_m+1)},\quad s_m=\sqrt{1+\varpi^2a_n^2z^2},\\
N_0z^3\in[\chi_n-2z^3,\chi_n-z^3],\quad \lfloor\chi\rfloor z^3\in[\chi_n-z^3,\chi_n],\quad
r(L)z^3=\chi_n+\zeta_nL_nz,\\
M_{\max}z^3=r(L)z^3+N_0z^3\,q_{0,n}z^2,\\
\Lambda z^3=\frac{q_{\min,n}N_0z^3-z-\zeta_na_nz^4}{\zeta_n},\quad \Kb z=\frac{L_n}{\Lambda z^3},\\
\frac{\varepsilon_W}{z^2}=\frac{\sez-1}{z^2}+N_0z^3\,g_q+f_qz^3 .
\end{gather*}
Margins of the wall:
\begin{itemize}
\item (W1) $\Lambda z^3-N_0z^3\,a_nz\ge0$ (positive on every box);
\item (W2) $(1-\varpi)a_n-\Delta_{\rm dr}/z>0$ with
$\Delta_{\rm dr}/z=\frac{q_{0,n}}{N_0z^3}\bigl(L_n+\Kb z\,(1+\zeta_na_nz^3)z/\zeta_n\bigr)+\Kb z\,g_qz^3$;
\item (W3) $N_0z^3\caz-z^3\bigl(\zeta_n\sez/q_{\min,n}+a_nz+\zeta_nz^2(2+\sez)+\zeta_nM_{\max}z^3a_n\bigr)
-z(1+\zeta_na_nz^3)/q_{\min,n}\ge0$;
\item (W4) $1-(\varepsilon_W/z^2)z^2>0$;
\item (W5) $\varpi a_n-g_qz^4\ge0$;
\item $\chi\ge4$: $\chi_n-4z^3\ge0$;
\item the tilt condition $t_{\rm up}<1/10$ of (E6): $1/10-a_nz/(1+\sez)>0$, where $a_nz/(1+\sez)$ is the half-angle
formula $t=\ta/(1+\se)$ evaluated on the enclosure of $a_nz$; this also gives $t_0\le1/10$ (the second half of (W5)).
(Version~1.0 checked $a_nz\le1/10$ instead.)
\end{itemize}
The bound: $B_Wz^3\le\alpha L_n+R$ with $\alpha:=c_{\rm col}/z+X/(\Lambda z^3)$, where
\begin{gather*}
c_{\rm col}/z=a_n+\sez q_{0,n}z+\sez^2\zeta_nz+\sez\zeta_nN_0z^3\,q_{0,n}a_nz,\qquad X=X_1+X_2+X_3,\\
X_1=(N_0z^3-z^3)\bigl((\varepsilon_W/z^2)z+a_n\bigr),\qquad X_2=(1+\zeta_na_nz^3)(1+a_nz)/q_{\min,n},\\
h_T=1+\zeta_n\sez/q_{\min,n}+a_nz+\zeta_nz^2(1+\sez)+\zeta_nM_{\max}z^3a_n,\\
w_Tz=z+(r(L)z^3+z^3)(q_{0,n}^2/\zeta_n+g_qz^3),
\end{gather*}
$X_3=h_T\,(w_Tz)\,z$ (this is $A_Tz^2$), and
\begin{gather*}
R=\lfloor\chi\rfloor z^3\bigl(1+\tfrac12a_nz\bigr)+w_0z^3(1+\zeta_nz^2w_0)+r(L)z^3\bigl(1+\tfrac12a_nz\bigr)+R_e,\\
R_e=(w_ez^2)z\,(4+2\zeta_nw_ez^2),\\
w_0=a_nz+N_0z^3\,q_{0,n}a_n,\qquad w_ez^2=z^2(1+\sez)+r(L)z^3\,a_n+N_0z^3\,q_{0,n}a_nz^2 .
\end{gather*}
For $B_W''$ (rows marked $\dagger$) two terms change:
$X_3=(h_T+1)(1+a_nz)z^2+\bigl(2+\zeta_n(w_Tz)z+\zeta_na_nz^3+\zeta_n^2a_nz^5\bigr)\bigl((w_Tz)z+a_nz^3\bigr)$ (this is
$A_T'z^2$ of Lemma~\ref{lem:P2}), and $R_e=(w_ez^2)z+2(1+a_nz)z^3$ (this is $(w_e+2(1+\taz))z^3$, Lemma~\ref{lem:Wpp}).

\emph{The three walls.} $Z_{\rm L}$: $\zeta_n=\ta/z^2$, $\chi_n\in[D_\ell z^3,\ D_\ell z^3+3.000001(\ta/z^2)z^5]$;
$Z_{\rm LR}$: $\zeta_n=\tau/z^2$, $\chi_n\in[D_rz^3+((\se-1)/z^4)z^7,\ D_rz^3+2\se z^3]$; $Z_{\rm UR}$: $\zeta_n=\ta/z^2$,
$\chi_n\in[D_rz^3,\ D_rz^3+\se z^3+(\ta/z^2)z]$. Finally
\[
\Psi\le w\sa z^3+(\text{other main terms})+(\alpha+R)_{Z_{\rm L}}+\max\bigl(\alpha_{Z_{\rm LR}},\alpha_{Z_{\rm UR}}\bigr)
+R_{Z_{\rm LR}}+R_{Z_{\rm UR}},
\]
where the maximum is taken of the upper ends. $\Psi_{\sup}$ is the largest upper end over all boxes, and
$C_E^{\rm cert}=\Psi_{\sup}c_y^{3/4}(1+5.000001/(c_yb_0^{4/5}))^{3/4}$ (upper end).

\section{The normalised inequalities of Corollary~\ref{cor:b35c}}\label{app:encC}
These are the inequalities of Appendix~\ref{app:enc} for the rows marked $\dagger$, with exactly the following
replacements; they are those of the program \texttt{cert\_v4.py}.
\begin{itemize}
\item Inputs: $\hat b\in[c_y^{-5/4}(1-4.5002z^4)^{5/4},\ c_y^{-5/4}]$ (instead of $5.000001$); the upper end of the range of
$\tau_n$ is
\[
\frac{z^2+\sqrt{z^4+2.500002(2-2.500002z^4)}}{2-2.500002z^4}+\frac{z^3}{16\hat b}
\]
(instead of $3.000001$); everything
else as in Appendix~\ref{app:enc}, including $\chi_n\in[D_\ell z^3,D_\ell z^3+3.000001(\ta/z^2)z^5]$ for $Z_{\rm L}$.
\item One wall: $a_w$ is $a_{w,\rm L}$, $a_{w,\rm LR}$ or $a_{w,\rm UR}$ for $Z_{\rm L}$, $Z_{\rm LR}$, $Z_{\rm UR}$, so that
\[
a_n\in\bigl[a_w\sqrt{\zeta_n}(1-\epsilon'),\,a_w\sqrt{\zeta_n}(1+\epsilon')+g_qz^4\bigr]
\]
and $t_{\rm up}$ are computed per wall. The new
margin (W6), normalised by $z$: $\nu_*/z=\varpi a_n-q_{0,n}^2z/\zeta_n>0$.
\item The bound: $B_{W4}z^3\le\alpha L_n+R$ with $\alpha$ as for $B_W''$ and $R=R_{\rm start}+R_{\rm end}+R_e$, where
\begin{gather*}
R_{\rm start}=\lfloor\chi\rfloor z^3\bigl(\tfrac12+a_nz\bigr)+\frac{z^2}{8a_n}+w_0z^3(1+\zeta_nz^2w_0),\\
R_{\rm end}=r(L)z^3\Bigl(\tfrac12+a_nz+\frac{\sez-1}{z^2}z^2\Bigr)+\frac{(\Kb z)z+2z^2}{8\nu_*/z}
+\tfrac12\Bigl(N_0z^3\,q_{0,n}z^2+2z^3+\frac{z(1+\zeta_na_nz^3)}{q_{\min,n}}\Bigr)\\
+\tfrac12\bigl(z^3+\zeta_nL_nz+\zeta_n(w_ez^2)z^3\bigr),\qquad
R_e=(w_ez^2)z+2(1+a_nz)z^3 .
\end{gather*}
\item $B_{W4}$, (W2) and (W3) are non-decreasing in $L$, and (W6) does not involve $L$; so $L_n=1$ is used as before, and
$C_E^{\rm cert}=\Psi_{\sup}c_y^{3/4}(1+4.5002/(c_yb_0^{4/5}))^{3/4}$.
\item Margins checked on every box: those of Appendix~\ref{app:enc} and, for each of the three walls, (W1)--(W6),
$\chi-4$ and $1/10-t_{\rm up}$.
\end{itemize}

\section*{Verification status}
\begin{itemize}
\item Theorems~\ref{thm:A} and~\ref{thm:sqrt}, Corollary~\ref{cor:B}, Proposition~\ref{prop:D} and
Lemmas~\ref{lem:translation}--\ref{lem:max}, \ref{lem:ends}, Corollary~\ref{cor:assembly} and
Proposition~\ref{prop:waste} are proved in full in this text. An earlier write-up of Sections~\ref{sec:block}--\ref{sec:count}
received an independent AI review, which found no critical or major issue; its minor remarks are included (the case
$(k,b)=(2,2)$, $L\ge0$ in Lemma~\ref{lem:band}, and the scope of Lemma~\ref{lem:max}, see Remark~\ref{rem:extra}). The
formulas were also tested in exact arithmetic: \eqref{eq:exact} and the bounds of Theorems~\ref{thm:A}, \ref{thm:sqrt}
and Proposition~\ref{prop:D} for all $2\le b\le k\le1500$, the angle of Lemma~\ref{lem:angle} for all $2\le b\le k\le250$,
and actual packings (separating axis tests) for all $b$ and $k\le60$ and some larger cases. Such tests are not proofs.
\item Theorem~\ref{thm:main} (Section~\ref{sec:stair}) is proved in full; it uses Lemma~\ref{lem:reduction},
Lemmas~\ref{lem:proj}--\ref{lem:stack}, Lemma~\ref{lem:ends}, Corollary~\ref{cor:assembly} (with
Lemma~\ref{lem:valid}) and Proposition~\ref{prop:waste}. An earlier write-up of Section~\ref{sec:stair} received two
independent blank-slate AI reviews, both of which judged it correct with no critical or major issue; their minor
remarks are included. The proof relies on one evaluation in interval arithmetic, $\Phi(10^{-4/3})\le5.0163$, done by the
author's program and by three independent computations; the other constants were certified in the same way. The
construction was tested in exact arithmetic, also by an independent reimplementation (Remark~\ref{rem:stairexact}).
To our knowledge the proof has not appeared before; it deserves priority in a human review.
\item Theorem~\ref{thm:main38} (Section~\ref{sec:38}) is proved in full; it uses Lemmas~\ref{lem:col}
and~\ref{lem:tiltS}, Lemma~\ref{lem:reduction}, Lemmas~\ref{lem:proj}--\ref{lem:stack}, Lemma~\ref{lem:ends},
Corollary~\ref{cor:assembly} and Proposition~\ref{prop:waste}.
The tiers marked $\dagger$ also use Lemmas~\ref{lem:P3a}, \ref{lem:P2}, \ref{lem:P3} and~\ref{lem:Wpp}, which are new in
version~1.1.
The tiers marked $\ddagger$ also use Lemmas~\ref{lem:SL}, \ref{lem:T1}, \ref{lem:W4}, \ref{lem:T2} and
Corollary~\ref{cor:b35c}, which are new in version~1.1.
The write-up of Section~\ref{sec:38} of version~1.0 received two independent blank-slate AI reviews, with their own
programs: one judged it correct after fixes, the other correct; all their fixes were included in version~1.0, and an
independent reimplementation found no counterexample. The changes of version~1.1 received two further independent
blank-slate AI reviews.
Both judged the parameter part correct after minor fixes and the new lemmas of Section~\ref{ssec:refine} correct.
All their fixes are included (among them the value of $\frac{32}5(\frac53)^{3/8}$, the rounding of the tables, the
precise statement of two hypotheses, the derivation of the range of $t/z^2$, and the complete list of
Appendix~\ref{app:enc}). The constants are certified by the author's covering and by three coverings
written independently from the text: two by verification agents (\texttt{mpmath.iv} and \texttt{python-flint} arb) and
one by a reviewer (\texttt{mpmath.iv}), and all three implement the same normalised inequalities (the lists of Appendices~\ref{app:enc} and~\ref{app:encC}), so they check the arithmetic and the coverage but are only a weak check of the transcription of these inequalities from the lemmas; the point checks of the un-normalised formulas (at \(90\), \(120\) and \(5850\) points) partly address that. All programs were written by AI systems of the same family. The constant
$C_E=15.54$ of version~1.0 is certified by three independent full-range coverings ($C_E\le14.964$, $14.963$ and
$14.9631$).
The range $b\ge b_0$ has not been tested by explicit computation, except for $15$ end packings of the tiers $\ddagger$
near $b_0$; the explicit checks cover the validity of the construction at smaller $b$ (Remark~\ref{rem:38exact}).
It deserves priority in a human review.
\item Theorem~\ref{thm:cert} is computer-assisted. Its finite part was checked by three programs, two of them written
independently of the generating program (Section~\ref{sec:cert}).
\item Section~\ref{sec:open} contains remarks only; Remarks~\ref{rem:37} and~\ref{rem:cg3} are not proved here.
\item Section~\ref{sec:lit} and the novelty statement rest on a literature search, which cannot prove absence.
\item No human expert has reviewed this text. The paper has not been refereed. All reviewers and checkers so far are
AI systems of the same family, which may share blind spots.
\end{itemize}

\subsection*{Verification of version~1.1}
The following checks of the changes of version~1.1 were made after the two reviews mentioned above. All programs were
written by AI systems of the same family; none has been reviewed by a human.
\begin{itemize}
\item \emph{Independent covering 1} (verification agent; written from the text alone; \texttt{mpmath.iv}, $120$ bits):
$512$ $z$-strips $\times$ $128$ $t$-pieces with adaptive refinement, about $65{,}537$ boxes per tier; gap-free check
passed; $z=0$ included. All nine tiers certified: $C_E^{\rm cert}$ at most the stated $C_E$, every margin positive; for
example $C_E^{\rm cert}\le18.71447$ (A1s), $14.14983$ (A5), $16.96349$ (B1s), $16.71729$ (B2). Pointwise check of the
un-normalised formulas at $90$ points: passed. Negative control: the parameters of version~1.0 with $b_0=10^9$ fail
(E-iii), as expected.
\item \emph{Independent covering 2} (verification agent; written from the text alone; \texttt{python-flint} arb,
$160$ bits, adaptive bisection): $389$--$3013$ leaves per tier; the partition checked by a Kraft sum equal to $1$ and by
the total area. All nine tiers certified; for example $C_E^{\rm cert}\le18.70756$ (A1s), $14.14845$ (A5), $16.95859$
(B1s), $16.71253$ (B2). Direct un-normalised evaluation at $120$ cases: passed.
\item \emph{Reviewer's covering} (blank-slate review; own code from the text; \texttt{mpmath.iv}, $120$ bits,
$320\times128$ boxes): all nine tiers agree with Table~\ref{tab:certtiers} to the printed digits.
For the tiers marked $\dagger$ it uses $A_T'$ of Lemma~\ref{lem:P2}.
\item \emph{Independent reimplementation} of ZC, ZC$'$, E38, E38$'$ and the full L-shaped packing, written from the text
alone, in exact rational arithmetic, with a separate exact checker (separating axes): $133$ instances,
$4.39\cdot10^7$ pieces, all valid; some of these instances (end packings where (E-iii) fails) use the convention
$\tau:=0$ of Remark~\ref{rem:38exact}.
The wall bounds $U_Z\le B_W''$ hold in $49$ of $49$ walls ((W1)--(W5) hold in $36$ of them), and Lemmas~\ref{lem:P2} and~\ref{lem:P3} hold term by term
in $46$ of $46$ accounted walls.
The first in-window tests with an aligned cut ($\tau>0$), eight end packings at $b=1.5\cdot10^6$--$4\cdot10^6$ with the
parameters of the tiers B1, B1s and B2: all valid; seven of the eight satisfy all hypotheses and are within the bound of
Lemma~\ref{lem:E38}; in the eighth (B1, $b=1.5\cdot10^6$, $y=59407/4$) hypothesis (E-v) fails for the lower right wall,
which was filled by rows only (a convention, not part of the construction), so the bound of Lemma~\ref{lem:E38} does not
apply there. Pointwise check of
Corollary~\ref{cor:b35} at $5850$ points: passed. $29$ of $30$ negative controls were rejected; the remaining one is a
valid packing.
\item \emph{Tiers marked $\ddagger$ (Sections~\ref{ssec:sawtooth} and~\ref{ssec:liftT2}).} Two blank-slate AI reviews:
one found Lemmas~\ref{lem:SL}, \ref{lem:T1}, \ref{lem:W4} and~\ref{lem:T2} correct and Corollary~\ref{cor:b35c} and the
tiers correct after fixes (included); the other found all parts correct. Exact tests (not proofs): $53$ walls ZC$'$
($34$ with all hypotheses true), no violation of Lemmas~\ref{lem:T1} and~\ref{lem:W4}, largest ratio $U_Z/B_{W4}=0.737$;
Lemma~\ref{lem:SL} in $21592$ cases; Lemma~\ref{lem:T2} in $376$ cases and in $3000$ cases forced to $y_1=y_0+1$; no
violation. Three independent coverings of the rows of Table~\ref{tab:certC}, written from the text, all passing
(\texttt{python-flint} arb: $C_E^{\rm cert}\le14.46634$, $14.79786$, $14.97606$; \texttt{mpmath.iv}: $14.469520$, $14.801342$,
$14.979582$; the reviewer's \texttt{mpmath.iv} covering reproduces $14.47251$, $14.80459$, $14.982873$); all three implement
the same normalised inequalities (Appendix~\ref{app:encC}), with point checks of the un-normalised formulas at $46$ and
$36$ points. A direct covering of the un-normalised bound for $b_0\le b\le10^{16}$, without Appendix~\ref{app:encC}
(\texttt{python-flint} arb; integer quantities enclosed exactly per box, actual wall lengths; only $J_0$ uses the
mean-value inequality), certifies $U/b^{3/5}\le13.910$, $14.235$, $14.319$ for C1, C2, C3, all margins positive; for
$b>10^{16}$ the normalised coverings apply. That the cited results of Section~\ref{sec:ends} hold for real lifts was
checked by two independent AI readings (Section~\ref{ssec:liftT2}). Exact end packings: $15$ complete end packings E38$'$
for the rows C1--C3 with lifts of Lemma~\ref{lem:T2} in the window, $b$ from $1.41\cdot10^7$ to $2.30\cdot10^7$, each with
$1.47\cdot10^7$--$2.37\cdot10^7$ pieces, checked exactly (separating axes, containment, exact $U$): all hypotheses hold,
no overlap, no failure; $U$ is at most $0.71$ times the bound of Lemma~\ref{lem:E38} and at most $0.69\,C_Eb^{3/5}$, and
$U_Z/B_{W4}$ lies between $0.44$ and $0.75$; $3$ of the $15$ have $y_1=y_0+1$. Of $30$ negative controls $24$ were
rejected; the other $6$ came from a flaw in the design of the controls and were rejected when applied to the whole
packing.
\end{itemize}

\emph{Programs and data.} \url{https://github.com/squarepacker/k2-minus-c-upper} (archived on Zenodo; the DOI is
given in the repository README). The data files \texttt{stair\_k*.json.gz} (the end pieces of the five certificates),
the checker \texttt{stair\_check.py}, the generator \texttt{stair\_dump.py} and the interval computation
\texttt{const\_stair.py} of the constants of Corollary~\ref{cor:b23} and Section~\ref{ssec:proofmain}, the interval
covering \texttt{cert3e.py} of the constants of Corollary~\ref{cor:b35} for the row v1.0, the author's covering
\texttt{cert\_v3.py} of all rows of Table~\ref{tab:certtiers} with the merge and theorem checks \texttt{finalize\_v3.py}
and the recorded slices, the wall filler \texttt{r38.py} with the exact checker \texttt{chk.py} and the driver
\texttt{run\_z.py}, and the exact tests of version~1.1 in \texttt{code/zc\_tests} (Remark~\ref{rem:38exact})
as well as the author's covering \texttt{cert\_v4.py} of the rows of Table~\ref{tab:certC} with its recorded slices and merges, and
the exact tests \texttt{sl\_t2\_test.py} of Lemmas~\ref{lem:SL} and~\ref{lem:T2}
accompany this text; \texttt{README\_code.txt} lists the commands and the expected outputs. The programs of the reviews and of
the independent coverings and reimplementations are not included.

\subsection*{Use of AI}
Developed with the assistance of Claude (Anthropic), including the proofs, the text and the programs; the author takes
full responsibility.

\subsection*{Licence}
This text is distributed under the Creative Commons Attribution 4.0 International licence (CC BY 4.0); the accompanying
programs are distributed under the MIT licence.

\begin{thebibliography}{WDL}
\raggedright
\bibitem[AMS]{AMS} M.~Z.~Arslanov, S.~A.~Mustafin and Z.~K.~Shangitbayev, Improved packings of $n(n-1)$ unit squares in a
square, Electron.\ J.\ Combin.\ \textbf{28}(4) (2021), \#P4.22.
\bibitem[Bui]{Bui} H.~D.~Bui, Square packing with $O(x^{0.6})$ wasted area, arXiv:2508.04603v2 (2025, revised 2026).
\bibitem[CG]{CG} F.~Chung and R.~Graham, Packing equal squares into a large square, J.\ Combin.\ Theory Ser.~A
\textbf{116} (2009), 1167--1175.
\bibitem[CG2]{CG2} F.~Chung and R.~Graham, Efficient packings of unit squares in a large square, Discrete Comput.\ Geom.\
\textbf{64} (2020), 690--699.
\bibitem[EG]{EG} P.~Erd\H{o}s and R.~L.~Graham, On packing squares with equal squares, J.\ Combin.\ Theory Ser.~A
\textbf{19} (1975), 119--123.
\bibitem[El]{El} D.~Ellsworth, Squares in Squares (web page, version of 2026-09-24),
\url{https://kingbird.myphotos.cc/packing/squares_in_squares.html}.
\bibitem[Ev]{Ev} evand/square-packing (online problem collection), \url{https://github.com/evand/square-packing},
consulted 2026-10-08.
\bibitem[Fr]{Fr} E.~Friedman, Packing unit squares in squares: a survey and new results, Electronic J.\ Combin.,
Dynamic Survey DS7, version of 2009 (errata 2023).
\bibitem[G\"o]{Go} F.~G\"obel, Geometrical packing and covering problems, in: Packing and Covering in Combinatorics
(A.~Schrijver, ed.), Math.\ Centre Tracts \textbf{106}, Mathematisch Centrum, Amsterdam, 1979.
\bibitem[Le]{Le} jlevy/squares (online collection), \url{https://github.com/jlevy/squares}, consulted 2026-10-08.
\bibitem[McC]{McC} R.~McClenagan, Optimally packing a large square by unit squares, arXiv:2602.01484 (2026).
\bibitem[McC2]{McC2} R.~McClenagan, Master's thesis, UNBC (University of Northern British Columbia), 2024.
\bibitem[Q]{Q} S.~Ryu, Packing $k^2-c$ unit squares: a lower bound of order $k^{1/4}$ for the deficiency, preprint,
version~1.0, 2026, doi:10.5281/zenodo.23212643 (all versions: doi:10.5281/zenodo.23212642); programs:
doi:10.5281/zenodo.23211207.
\bibitem[R]{R} S.~Ryu, Packing $k^2-c$ unit squares: $s(k^2-c)=k$ for all large $k$, preprint, version~1.2, 2026,
doi:10.5281/zenodo.23194104; programs and data: doi:10.5281/zenodo.23194031 and
\url{https://github.com/squarepacker/k2-minus-c} (release v1.2).
\bibitem[RV]{RV} K.~F.~Roth and R.~C.~Vaughan, Inefficiency in packing squares with unit squares, J.\ Combin.\ Theory
Ser.~A \textbf{24} (1978), 170--186.
\bibitem[T]{T} S.~Ryu, Packing $k^2-c$ unit squares: a lower bound of order $k^{1/3}$ for the deficiency, preprint,
version~1.0, 2026, doi:10.5281/zenodo.23213564 (all versions: doi:10.5281/zenodo.23213563); programs:
doi:10.5281/zenodo.23212059 (all versions: doi:10.5281/zenodo.23212058).
\bibitem[U]{U} S.~Ryu, Packing $k^2-c$ unit squares: an upper bound of order $k^{3/8}$ for the deficiency, preprint,
version~1.0, 2026, doi:10.5281/zenodo.23256655 (all versions: doi:10.5281/zenodo.23256654). Version~1.0 of this paper.
\bibitem[WDL]{WDL} S.~Wang, T.~Dong and J.~Li, A new result on packing unit squares into a large square,
arXiv:1603.02368v2 (2016).
\end{thebibliography}
\end{document}
